% This paper has been transcribed in Plain TeX by
% David R. Wilkins
% School of Mathematics, Trinity College, Dublin 2, Ireland
% (dwilkins@maths.tcd.ie)
%
% Trinity College, 2000.

\magnification=\magstep1
\vsize=227 true mm \hsize=170 true mm
   \voffset=-0.4 true mm \hoffset=-5.4 true mm

\def\folio{\ifnum\pageno>0 \number\pageno \else
   \ifnum\pageno<0 \romannumeral-\pageno \else\fi\fi}

\font\Largebf=cmbx10  scaled \magstep2
\font\largerm=cmr12
\font\largeit=cmti12
\font\largesc=cmcsc10  scaled \magstep1
\font\sc=cmcsc10

\def\multieqalign#1{\null\,\vcenter{\openup1\jot \mathsurround=0pt
   \ialign{\strut\hfil$\displaystyle{##}$&$\displaystyle{{}##}$\hfil
   &&\quad \strut\hfil$\displaystyle{##}$&$\displaystyle{{}##}$\hfil
    \crcr#1\crcr}}\,}

\def\sum{\mathop{\Sigma}\nolimits}

\def\sectiontitle{\it\noindent}

\pageno=0

\null\vskip72pt

\centerline{\Largebf ON A GENERAL METHOD IN DYNAMICS}

\vskip24pt

\centerline{\Largebf By}

\vskip24pt

\centerline{\Largebf William Rowan Hamilton}

\vskip24pt

\centerline{\largerm (Philosophical Transactions of the Royal Society,
   part~II for 1834, pp.\ 247--308.)}

\vskip36pt

\vfill

\centerline{\largerm Edited by David R. Wilkins}

\vskip 12pt

\centerline{\largerm 2000}

\vskip36pt\eject

\pageno=-1

\null\vskip36pt

\centerline{\Largebf NOTE ON THE TEXT}

\bigskip

This edition is based on the original publication in the {\it
Philosophical Transactions of the Royal Society}, part~II for
1834.


The following errors in the original published text have been
corrected:

\smallskip

\item{}
a term $w^{(n)}$ in the last summand on the right hand side of
equation (S${}^5$.) has been corrected to $w^{(n-1)}$;

\smallskip

\item{}
a minus sign ($-$) missing from equation (K${}^6$.) has been
inserted.

\smallskip

The paper {\it On a General Method in Dynamics} has also
been republished in {\it The Mathematical Papers of Sir
William Rowan Hamilton, Volume II: Dynamics}, edited for
the Royal Irish Academy by A.~W. Conway and A.~J. McConnell, and
published by Cambridge University Press in 1940.

\bigbreak\bigskip

\line{\hfil David R. Wilkins}

\vskip3pt

\line{\hfil Dublin, February 2000}

\vfill\eject

\pageno=1

\null\vskip36pt

\noindent
{\largeit
On a General Method in Dynamics; by which the Study of the
Motions of all free Systems of attracting or repelling Points is
reduced to the Search and Differentiation of one central Relation,
or characteristic Function.
By {\largesc William Rowan Hamilton}, Member of several scientific
Societies in the British Dominions, and of the American Academy
of Arts and Sciences, Andrews' Professor of Astronomy in the
University of Dublin, and Royal Astronomer of Ireland.
Communicated by Captain {\largesc Beaufort}, R.N. F.R.S.}

\vskip12pt

\centerline{Received April 1,---Read April 10, 1834.}
\vskip12pt
\centerline{[{\it Philosophical Transactions of the Royal Society},
   part~II for 1834, pp.\ 247--308.]}

\bigbreak

{\sectiontitle
Introductory Remarks.\par}

\nobreak\bigskip

The theoretical development of the laws of motion of bodies is a
problem of such interest and importance, that it has engaged the
attention of all the most eminent mathematicians, since the
invention of dynamics as a mathematical science by {\sc Galileo},
and especially since the wonderful extension which was given to that
science by {\sc Newton}.  Among the successors of those illustrious
men, {\sc Lagrange} has perhaps done more than any other analyst, to
give extent and harmony to such deductive researches, by showing
that the most varied consequences respecting the motions of
systems of bodies may be derived from one radical formula; the
beauty of the method so suiting the dignity of the results, as to
make of his great work a kind of scientific poem.  But the
science of force, or of power acting by law in space and time,
has undergone already another revolution, and has become already
more dynamic, by having almost dismissed the conceptions of solidity
and cohesion, and those other material ties, or geometrically
imaginably conditions, which {\sc Lagrange} so happily reasoned on,
and by tending more and more to resolve all connexions and actions
of bodies into attractions and repulsions of points: and while the
science is advancing thus in one direction by the improvement of
physical views, it may advance in another direction also by the
invention of mathematical methods.  And the method proposed in
the present essay, for the deductive study of the motions of
attracting or repelling systems, will perhaps be received with
indulgence, as an attempt to assist in carrying forward so high
an inquiry.

In the methods commonly employed, the determination of the motion
of a free point in space, under the influence of accelerating
forces, depends on the integration of three equations in ordinary
differentials of the second order; and the determination of the
motions of a system of free points, attracting or repelling one
another, depends on the integration of a system of such
equations, in number threefold the number of the attracting or
repelling points, unless we previously diminish by unity this
latter number, by considering only relative motions.  Thus, in
the solar system, when we consider only the mutual attractions of
the sun and the ten known planets, the determination of the
motions of the latter about the former is reduced, by the usual
methods, to the integration of a system of thirty ordinary
differential equations of the second order, between the
coordinates and the time; or, by a transformation of {\sc Lagrange},
to the integration of a system of sixty ordinary differential
equations of the first order, between the time and the elliptic
elements: by which integrations, the thirty varying coordinates,
or the sixty varying elements, are to be found as functions of
the time.  In the method of the present essay, this problem is
reduced to the search and differentiation of a single function,
which satisfies two partial differential equations of the first
order and of the second degree: and every other dynamical
problem, respecting the motions of any system, however numerous,
of attracting or repelling points, (even if we suppose those
points restricted by any conditions of connexion consistent with
the law of living force,) is reduced, in like manner, to the
study of one central function, of which the form marks out and
characterizes the properties of the moving system, and is to be
determined by a pair of partial differential equations of the
first order, combined with some simple considerations.  The
difficulty is therefore at least transferred from the integration
of many equations of one class to the integration of two of
another: and even if it should be thought that no practical
facility is gained, yet an intellectual pleasure may result from
the reduction of the most complex and, probably, of all
researches respecting the forces and motions of body, to the
study of one characteristic function,\footnote*{{\sc Lagrange}
and, after him, {\sc Laplace} and others, have employed a single
function to express the different forces of a system, and so to
form in an elegant manner the differential equations of its
motion.  By this conception, great simplicity has been given
to the statement of the problem of dynamics; but the solution
of that problem, or the expression of the motions themselves,
and of their integrals, depends on a very different and hitherto
unimagined function, as it is the purpose of this essay to show.}
the unfolding of one central relation.

The present essay does not pretend to treat fully of this
extensive subject,---a task which may require the labours of many
years and many minds; but only to suggest the thought and propose
the path to others.  Although, therefore, the method may be used
in the most varied dynamical researches, it is at present only
applied to the orbits and perturbations of a system with any laws
of attraction or repulsion, and with one predominant mass or
centre of predominant energy; and only so far, even in this one
research, as appears sufficient to make the principle itself
understood.  It may be mentioned here, that this dynamical
principle is only another form of that idea which has already
been applied to optics in the {\it Theory of systems of rays},
and that an intention of applying it to the motion of systems of
bodies was announced\footnote\dag{Transactions of the Royal
Irish Academy, Vol.~{\sc xv}, page~80.  A notice of this dynamical
principle was also lately given in an article ``On a general
Method of expressing the Paths of Light and of the Planets,''
published in the Dublin University Review for October 1833.}
at the publication of that theory.  And besides the idea itself,
the manner of calculation also, which has been thus exemplified
in the sciences of optics and dynamics, seems not confined to
those two sciences, but capable of other applications; and the
peculiar combination which it involves, of the principles of
variations with those of partial differentials, for the
determination and use of an important class of integrals, may
constitute, when it shall be matured by the future labours of
mathematicians, a separate branch of analysis.

\nobreak\medskip

\line{\hfil WILLIAM R. HAMILTON.}

\nobreak\bigskip

{\it Observatory, Dublin, March\/} 1834.

\vfill\eject

{\sectiontitle
Integration of the Equations of Motion of a System, characteristic
Function of such Motion, and Law of varying Action.\par}

\nobreak\bigskip

1.
The known differential equations of motion of a system of free
points, repelling or attracting one another according to any
functions of their distances, and not disturbed by any foreign
force, may be comprised in the following formula:
$$\sum \mathbin{.} m (x'' \, \delta x + y'' \, \delta y + z'' \, \delta z)
   = \delta U.
   \eqno {\rm (1.)}$$
In this formula the sign of summation $\sum$ extends to all the
points of the system; $m$ is, for any one such point, the
constant called its mass; $x''$, $y''$, $z''$, are its component
accelerations, or the second differential coefficients of its
rectangular coordinates $x$, $y$, $z$, taken with respect to the
time; $\delta x$, $\delta y$, $\delta z$, are any arbitrary
infinitesimal displacements which the point can be imagined to
receive in the same three rectangular directions; and $\delta U$
is the infinitesimal variation corresponding, of a function~$U$
of the masses and mutual distances of the several points of the
system, of which the form depends on the laws of their mutual
actions, by the equation
$$U = \sum \mathbin{.} m m_\prime f(r),
   \eqno {\rm (2.)}$$
$r$ being the distance between any two points $m$, $m_\prime$,
and the function $f(r)$ being such that the derivative or
differential coefficient $f'(r)$ expresses the law of their
repulsion, being negative in the case of attraction.  The
function which has been here called $U$ may be named the
{\it force-function\/} of a system: it is of great utility in
theoretical mechanics, into which it was introduced by
{\sc Lagrange}, and it furnishes the following elegant forms
for the differential equations of motion, included in the
formula (1.):
$$\left. \multieqalign{
m_1 x_1'' &= {\delta U \over \delta x_1}; &
m_2 x_2'' &= {\delta U \over \delta x_2}; \quad \ldots&
m_n x_n'' &= {\delta U \over \delta x_n}; \cr
m_1 y_1'' &= {\delta U \over \delta y_1}; &
m_2 y_2'' &= {\delta U \over \delta y_2}; \quad \ldots&
m_n y_n'' &= {\delta U \over \delta y_n}; \cr
m_1 z_1'' &= {\delta U \over \delta z_1}; &
m_2 z_2'' &= {\delta U \over \delta z_2}; \quad \ldots&
m_n z_n'' &= {\delta U \over \delta z_n}; \cr}
   \right\}
   \eqno {\rm (3.)}$$
the second members of these equations being the partial
differential coefficients of the first order of the function~$U$.
But notwithstanding the elegance and simplicity of this known
manner of stating the principal problem of dynamics, the
difficulty of solving that problem, or even of expressing its
solution, has hitherto appeared insuperable; so that only seven
intermediate integrals, or integrals of the first order, with as
many arbitrary constants, have hitherto been found for these
general equations of motion of a system of $n$ points, instead of
$3n$ intermediate and $3n$ final integrals, involving ultimately
$6n$ constants; nor has any integral been
found which does not need to be integrated again.  No general
solution has been obtained assigning (as a complete solution
ought to do) $3n$ relations between the $n$ masses
$m_1, m_2,\ldots \, m_n$, the $3n$ varying coordinates
$x_1, y_1, z_1,\ldots\, x_n, y_n, z_n$, the varying time~$t$,
and the $6n$ initial data of the problem, namely, the initial
coordinates $a_1, b_1, c_1,\ldots\, a_n, b_n, c_n$, and their
initial rates of increase
$a_1', b_1', c_1',\ldots\, a_n', b_n', c_n'$; the quantities
called here initial being those which correspond to the arbitrary
origin of time.  It is, however, possible (as we shall see) to
express these long-sought relations by the partial differential
coefficients of a new central or radical function, to the search
and employment of which the difficulty of mathematical dynamics
becomes henceforth reduced.

\bigbreak

2.
If we put for abridgement
$$T = {\textstyle {1 \over 2}} \sum \mathbin{.} m
         (x'^2 + y'^2 + z'^2),
   \eqno {\rm (4.)}$$
so that $2T$ denotes, as in the M\'{e}canique Analytique,
the whole living force of the system; ($x'$,~$y'$,~$z'$, being
here, according to the analogy of our foregoing notation, the
rectangular components of velocity of the point~$m$, or the first
differential coefficients of its coordinates taken with respect to
the time;) an easy and well known combination of the
differential equations of motion, obtained by changing in the
formula (1.) the variations to the differentials of the
coordinates, may be expressed in the following manner,
$$dT = dU,
   \eqno {\rm (5.)}$$
and gives, by integration, the celebrated law of living force,
under the form
$$T = U + H.
   \eqno {\rm (6.)}$$

In this expression, which is one of the seven known integrals
already mentioned, the quantity~$H$ is independent of the time,
and does not alter in the passage of the points of the system
from one set of positions to another.  We have, for example, an
initial equation of the same form, corresponding to the origin of
time, which may be written thus,
$$T_0 = U_0 + H.
   \eqno {\rm (7.)}$$

The quantity~$H$ may, however, receive any arbitrary increment
whatever, when we pass in thought from a system moving in one
way, to the same system moving in another, with the same
dynamical relations between the accelerations and positions of
its points, but with different initial data; but the increment of
$H$, thus obtained, is evidently connected with the analogous
increments of the functions $T$ and $U$, by the relation
$$\Delta T = \Delta U + \Delta H,
   \eqno {\rm (8.)}$$
which, for the case of infinitesimal variations, may be
conveniently be written thus,
$$\delta T = \delta U + \delta H;
   \eqno {\rm (9.)}$$
and this last relation, when multiplied by $dt$, and integrated,
conducts to an important result.  For it thus becomes, by (4.)
and (1.),
$$\int \sum \mathbin{.}
     m (    dx \mathbin{.} \delta x'
          + dy \mathbin{.} \delta y'
          + dz \mathbin{.} \delta z' )
   = \int \sum \mathbin{.}
     m (    dx' \mathbin{.} \delta x
          + dy' \mathbin{.} \delta y
          + dz' \mathbin{.}  \delta z )
      + \int \delta H \mathbin{.} dt,
   \eqno {\rm (10.)}$$
that is, by the principles of the calculus of variations,
$$\delta V
   =   \sum \mathbin{.} m (x' \, \delta x + y' \, \delta y + z' \, \delta z)
     - \sum \mathbin{.} m (a' \, \delta a + b' \, \delta b + c' \, \delta c)
     + t \, \delta H,
   \eqno {\rm (A.)}$$
if we denote by $V$ the integral
$$V = \int \sum \mathbin{.} m ( x' \,dx + y' \,dy + z' \,dz )
   = \int_0^t 2 T \,dt,
   \eqno {\rm (B.)}$$
namely, the accumulated living force, called often the action of
the system, from its initial to its final position.

If, then, we consider (as it is easy to see that we may) the
action~$V$ as a function of the initial and final coordinates,
and of the quantity~$H$, we shall have, by (A.), the following
groups of equations; first, the group,
$$\left. \multieqalign{
{\delta V \over \delta x_1} &= m_1 x_1'; &
{\delta V \over \delta x_2} &= m_2 x_2'; \quad \ldots &
{\delta V \over \delta x_n} &= m_n x_n'; \cr
{\delta V \over \delta y_1} &= m_1 y_1'; &
{\delta V \over \delta y_2} &= m_2 y_2'; \quad \ldots &
{\delta V \over \delta y_n} &= m_n y_n'; \cr
{\delta V \over \delta z_1} &= m_1 z_1'; &
{\delta V \over \delta z_2} &= m_2 z_2'; \quad \ldots &
{\delta V \over \delta z_n} &= m_n z_n'. \cr}
   \right\}
   \eqno {\rm (C.)}$$

Secondly, the group,
$$\left. \multieqalign{
{\delta V \over \delta a_1} &= - m_1 a_1'; &
{\delta V \over \delta a_2} &= - m_2 a_2'; \quad \ldots &
{\delta V \over \delta a_n} &= - m_n a_n'; \cr
{\delta V \over \delta b_1} &= - m_1 b_1'; &
{\delta V \over \delta b_2} &= - m_2 b_2'; \quad \ldots &
{\delta V \over \delta b_n} &= - m_n b_n'; \cr
{\delta V \over \delta c_1} &= - m_1 c_1'; &
{\delta V \over \delta c_2} &= - m_2 c_2'; \quad \ldots &
{\delta V \over \delta c_n} &= - m_n c_n'; \cr}
   \right\}
   \eqno {\rm (D.)}$$
and finally, the equation,
$${\delta V \over \delta H} = t.
   \eqno {\rm (E.)}$$
So that if this function~$V$ were known, it would only remain
to eliminate $H$ between the $3n + 1$ equations (C.) and (E.),
in order to obtain all the $3n$ intermediate integrals, or
between (D.) and (E.) to obtain all the $3n$ final integrals of
the differential equations of motion; that is, ultimately, to
obtain the $3n$ sought relations between the $3n$ varying
coordinates and the time, involving also the masses and the $6n$
initial data above mentioned; the discovery of which relations
would be (as we have said) the general solution of the general
problem of dynamics.  We have, therefore, at least reduced that
general problem to the search and differentiation of a single
function~$V$, which we shall call on this account the
{\sc characteristic function} of motion of a system; and the
equation (A.), expressing the fundamental law of its variation,
we shall call the {\it equation of the characteristic function},
or the {\sc law of varying action}.

\medskip

3.
To show more clearly that the action or accumulated living force
of a system, or in other words, the integral of the product of
the living force by the element of the time, may be regarded as
a function of the $6n + 1$ quantities already mentioned, namely,
of the initial and final coordinates, and of the quantity~$H$,
we may observe, that whatever depends on the manner and time of
motion of the system may be considered as such a function;
because the initial form of the law of living force, when
combined with the $3n$ known or unknown relations between the
time, the initial data, and the varying coordinates, will always
furnish $3n + 1$ relations, known or unknown, to connect the time
and the initial components of velocities with the initial and
final coordinates, and with $H$.  Yet from not having formed the
conception of the action as a {\it function\/} of this kind, the
consequences that have been here deduced from the formula (A.)
for the variation of that definite integral appear to have
escaped the notice of {\sc Lagrange}, and of the other illustrious
analysts who have written on theoretical mechanics; although they
were in possession of a formula for the variation of this
integral not greatly differing from ours.  For although {\sc Lagrange}
and others, in treating of the motion of a system, have shown
that the variation of this definite integral vanishes when the
extreme coordinates and the constant $H$ are given, they appear
to have deduced from this result only the well known law of
{\it least action\/}; namely, that if the points or bodies of a
system be imagined to move from a given set of initial to a given
set of final positions, not as they do nor even as they could
move consistently with the general dynamical laws or differential
equations of motion, but so as not to violate any supposed
geometrical connexions, nor that one dynamical relation between
velocities and configurations which constitutes the law of
living force; and if, besides, this geometrically imaginable, but
dynamically impossible motion, be made to differ infinitely
{\it little\/} from the actual manner of motion of the system,
between the given extreme positions; then the varied value of the
definite integral called action, or the accumulated living force
of the system in the motion thus imagined, will differ infinitely
{\it less\/} from the actual value of that integral.  But when
this well known law of least, or as it might be better called, of
{\it stationary action}, is applied to the determination of the
actual motion of the system, it serves only to form, by the rules
of the calculus of variations, the differential equations of
motion of the second order, which can always be otherwise found.
It seems, therefore, to be with reason that {\sc Lagrange},
{\sc Laplace}, and {\sc Poisson} have spoken lightly of the utility
of this principle in the present state of dynamics.  A different
estimate, perhaps, will be formed of that other principle which
has been introduced in the present paper, under the name of the
{\it law of varying action}, in which we pass from an actual
motion to another motion dynamically possible, by varying the
extreme positions of the system, and (in general) the
quantity~$H$, and which serves to express, by means of a single
function, not the mere differential equations of motion, but
their intermediate and their final integrals.

\bigbreak

{\sectiontitle
Verification of the foregoing Integrals.\par}

\nobreak\bigskip

4.
A verification, which ought not to be neglected, and at the same
time an illustration of this new principle, may be obtained by
deducing the known differential equations of motion from our
system of intermediate integrals, and by showing the consistence
of these again with our final integral system.  As preliminary to
such verification, it is useful to observe that the final
equation (6.) of living force, when combined with the system
(C.), takes this new form,

\vfill\eject  % Page break necessary with current page size

$${\textstyle {1 \over 2}} \sum \mathbin{.} {1 \over m} \left\{
        \left( {\delta V \over \delta x} \right)^2
      + \left( {\delta V \over \delta y} \right)^2
      + \left( {\delta V \over \delta z} \right)^2
      \right\} = U + H;
   \eqno {\rm (F.)}$$
and that the initial equation (7.) of living force becomes by
(D.)
$${\textstyle {1 \over 2}} \sum \mathbin{.} {1 \over m} \left\{
        \left( {\delta V \over \delta a} \right)^2
      + \left( {\delta V \over \delta b} \right)^2
      + \left( {\delta V \over \delta c} \right)^2
      \right\} = U_0 + H.
   \eqno {\rm (G.)}$$
These two partial differential equations, inital and final, of
the first order and the second degree, must both be identically
satisfied by the characteristic function~$V$: they furnish (as we
shall find) the principal means of discovering the form of that
function, and are of essential importance in its theory.  If the
form of this function were known, we might eliminate $3n - 1$ of
the $3n$ initial coordinates between the $3n$ equations (C.); and
although we cannot yet perform the actual process of this
elimination, we are entitled to assert that it would remove along
with the others the remaining initial coordinate, and would conduct
to the equation (6.) of final living force, which might then be
transformed into the equation (F.).  In like manner we may
conclude that all the $3n$ final coordinates could be eliminated
together from the $3n$ equations (D.), and that the result would
be the initial equation (7.) of living force, or the transformed
equation (G.).  We may therefore consider the law of living
force, which assisted us in discovering the properties of our
characteristic function~$V$, as included reciprocally in those
properties, and as resulting by elimination, in every particular
case, from the systems (C.) and (D.); and in treating of either
of these systems, or in conducting any other dynamical
investigation by the method of this characteristic function, we
are at liberty to employ the partial differential equations (F.)
and (G.) which that function must necessarily satisfy.

It will now be easy to deduce, as we proposed, the known
equations of motion (3.) of the second order, by differentiation
and elimination of constants, from our intermediate integral
system (C.), (E.), or even from a part of that system, namely,
from the group (C.), when combined with the equation (F.).  For
we thus obtain
$$\left. \eqalign{
m_1 x_1'' = {d \over dt} {\delta V \over \delta x_1}
   &=   x_1' {\delta^2 V \over \delta x_1^2}
      + x_2' {\delta^2 V \over \delta x_1 \, \delta x_2}
      + \cdots
      + x_n' {\delta^2 V \over \delta x_1 \, \delta x_n} \cr
   &\mathrel{\phantom{=}} \mathord{}
      + y_1' {\delta^2 V \over \delta x_1 \, \delta y_1}
      + y_2' {\delta^2 V \over \delta x_1 \, \delta y_2}
      + \cdots
      + y_n' {\delta^2 V \over \delta x_1 \, \delta y_n} \cr
   &\mathrel{\phantom{=}} \mathord{}
      + z_1' {\delta^2 V \over \delta x_1 \, \delta z_1}
      + z_2' {\delta^2 V \over \delta x_1 \, \delta z_2}
      + \cdots
      + z_n' {\delta^2 V \over \delta x_1 \, \delta z_n} \cr
   &=   {1 \over m_1} {\delta V \over \delta x_1}
            {\delta^2 V \over \delta x_1^2}
      + {1 \over m_2} {\delta V \over \delta x_2}
            {\delta^2 V \over \delta x_1 \, \delta x_2}
      + \cdots
      + {1 \over m_n} {\delta V \over \delta x_n}
            {\delta^2 V \over \delta x_1 \, \delta x_n} \cr
   &\mathrel{\phantom{=}} \mathord{}
      + {1 \over m_1} {\delta V \over \delta y_1}
            {\delta^2 V \over \delta x_1 \, \delta y_1}
      + {1 \over m_2} {\delta V \over \delta y_2}
            {\delta^2 V \over \delta x_1 \, \delta y_2}
      + \cdots
      + {1 \over m_n} {\delta V \over \delta y_n}
            {\delta^2 V \over \delta x_1 \, \delta y_n} \cr
   &\mathrel{\phantom{=}} \mathord{}
      + {1 \over m_1} {\delta V \over \delta z_1}
            {\delta^2 V \over \delta x_1 \, \delta z_1}
      + {1 \over m_2} {\delta V \over \delta z_2}
            {\delta^2 V \over \delta x_1 \, \delta z_2}
      + \cdots
      + {1 \over m_n} {\delta V \over \delta z_n}
            {\delta^2 V \over \delta x_1 \, \delta z_n} \cr
   &=  {\delta \over \delta x_1} \sum \mathbin{.} {1 \over 2m} \left\{
        \left( {\delta V \over \delta x} \right)^2
      + \left( {\delta V \over \delta y} \right)^2
      + \left( {\delta V \over \delta z} \right)^2
      \right\}
      = {\delta \over \delta x_1} (U + H);\cr}
   \right\}
   \eqno {\rm (11.)}$$
that is, we obtain
$$m_1 x_1'' = {\delta U \over \delta x_1}.
   \eqno {\rm (12.)}$$
And in like manner we might deduce, by differentiation, from the
integrals (C.) and from (F.) all the other known differential
equations of motion, of the second order, contained in the set
marked (3.); or, more concisely, we may deduce at once the formula
(1.), which contains all those known equations, by observing that
the intermediate integrals (C.), when combined with the relation
(F.), give
$$\left. \eqalign{
\sum \mathbin{.} m (x'' \, \delta x + y'' \, \delta y + z'' \, \delta z)
  \hskip-10em & \cr
   &= \sum \left(
        {d \over dt} {\delta V \over \delta x} \mathbin{.} \delta x
      + {d \over dt} {\delta V \over \delta y} \mathbin{.} \delta y
      + {d \over dt} {\delta V \over \delta z} \mathbin{.} \delta z
      \right) \cr
   &= \sum \mathbin{.} {1 \over m} \left(
        {\delta V \over \delta x} {\delta \over \delta x}
      + {\delta V \over \delta y} {\delta \over \delta y}
      + {\delta V \over \delta z} {\delta \over \delta z}
      \right)
      \sum \left(
        {\delta V \over \delta x} \delta x
      + {\delta V \over \delta y} \delta y
      + {\delta V \over \delta z} \delta z
      \right) \cr
   &= \sum \left(
        \delta x \, {\delta \over \delta x}
      + \delta y \, {\delta \over \delta y}
      + \delta z \, {\delta \over \delta z}
      \right)
      \sum \mathbin{.} {1 \over 2m} \left\{
        \left( {\delta V \over \delta x} \right)^2
      + \left( {\delta V \over \delta y} \right)^2
      + \left( {\delta V \over \delta z} \right)^2
      \right\} \cr
   &= \sum \left(
        \delta x \, {\delta \over \delta x}
      + \delta y \, {\delta \over \delta y}
      + \delta z \, {\delta \over \delta z}
      \right) (U + H) \cr
   &= \delta U.\cr}
   \right\}
   \eqno {\rm (13.)}$$

\bigbreak

5.
Again, we were to show that our intermediate integral system,
composed of the equations (C.) and (E.), with the $3n$ arbitrary
constants $a_1, b_1, c_1,\ldots\,a_n, b_n, c_n$, (and involving
also the auxiliary constant $H$,) is consistent with our final
integral system of equations (D.) and (E.), which contain $3n$
other arbitrary constants, namely
$a_1', b_1', c_1',\ldots\, a_n', b_n', c_n'$.  The immediate
differentials of the equations  (C.), (D.), (E.), taken with
respect to the time, are, for the first group,
$$\left. \multieqalign{
{d \over dt} {\delta V \over \delta x_1} &= m_1 x_1''; &
{d \over dt} {\delta V \over \delta x_2} &= m_2 x_2''; \quad \ldots &
{d \over dt} {\delta V \over \delta x_n} &= m_n x_n''; \cr
{d \over dt} {\delta V \over \delta y_1} &= m_1 y_1''; &
{d \over dt} {\delta V \over \delta y_2} &= m_2 y_2''; \quad \ldots &
{d \over dt} {\delta V \over \delta y_n} &= m_n y_n''; \cr
{d \over dt} {\delta V \over \delta z_1} &= m_1 z_1''; &
{d \over dt} {\delta V \over \delta z_2} &= m_2 z_2''; \quad \ldots &
{d \over dt} {\delta V \over \delta z_n} &= m_n z_n''; \cr}
   \right\}
   \eqno {\rm (H.)}$$
for the second group,
$$\left. \multieqalign{
{d \over dt} {\delta V \over \delta a_1} &= 0; &
{d \over dt} {\delta V \over \delta a_2} &= 0; \quad \ldots &
{d \over dt} {\delta V \over \delta a_n} &= 0; \cr
{d \over dt} {\delta V \over \delta b_1} &= 0; &
{d \over dt} {\delta V \over \delta b_2} &= 0; \quad \ldots &
{d \over dt} {\delta V \over \delta b_n} &= 0; \cr
{d \over dt} {\delta V \over \delta c_1} &= 0; &
{d \over dt} {\delta V \over \delta c_2} &= 0; \quad \ldots &
{d \over dt} {\delta V \over \delta c_n} &= 0; \cr}
   \right\}
   \eqno {\rm (I.)}$$
and finally, for the last equation,
$${d \over dt} {\delta V \over \delta H} = 1.
   \eqno {\rm (K.)}$$

By combining the equations (C.) with their differentials (H.),
and with the relation (F.), we deduced, in the foregoing number,
the known equations of motion (3.); and we are now to show the
consistence of the same intermediate integrals (C.) with the
group of differentials (I.) which have been obtained from the
final integrals.

The first equation of the group (I.) may be developed thus:
$$\left. \eqalign{0
   &=   x_1' {\delta^2 V \over \delta a_1 \, \delta x_1}
      + x_2' {\delta^2 V \over \delta a_1 \, \delta x_2}
      + \cdots
      + x_n' {\delta^2 V \over \delta a_1 \, \delta x_n} \cr
   &\mathrel{\phantom{=}} \mathord{}
      + y_1' {\delta^2 V \over \delta a_1 \, \delta y_1}
      + y_2' {\delta^2 V \over \delta a_1 \, \delta y_2}
      + \cdots
      + y_n' {\delta^2 V \over \delta a_1 \, \delta y_n} \cr
   &\mathrel{\phantom{=}} \mathord{}
      + z_1' {\delta^2 V \over \delta a_1 \, \delta z_1}
      + z_2' {\delta^2 V \over \delta a_1 \, \delta z_2}
      + \cdots
      + z_n' {\delta^2 V \over \delta a_1 \, \delta z_n} \cr}
   \right\}
   \eqno {\rm (14.)}$$
and the others may be similarly developed.  In order, therefore,
to show that they are satisfied by the group (C.), it is
sufficient to prove that the following equations are true,
$$\left. \eqalign{
0 &= {\delta \over \delta a_i} \sum \mathbin{.} {1 \over 2m} \left\{
        \left( {\delta V \over \delta x} \right)^2
      + \left( {\delta V \over \delta y} \right)^2
      + \left( {\delta V \over \delta z} \right)^2
      \right\},\cr
0 &= {\delta \over \delta b_i} \sum \mathbin{.} {1 \over 2m} \left\{
        \left( {\delta V \over \delta x} \right)^2
      + \left( {\delta V \over \delta y} \right)^2
      + \left( {\delta V \over \delta z} \right)^2
      \right\},\cr
0 &= {\delta \over \delta c_i} \sum \mathbin{.} {1 \over 2m} \left\{
        \left( {\delta V \over \delta x} \right)^2
      + \left( {\delta V \over \delta y} \right)^2
      + \left( {\delta V \over \delta z} \right)^2
      \right\},\cr}
   \right\}
   \eqno {\rm (L.)}$$
the integer~$i$ receiving any value from $1$ to $n$ inclusive;
which may be shown at once, and the required verification thereby
be obtained, if we merely take the variation of the relation
(F.) with respect to the initial coordinates, as in the former
verification we took its variation with respect to the final
coordinates, and so obtained results which agreed with the known
equations of motion, and which may be thus collected,
$$\left. \eqalign{
{\delta \over \delta x_i} \sum \mathbin{.} {1 \over 2m} \left\{
        \left( {\delta V \over \delta x} \right)^2
      + \left( {\delta V \over \delta y} \right)^2
      + \left( {\delta V \over \delta z} \right)^2
      \right\}
   &= {\delta U \over \delta x_i};\cr
{\delta \over \delta y_i} \sum \mathbin{.} {1 \over 2m} \left\{
        \left( {\delta V \over \delta x} \right)^2
      + \left( {\delta V \over \delta y} \right)^2
      + \left( {\delta V \over \delta z} \right)^2
      \right\}
   &= {\delta U \over \delta y_i};\cr
{\delta \over \delta z_i} \sum \mathbin{.} {1 \over 2m} \left\{
        \left( {\delta V \over \delta x} \right)^2
      + \left( {\delta V \over \delta y} \right)^2
      + \left( {\delta V \over \delta z} \right)^2
      \right\}
   &= {\delta U \over \delta z_i}.\cr}
   \right\}
   \eqno {\rm (M.)}$$

The same relation (F.), by being varied with respect to the
quantity $H$, conducts to the expression
$${\delta \over \delta H} \sum \mathbin{.} {1 \over 2m} \left\{
        \left( {\delta V \over \delta x} \right)^2
      + \left( {\delta V \over \delta y} \right)^2
      + \left( {\delta V \over \delta z} \right)^2
      \right\}
   = 1;
   \eqno {\rm (N.)}$$
and this, when developed, agrees with the equation (K.), which
is a new verification of the consistence of our foregoing
results.  Nor would it have been much more difficult, by the help
of the foregoing principles, to have integrated directly our
integrals of the first order, and so to have deduced in a
different way our final integral system.

\bigbreak

6.
It may be considered as still another verification of our own
general integral equations, to show that they include not only
the known law of living force, or the integral expressing that
law, but also the six other known integrals of the first order,
which contain the law of motion of the centre of gravity, and the
law of description of areas.  For this purpose, it is only
necessary to observe that it evidently follows from the
conception of our characteristic function $V$, that the function
depends on the initial and final positions of the attracting or
repelling points of a system, not as referred to any foreign
standard, but only as compared with one another; and therefore
that this function will not vary, if without making any real
change in either initial or final configuration, or in the
relation of these to each other, we alter at once all the initial
and all the final positions of the points of the system, by any
common motion, whether of translation or of rotation.  Now by
considering these coordinate translations, we obtain the three
following partial differential equations of the first order,
which the function $V$ must satisfy,
$$\left. \eqalign{
\sum {\delta V \over \delta x} + \sum {\delta V \over \delta a}
   &= 0;\cr
\sum {\delta V \over \delta y} + \sum {\delta V \over \delta b}
   &= 0;\cr
\sum {\delta V \over \delta z} + \sum {\delta V \over \delta c}
   &= 0;\cr}
   \right\}
   \eqno {\rm (O.)}$$
and by considering three coordinate rotations, we obtain these
three other relations between the partial differential
coefficients of the same order of the same characteristic
function,
$$\left. \eqalign{
         \sum \left( x {\delta V \over \delta y}
                   - y {\delta V \over \delta x} \right)
       + \sum \left( a {\delta V \over \delta b}
                   - b {\delta V \over \delta a} \right)
   &= 0;\cr
         \sum \left( y {\delta V \over \delta z}
                   - z {\delta V \over \delta y} \right)
       + \sum \left( b {\delta V \over \delta c}
                   - c {\delta V \over \delta b} \right)
   &= 0;\cr
         \sum \left( z {\delta V \over \delta x}
                   - x {\delta V \over \delta z} \right)
       + \sum \left( c {\delta V \over \delta a}
                   - a {\delta V \over \delta c} \right)
   &= 0;\cr}
   \right\}
   \eqno {\rm (P.)}$$
and if we change the final coefficients of $V$ to the final
components of momentum, and the initial coefficients to the
initial components taken negatively, according to the dynamical
properties of this function expressed by the integrals (C.) and
(D.), we shall change these partial differential equations (O.)
(P.), to the following,
$$\sum \mathbin{.} m x' = \sum \mathbin{.} m a';\quad
  \sum \mathbin{.} m y' = \sum \mathbin{.} m b';\quad
  \sum \mathbin{.} m z' = \sum \mathbin{.} m c';
   \eqno {\rm (15.)}$$
and
$$\left. \eqalign{
\sum \mathbin{.} m (x y' - y x') &= \sum \mathbin{.} m (a b' - b a');\cr
\sum \mathbin{.} m (y z' - z y') &= \sum \mathbin{.} m (b c' - c b');\cr
\sum \mathbin{.} m (z x' - x z') &= \sum \mathbin{.} m (c a' - a c').\cr}
   \right\}
   \eqno {\rm (16.)}$$

In this manner, therefore, we can deduce from the properties of
our characteristic function the six other known integrals above
mentioned, in addition to that seventh which contains the law of
living force, and which assisted in the discovery of our method.

\bigbreak

{\sectiontitle
Introduction of relative or polar Coordinates, or other marks of
position of a System.\par}

\nobreak\bigskip

7.
The property of our characteristic function, by which it depends
only on the internal or mutual relations between the positions
initial and final of the points of an attracting or repelling
system, suggests an advantage in employing internal or relative
coordinates; and from the analogy of other applications of
algebraical methods to researches of a geometrical kind, it may
be expected that polar and other marks of position will also
often be found useful.  Supposing, therefore, that the $3n$ final
coordinates $x_1 \, y_1 \, z_1 \,\ldots\, x_n \, y_n \, z_n$
have been expressed as functions of $3n$ other variables
$\eta_1 \, \eta_2 \,\ldots\, \eta_{3n}$, and that the $3n$ initial
coordinates have in like manner been expressed as functions of
$3n$ similar quantities, which we shall call
$e_1 \, e_2 \, \ldots\, e_{3n}$, we shall proceed to assign a general
method for introducing these new marks of position into the
expressions of our fundamental relations.

For this purpose we have only to transform the law of varying
action, or the fundamental formula (A.), by transforming the two
sums,
$$\sum \mathbin{.} m (x' \, \delta x + y' \, \delta y + z' \, \delta z),
      \quad \hbox{and} \quad
  \sum \mathbin{.} m (a' \, \delta a + b' \, \delta b + c' \, \delta c),$$
which it involves, and which are respectively equivalent to the
following more developed expressions,
$$\left. \eqalign{
\sum \mathbin{.} m (x' \, \delta x + y' \, \delta y + z' \, \delta z)
   = m_1
      & ( x_1' \, \delta x_1 + y_1' \, \delta y_1 + z_1' \, \delta z_1 ) \cr
   \mathord{}
   + m_2
      & ( x_2' \, \delta x_2 + y_2' \, \delta y_2 + z_2' \, \delta z_2 ) \cr
   \mathord{}
   + \hbox{\&c.} + m_n
      & ( x_n' \, \delta x_n + y_n' \, \delta y_n + z_n' \, \delta z_n );\cr}
   \right\}
   \eqno {\rm (17.)}$$
$$\left. \eqalign{
\sum \mathbin{.} m (a' \, \delta a + b' \, \delta b + c' \, \delta c)
   = m_1
      & ( a_1' \, \delta a_1 + b_1' \, \delta b_1 + c_1' \, \delta c_1 ) \cr
   \mathord{}
   + m_2
      & ( a_2' \, \delta a_2 + b_2' \, \delta b_2 + c_2' \, \delta c_2 ) \cr
   \mathord{}
   + \hbox{\&c.} + m_n
      & ( a_n' \, \delta a_n + b_n' \, \delta b_n + c_n' \, \delta c_n ).\cr}
   \right\}
   \eqno {\rm (18.)}$$
Now $x_i$ being by supposition a function of the $3n$ new marks
of position $\eta_1 \, \ldots\, \eta_{3n}$, its variation
$\delta x_i$, and its differential coefficient $x_i'$ may be thus
expressed:
$$\delta x_i =   {\delta x_i \over \delta \eta_1} \delta \eta_1
               + {\delta x_i \over \delta \eta_2} \delta \eta_2
               + \cdots
               + {\delta x_i \over \delta \eta_{3n}} \delta \eta_{3n};
   \eqno {\rm (19.)}$$
$$x_i'       =   {\delta x_i \over \delta \eta_1} \eta_1'
               + {\delta x_i \over \delta \eta_2} \eta_2'
               + \cdots
               + {\delta x_i \over \delta \eta_{3n}} \eta_{3n}';
   \eqno {\rm (20.)}$$
and similarly for $y_i$ and $z_i$.  If, then, we consider $x_i'$
as a function, by (20.), of $\eta_1' \, \ldots\, \eta_{3n}'$,
involving also in general $\eta_1 \, \ldots\, \eta_{3n}$, and if we
take its partial differential coefficients of the first order
with respect to $\eta_1' \,\ldots\, \eta_{3n}'$, we find the
relations,
$${\delta x_i' \over \delta \eta_1'}
   = {\delta x_i \over \delta \eta_1};\quad
  {\delta x_i' \over \delta \eta_2'}
   = {\delta x_i \over \delta \eta_2};\quad \ldots \quad
  {\delta x_i' \over \delta \eta_{3n}'}
   = {\delta x_i \over \delta \eta_{3n}};
   \eqno {\rm (21.)}$$
and therefore we obtain these new expressions for the variations
$\delta x_i$, $\delta y_i$, $\delta z_i$,
$$\left. \eqalign{
\delta x_i =   {\delta x_i' \over \delta \eta_1'} \delta \eta_1
             + {\delta x_i' \over \delta \eta_2'} \delta \eta_2
             + \cdots
             + {\delta x_i' \over \delta \eta_{3n}'} \delta \eta_{3n},\cr
\delta y_i =   {\delta y_i' \over \delta \eta_1'} \delta \eta_1
             + {\delta y_i' \over \delta \eta_2'} \delta \eta_2
             + \cdots
             + {\delta y_i' \over \delta \eta_{3n}'} \delta \eta_{3n},\cr
\delta z_i =   {\delta z_i' \over \delta \eta_1'} \delta \eta_1
             + {\delta z_i' \over \delta \eta_2'} \delta \eta_2
             + \cdots
             + {\delta z_i' \over \delta \eta_{3n}'} \delta \eta_{3n}.\cr}
   \right\}
   \eqno {\rm (22.)}$$

Substituting these expressions (22.) for the variations in the
sum (17.), we easily transform it into the following,
$$\left. \eqalign{
\sum \mathbin{.} m (x' \, \delta x + y' \, \delta y + z' \, \delta z)
   = &\sum \mathbin{.} m \left(   x' {\delta x' \over \delta \eta_1'}
                      + y' {\delta y' \over \delta \eta_1'}
                      + z' {\delta z' \over \delta \eta_1'}
                      \right) \mathbin{.} \delta \eta_1 \cr
   + &\sum \mathbin{.} m \left(   x' {\delta x' \over \delta \eta_2'}
                      + y' {\delta y' \over \delta \eta_2'}
                      + z' {\delta z' \over \delta \eta_2'}
                      \right) \mathbin{.} \delta \eta_2 \cr
   \mathord{}
   + \hbox{\&c.}
   + &\sum \mathbin{.} m \left(   x' {\delta x' \over \delta \eta_{3n}'}
                      + y' {\delta y' \over \delta \eta_{3n}'}
                      + z' {\delta z' \over \delta \eta_{3n}'}
                      \right) \mathbin{.} \delta \eta_{3n} \cr
   = &  {\delta T \over \delta \eta_1'} \delta \eta_1
      + {\delta T \over \delta \eta_2'} \delta \eta_2
      + \cdots
      + {\delta T \over \delta \eta_{3n}'} \delta \eta_{3n};\cr}
   \right\}
   \eqno {\rm (23.)}$$
$T$ being the same quantity as before, namely, the half of the
final living force of system, but being now considered as a
function of $\eta_1' \,\ldots\, \eta_{3n}'$, involving also the
masses, and in general $\eta_1 \,\ldots\, \eta_{3n}$, and obtained
by substituting for the quantities $x'$~$y'$~$z'$ their values
of the form (20.) in the equation of definition
$$T = {\textstyle {1 \over 2}} \sum \mathbin{.} m
         (x'^2 + y'^2 + z'^2).
   \eqno {\rm (4.)}$$

In like manner we find this transformation for the sum (18.),
$$\sum \mathbin{.} m (a' \, \delta a + b' \, \delta b + c' \, \delta c)
   =   {\delta T_0 \over \delta e_1'} \delta e_1
     + {\delta T_0 \over \delta e_2'} \delta e_2
     + \cdots
     + {\delta T_0 \over \delta e_{3n}'} \delta e_{3n}.
   \eqno {\rm (24.)}$$

The law of varying action, or the formula (A.), becomes
therefore, when expressed by the present more general coordinates
or marks of position,
$$\delta V =   \sum \mathbin{.} {\delta T \over \delta \eta'} \delta \eta
             - \sum \mathbin{.} {\delta T \over \delta e'} \delta e
             + t \, \delta H;
   \eqno {\rm (Q.)}$$
and instead of the groups (C.) and (D.), into which, along with
the equation (E.), this law resolved itself before, it gives now
these other groups,
$${\delta V \over \delta \eta_1}
   = {\delta T \over \delta \eta_1'}; \quad
  {\delta V \over \delta \eta_2}
   = {\delta T \over \delta \eta_2'}; \quad \cdots \quad
  {\delta V \over \delta \eta_{3n}}
   = {\delta T \over \delta \eta_{3n}'};
   \eqno {\rm (R.)}$$
and
$${\delta V \over \delta e_1}
   = - {\delta T_0 \over \delta e_1'}; \quad
  {\delta V \over \delta e_2}
   = - {\delta T_0 \over \delta e_2'}; \quad \cdots \quad
  {\delta V \over \delta e_{3n}}
   = - {\delta T_0 \over \delta e_{3n}'}.
   \eqno {\rm (S.)}$$

The quantities $e_1 \, e_2 \,\ldots\, e_{3n}$, and
$e_1' \, e_2' \,\ldots\, e_{3n}'$, are now the initial data
respecting the manner of motion of the system; and the $3n$ final
integrals, connecting these $6n$ initial data, and the $n$
masses, with the time~$t$, and with the $3n$ final or varying
quantities $\eta_1 \, \eta_2 \,\ldots\, \eta_{3n}$, which mark the
varying positions of the $n$ moving points of the system, are now
to be obtained by eliminating the auxiliary constant $H$ between
the $3n + 1$ equations (S.) and (E.); while the $3n$ intermediate
integrals, or integrals of the first order, which connect the
same varying marks of position and their first differential
coefficients with the time, the masses, and the initial marks of
position, are the result of elimination of the same auxiliary
constant $H$ between the equations (R.) and (E.).  Our
fundamental formula, and intermediate and final integrals, can
therefore be very simply expressed with any new sets of
coordinates; and the partial differential equations (F.) (G.),
which our characteristic function~$V$ must satisfy, and which
are, as we have said, essential in the theory of that function,
can also easily be expressed with any such transformed
coordinates, by merely combining the final and initial
expressions of the law of living force,
$$T = U + H,
   \eqno {\rm (6.)}$$
$$T_0 = U_0 + H,
   \eqno {\rm (7.)}$$
with the new groups (R.) and (S.).  For this purpose we must now
consider the function~$U$, of the masses and mutual
distances of the several points of the system, as depending on
the new marks of position
$\eta_1 \, \eta_2 \,\ldots\, \eta_{3n}$; and the analogous function
$U_0$, as depending similarly on the initial quantities
$e_1 \, e_2 \,\ldots\, e_{3n}$; we must also suppose that $T$ is
expressed (as it may) as a function of its own coefficients,
$\displaystyle
{\delta T \over \delta \eta_1'},
{\delta T \over \delta \eta_2'},\ldots\,
{\delta T \over \delta \eta_{3n}'}$,
which will always be, with respect to these, homogeneous of the
second dimension, and may also involve explicitly the quantities
$\eta_1 \, \eta_2 \,\ldots\, \eta_{3n}$;
and that $T_0$ is expressed as a similar function of its
coefficients
$\displaystyle
{\delta T_0 \over \delta e_1'},
{\delta T_0 \over \delta e_2'},\ldots\,
{\delta T_0 \over \delta e_{3n}'}$;
so that
$$\left. \eqalign{
T   &= F \left( {\delta T \over \delta \eta_1'},
                {\delta T \over \delta \eta_2'},\ldots\,
                {\delta T \over \delta \eta_{3n}'}
                \right),\cr
T_0 &= F \left( {\delta T_0 \over \delta e_1'},
                {\delta T_0 \over \delta e_2'},\ldots\,
                {\delta T_0 \over \delta e_{3n}'}
                \right);\cr}
   \right\}
   \eqno {\rm (25.)}$$
and that then these coefficients of $T$ and $T_0$ are changed to
their values (R.) and (S.), so as to give, instead of (F.) and
(G.), two other transformed equations, namely,
$$F \left( {\delta V \over \delta \eta_1},
           {\delta V \over \delta \eta_2},\ldots\,
           {\delta V \over \delta \eta_{3n}}
           \right)
   = U + H,
   \eqno {\rm (T.)}$$
and, on account of the homogeneity and dimension of $T_0$,
$$F \left( {\delta V \over \delta e_1},
           {\delta V \over \delta e_2},\ldots\,
           {\delta V \over \delta e_{3n}}
           \right)
   = U_0 + H.
   \eqno {\rm (U.)}$$

\bigbreak

8.
Nor is there any difficulty in deducing analogous transformations
for the known differential equations of motion of the second
order, of any system of free points, by taking the variation of
the new form (T.) of the law of living force, and by attending to
the dynamical meanings of the coefficients of our characteristic
function.  For if we observe that the final living force $2T$,
when considered as a function of
$\eta_1 \, \eta_2 \,\ldots\, \eta_{3n}$,
and of
$\eta_1' \, \eta_2' \,\ldots\, \eta_{3n}'$,
is necessarily homogeneous of the second dimension with respect
to the latter set of variables, and must therefore satisfy the
condition
$$2T =   \eta_1'    {\delta T \over \delta \eta_1'}
       + \eta_2'    {\delta T \over \delta \eta_2'}
       + \cdots
       + \eta_{3n}' {\delta T \over \delta \eta_{3n}'},
   \eqno {\rm (26.)}$$
we shall perceive that its total variation,
$$\left. \eqalign{
\delta T
  &=   {\delta T \over \delta \eta_1} \delta \eta_1
     + {\delta T \over \delta \eta_2} \delta \eta_2
     + \cdots
     + {\delta T \over \delta \eta_{3n}} \delta \eta_{3n} \cr
  &  + {\delta T \over \delta \eta_1'} \delta \eta_1'
     + {\delta T \over \delta \eta_2'} \delta \eta_2'
     + \cdots
     + {\delta T \over \delta \eta_{3n}'} \delta \eta_{3n}',\cr}
   \right\}
   \eqno {\rm (27.)}$$
may be put under the form
$$\left. \eqalign{
\delta T
   &=   \eta_1'    \, \delta {\delta T \over \delta \eta_1'}
      + \eta_2'    \, \delta {\delta T \over \delta \eta_2'}
      + \cdots
      + \eta_{3n}' \, \delta {\delta T \over \delta \eta_{3n}'} \cr
   &\mathrel{\phantom{=}} \mathord{}
      - {\delta T \over \delta \eta_1} \delta \eta_1
      - {\delta T \over \delta \eta_2} \delta \eta_2
      - \cdots
      - {\delta T \over \delta \eta_{3n}} \delta \eta_{3n} \cr
   &=   \sum \mathbin{.} \eta' \, \delta {\delta T \over \delta \eta'}
      - \sum \mathbin{.} {\delta T \over \delta \eta} \delta \eta \cr
   &=   \sum \mathbin{.} \left( \eta' \, \delta {\delta V \over \delta \eta}
      - {\delta T \over \delta \eta} \delta \eta \right),\cr}
   \right\}
   \eqno {\rm (28.)}$$
and therefore that the total variation of the new partial
differential equation (T.) may be thus written,
$$\sum \mathbin{.} \left( \eta' \, \delta {\delta V \over \delta \eta}
      - {\delta T \over \delta \eta} \delta \eta \right)
   = \sum \mathbin{.} {\delta U \over \delta \eta} \delta \eta
      + \delta H:
   \eqno {\rm (V.)}$$
in which, if we observe that
$\displaystyle \eta' = {d \eta \over dt}$,
and that the quantities of the form $\eta$ are the only ones
which vary with the time, we shall see that
$$\sum \mathbin{.} \eta' \delta{\delta V \over \delta \eta}
   = \sum \left(
      {d \over dt} {\delta V \over \delta \eta} \mathbin{.} \delta \eta
      + {d \over dt} {\delta V \over \delta e} \mathbin{.} \delta e
      \right)
      + {d \over dt} {\delta V \over \delta H} \mathbin{.} \delta H,
   \eqno {\rm (29.)}$$
because the identical equation $\delta dV = d \delta V$ gives,
when developed,
$$\sum \left( \delta {\delta V \over \delta \eta} \mathbin{.} d\eta
      + \delta {\delta V \over \delta e} \mathbin{.} de \right)
      + \delta {\delta V \over \delta H} \mathbin{.} dH
  = \sum \left( d {\delta V \over \delta \eta} \mathbin{.} \delta \eta
      + d {\delta V \over \delta e} \mathbin{.} \delta e \right)
      + d {\delta V \over \delta H} \mathbin{.} \delta H.
   \eqno {\rm (30.)}$$
Decomposing, therefore, the expression (V.), for the variation of
half the living force, into as many separate equations as it
contains independent variations, we obtain, not only the equation
$${d \over dt} {\delta V \over \delta H} = 1,
   \eqno {\rm (K.)}$$
which had already presented itself, and the group
$${d \over dt} {\delta V \over \delta e_1} = 0,\quad
  {d \over dt} {\delta V \over \delta e_2} = 0,\quad \cdots \quad
  {d \over dt} {\delta V \over \delta e_{3n}} = 0,
   \eqno {\rm (W.)}$$
which might have been at once obtained by differentiation from
the final integrals (S.), but also a group of $3n$ other
equations of the form
$${d \over dt} {\delta V \over \delta \eta}
      - {\delta T \over \delta \eta}
   = {\delta U \over \delta \eta},
   \eqno {\rm (X.)}$$
which give, by the intermediate integrals (R.),
$${d \over dt} {\delta T \over \delta \eta'}
      - {\delta T \over \delta \eta}
   = {\delta U \over \delta \eta}:
   \eqno {\rm (Y.)}$$
that is, more fully,
$$\left. \eqalign{
{d \over dt} {\delta T \over \delta \eta_1'}
      - {\delta T \over \delta \eta_1}
   &= {\delta U \over \delta \eta_1};\cr
{d \over dt} {\delta T \over \delta \eta_2'}
      - {\delta T \over \delta \eta_2}
   &= {\delta U \over \delta \eta_2};\cr
\cdots \cdots \cr
{d \over dt} {\delta T \over \delta \eta_{3n}'}
      - {\delta T \over \delta \eta_{3n}}
   &= {\delta U \over \delta \eta_{3n}}.\cr}
   \right\}
   \eqno {\rm (Z.)}$$

These last transformations of the differential equations of
motion of the second order, of an attracting or repelling system,
coincide in all respects (a slight difference of notation
excepted,) with the elegant canonical forms in the
{\it M\'{e}canique Analytique\/} of {\sc Lagrange}; but it seemed
worth while to deduce them here anew, from the properties of our
characteristic function.  And if we were to suppose (as it has
often been thought convenient and even necessary to do,) that the
$n$ points of a system are not entirely free, nor subject only to
their own mutual attractions or repulsions, but connected by any
geometrical conditions, and influenced by any foreign agencies,
consistent with the law of conservation of living force; so that
the number of independent marks of position would be now less
numerous, and the force-function $U$ less simple than before; it
might still be proved, by a reasoning very similar to the foregoing,
that on these suppositions also (which however, the dynamical
spirit is tending more and more to exclude,) the accumulated
living force or action~$V$ of the system is a {\it characteristic
motion-function\/} of the kind already explained; having the same
law and formula of variation, which are susceptible of the same
transformations; obliged to satisfy in the same way a final and
an initial relation between its partial differential coefficients
of the first order; conducting, by the variation of one of these
two relations, to the same canonical forms assigned by {\sc Lagrange}
for the differential equations of motion; and furnishing, on the
same principles as before, their intermediate and their final
integrals.  To those imaginable cases, indeed, in which the law of
living force no longer holds, our method also would not apply;
but it appears to be the growing conviction of the persons who
have meditated the most profoundly on the mathematical dynamics
of the universe, that these are cases suggested by insufficient
views of the mutual actions of body.

\bigbreak

9.
It results from the foregoing remarks, that in order to apply our
method of the characteristic function to any problem of dynamics
respecting any moving system, the known law of living force is to
be combined with our law of varying action; and that the general
expression of this latter law is to be obtained in the following
manner.  We are first to express the quantity $T$, namely, the
half of the living force of the system, as a function (which will
always be homogeneous of the second dimension,) of the
differential coefficients or rates of increase
$\eta_1'$, $\eta_2'$ \&c., of any rectangular coordinates, or
other marks of position of the system: we are next to take the
variation of this homogeneous function with respect to those
rates of increase,  and to change the variations of those rates
$\delta \eta_1'$, $\delta \eta_2'$, \&c.,
to the variations $\delta \eta_1$, $\delta \eta_2$, \&c.,
of the marks of position themselves; and then to subtract the
initial from the final value of the result, and to equate the
remainder to $\delta V - t \, \delta H$.  A slight consideration
will show that this general rule or process for obtaining the
variation of the characteristic function $V$, is applicable even
when the marks of position $\eta_1$, $\eta_2$, \&c. are not all
independent of each other; which will happen when they have been
made, from any motive of convenience, more numerous than the
rectangular coordinates of the several points of the system.  For
if we suppose that the $3n$ rectangular coordinates
$x_1 \, y_1 \, z_1 \,\ldots\, x_n \, y_n \, z_n$
have been expressed by any transformation as functions of
$3n + k$ other marks of position,
$\eta_1 \, \eta_2 \,\ldots\, \eta_{3n+k}$,
which must therefore be connected by $k$ equations of condition,
$$\left. \eqalign{
0 &= \phi_1( \eta_1, \eta_2,\ldots\, \eta_{3n+k} ),\cr
0 &= \phi_2( \eta_1, \eta_2,\ldots\, \eta_{3n+k} ),\cr
\noalign{\hbox{$\cdots\cdots$}}
0 &= \phi_k( \eta_1, \eta_2,\ldots\, \eta_{3n+k} ),\cr}
   \right\}
   \eqno {\rm (31.)}$$
giving $k$ of the new marks of position as functions of the
remaining $3n$,
$$\left. \eqalign{
\eta_{3n+1} &= \psi_1( \eta_1, \eta_2,\ldots\, \eta_{3n} ),\cr
\eta_{3n+2} &= \psi_2( \eta_1, \eta_2,\ldots\, \eta_{3n} ),\cr
\noalign{\hbox{$\cdots\cdots$}}
\eta_{3n+k} &= \psi_k( \eta_1, \eta_2,\ldots\, \eta_{3n} ),\cr}
   \right\}
   \eqno {\rm (32.)}$$
the expression
$$T = {\textstyle {1 \over 2}} \sum \mathbin{.} m
         (x'^2 + y'^2 + z'^2),
   \eqno {\rm (4.)}$$
will become, by the introduction of these new variables, a
homogeneous function of the second dimension of the $3n + k$
rates of increase
$\eta_1', \eta_2',\ldots\, \eta_{3n+k}'$,
involving also in general
$\eta_1, \eta_2,\ldots\, \eta_{3n+k}$,
and having a variation which may be thus expressed:
$$\left. \eqalign{
\delta T
   = & \left( {\delta T \over \delta \eta_1'} \right)
            \delta \eta_1'
     + \left( {\delta T \over \delta \eta_2'} \right)
            \delta \eta_2'
     + \cdots
     + \left( {\delta T \over \delta \eta_{3n+k}'} \right)
            \delta \eta_{3n+k}' \cr
   + & \left( {\delta T \over \delta \eta_1} \right)
            \delta \eta_1
     + \left( {\delta T \over \delta \eta_2} \right)
            \delta \eta_2
     + \cdots
     + \left( {\delta T \over \delta \eta_{3n+k}} \right)
            \delta \eta_{3n+k};\cr}
   \right\}
   \eqno {\rm (33.)}$$
or in this other way,
$$\left. \eqalign{
\delta T
   = & {\delta T \over \delta \eta_1'}
            \delta \eta_1'
     + {\delta T \over \delta \eta_2'}
            \delta \eta_2'
     + \cdots
     + {\delta T \over \delta \eta_{3n}'}
            \delta \eta_{3n}' \cr
   + & {\delta T \over \delta \eta_1}
            \delta \eta_1
     + {\delta T \over \delta \eta_2}
            \delta \eta_2
     + \cdots
     + {\delta T \over \delta \eta_{3n}}
            \delta \eta_{3n},\cr}
   \right\}
   \eqno {\rm (34.)}$$
on account of the relations (32.) which give, when differentiated
with respect to the time,
$$\left. \eqalign{
\eta_{3n+1}'
   &=   \eta_1'    {\delta \psi_1 \over \delta \eta_1}
      + \eta_2'    {\delta \psi_1 \over \delta \eta_2}
      + \cdots
      + \eta_{3n}' {\delta \psi_1 \over \delta \eta_{3n}},\cr
\eta_{3n+2}'
   &=   \eta_1'    {\delta \psi_2 \over \delta \eta_1}
      + \eta_2'    {\delta \psi_2 \over \delta \eta_2}
      + \cdots
      + \eta_{3n}' {\delta \psi_2 \over \delta \eta_{3n}},\cr
\noalign{\hbox{$\cdots\cdots$}}
\eta_{3n+k}'
   &=   \eta_1'    {\delta \psi_k \over \delta \eta_1}
      + \eta_2'    {\delta \psi_k \over \delta \eta_2}
      + \cdots
      + \eta_{3n}' {\delta \psi_k \over \delta \eta_{3n}},\cr}
   \right\}
   \eqno {\rm (35.)}$$
and therefore, attending only to the variations of quantities of
the form $\eta'$,
$$\left. \eqalign{
\delta \eta_{3n+1}'
   &=   {\delta \psi_1 \over \delta \eta_1} \delta \eta_1'
      + {\delta \psi_1 \over \delta \eta_2} \delta \eta_2'
      + \cdots
      + {\delta \psi_1 \over \delta \eta_{3n}} \delta \eta_{3n}',\cr
\delta \eta_{3n+2}'
   &=   {\delta \psi_2 \over \delta \eta_1} \delta \eta_1'
      + {\delta \psi_2 \over \delta \eta_2} \delta \eta_2'
      + \cdots
      + {\delta \psi_2 \over \delta \eta_{3n}} \delta \eta_{3n}',\cr
\noalign{\hbox{$\cdots\cdots$}}
\delta \eta_{3n+k}'
   &=   {\delta \psi_k \over \delta \eta_1} \delta \eta_1'
      + {\delta \psi_k \over \delta \eta_2} \delta \eta_2'
      + \cdots
      + {\delta \psi_k \over \delta \eta_{3n}} \delta \eta_{3n}'.\cr}
   \right\}
   \eqno {\rm (36.)}$$

Comparing the two expressions (33.) and (34.), we find by (36.)
the relations
$$\left. \eqalign{
{\delta T \over \delta \eta_1'}
   &=   \left( {\delta T \over \delta \eta_1'} \right)
      + \left( {\delta T \over \delta \eta_{3n+1}'} \right)
         {\delta \psi_1 \over \delta \eta_1}
      + \left( {\delta T \over \delta \eta_{3n+2}'} \right)
         {\delta \psi_2 \over \delta \eta_1}
      + \cdots
      + \left( {\delta T \over \delta \eta_{3n+k}'} \right)
         {\delta \psi_k \over \delta \eta_1},\cr
{\delta T \over \delta \eta_2'}
   &=   \left( {\delta T \over \delta \eta_2'} \right)
      + \left( {\delta T \over \delta \eta_{3n+1}'} \right)
         {\delta \psi_1 \over \delta \eta_2}
      + \left( {\delta T \over \delta \eta_{3n+2}'} \right)
         {\delta \psi_2 \over \delta \eta_2}
      + \cdots
      + \left( {\delta T \over \delta \eta_{3n+k}'} \right)
         {\delta \psi_k \over \delta \eta_2},\cr
\noalign{\hbox{$\cdots\cdots$}}
{\delta T \over \delta \eta_{3n}'}
   &=   \left( {\delta T \over \delta \eta_{3n}'} \right)
      + \left( {\delta T \over \delta \eta_{3n+1}'} \right)
         {\delta \psi_1 \over \delta \eta_{3n}}
      + \left( {\delta T \over \delta \eta_{3n+2}'} \right)
         {\delta \psi_2 \over \delta \eta_{3n}}
      + \cdots
      + \left( {\delta T \over \delta \eta_{3n+k}'} \right)
         {\delta \psi_k \over \delta \eta_{3n}};\cr}
   \right\}
   \eqno {\rm (37.)}$$
which give, by (32.),
$$    {\delta T \over \delta \eta_1'} \delta \eta_1
    + {\delta T \over \delta \eta_2'} \delta \eta_2
    + \cdots
    + {\delta T \over \delta \eta_{3n}'} \delta \eta_{3n}
  =   \left( {\delta T \over \delta \eta_1'} \right)
            \delta \eta_1
    + \left( {\delta T \over \delta \eta_2'} \right)
            \delta \eta_2
    + \cdots
    + \left( {\delta T \over \delta \eta_{3n+k}'} \right)
            \delta \eta_{3n+k};
   \eqno {\rm (38.)}$$
we may therefore put the expression (Q.) under the following
more general form,
$$\delta  V
   =  \sum \mathbin{.}
      \left( {\delta T \over \delta \eta'} \right) \delta \eta
    - \sum \mathbin{.}
      \left( {\delta T_0 \over \delta e'}  \right) \delta e
    + t \, \delta H,
   \eqno {\rm (A^1.)}$$
the coefficients
$\displaystyle \left( {\delta T \over \delta \eta'} \right)$
being formed by treating all the $3n + k$ quantities
$\eta_1', \eta_2',\ldots, \eta_{3n+k}'$,
as independent; which was the extension above announced, of the
rule for forming the variation of the characteristic function $V$.

We cannot, however, immediately decompose this new expression
(A${}^1$.) for $\delta V$, as we did the expression (Q.), by
treating all the variations $\delta \eta$, $\delta e$, as
independent; but we may decompose it so, if we previously combine
it with the final equations of condition (31.), and with the
analogous initial equations of condition, namely,
$$\left. \eqalign{
0 &= \Phi_1( e_1, e_2,\ldots\, e_{3n+k} ),\cr
0 &= \Phi_2( e_1, e_2,\ldots\, e_{3n+k} ),\cr
\noalign{\hbox{$\cdots\cdots$}}
0 &= \Phi_k( e_1, e_2,\ldots\, e_{3n+k} ),\cr}
   \right\}
   \eqno {\rm (39.)}$$
which we may do by adding the variations of the connecting
functions $\phi_1,\ldots, \phi_k$, $\Phi_1,\ldots\, \Phi_k$
multiplied respectively by the factors to be determined,
$\lambda_1,\ldots\, \lambda_k$, $\Lambda_1,\ldots\, \Lambda_k$.
In this manner the law of varying action takes this new form,
$$\delta  V
   =  \sum \mathbin{.}
      \left( {\delta T \over \delta \eta'} \right) \delta \eta
    - \sum \mathbin{.} \left( {\delta T_0 \over \delta e'}  \right) \delta e
  + t \, \delta H
  + \sum \mathbin{.} \lambda \, \delta \phi
  + \sum \mathbin{.} \Lambda \, \delta \Phi;
   \eqno {\rm (B^1.)}$$
and decomposes itself into $6n + 2k + 1$ separate expressions,
for the partial differential coefficients of the first order of
the characteristic function $V$, namely, into the following,
$$\left. \eqalign{
{\delta V \over \delta \eta_1}
   &= \left( {\delta T \over \delta \eta_1'} \right)
      + \lambda_1 {\delta \phi_1 \over \delta \eta_1}
      + \lambda_2 {\delta \phi_2 \over \delta \eta_1}
      + \cdots
      + \lambda_k {\delta \phi_k \over \delta \eta_1},\cr
{\delta V \over \delta \eta_2}
   &= \left( {\delta T \over \delta \eta_2'} \right)
      + \lambda_1 {\delta \phi_1 \over \delta \eta_2}
      + \lambda_2 {\delta \phi_2 \over \delta \eta_2}
      + \cdots
      + \lambda_k {\delta \phi_k \over \delta \eta_2},\cr
\noalign{\hbox{$\cdots\cdots$}}
{\delta V \over \delta \eta_{3n+k}}
   &= \left( {\delta T \over \delta \eta_{3n+k}'} \right)
      + \lambda_1 {\delta \phi_1 \over \delta \eta_{3n+k}}
      + \cdots
      + \lambda_k {\delta \phi_k \over \delta \eta_{3n+k}},\cr}
   \right\}
   \eqno {\rm (C^1.)}$$
and
$$\left. \eqalign{
{\delta V \over \delta e_1}
   &= - \left( {\delta T \over \delta e_1'} \right)
      + \Lambda_1 {\delta \Phi_1 \over \delta e_1}
      + \Lambda_2 {\delta \Phi_2 \over \delta e_1}
      + \cdots
      + \Lambda_k {\delta \Phi_k \over \delta e_1},\cr
{\delta V \over \delta e_2}
   &= - \left( {\delta T \over \delta e_2'} \right)
      + \Lambda_1 {\delta \Phi_1 \over \delta e_2}
      + \Lambda_2 {\delta \Phi_2 \over \delta e_2}
      + \cdots
      + \Lambda_k {\delta \Phi_k \over \delta e_2},\cr
\noalign{\hbox{$\cdots\cdots$}}
{\delta V \over \delta e_{3n+k}}
   &= - \left( {\delta T \over \delta e_{3n+k}'} \right)
      + \Lambda_1 {\delta \Phi_1 \over \delta e_{3n+k}}
      + \cdots
      + \Lambda_k {\delta \Phi_k \over \delta e_{3n+k}},\cr}
   \right\}
   \eqno {\rm (D^1.)}$$
besides the old equation (E.).  The analogous introduction of
multipliers in the canonical forms of {\sc Lagrange}, for the
differential equations of motion of the second order, by which a
sum such as
$\displaystyle \sum \mathbin{.} \lambda {\delta \phi \over \delta \eta}$
is added to
$\displaystyle {\delta U \over \delta \eta}$
in the second member of the formula (Y.), is also easily
justified on the principles of the present essay.

\bigbreak

{\sectiontitle
Separation of the relative motion of a system from the motion
of its centre of gravity; characteristic function for such motion,
and law of its variation.\par}

\nobreak\bigskip

10.
As an example of the foregoing transformations, and at the same
time as an important application, we shall now introduce relative
coordinates, $x_\prime$~$y_\prime$~$z_\prime$, referred to an
internal origin
$x_{\prime\prime}$~$y_{\prime\prime}$~$z_{\prime\prime}$;
that is, we shall put
$$x_i = x_{\prime i} + x_{\prime\prime},\quad
  y_i = y_{\prime i} + y_{\prime\prime},\quad
  z_i = z_{\prime i} + z_{\prime\prime},
   \eqno {\rm (40.)}$$
and in like manner
$$a_i = a_{\prime i} + a_{\prime\prime},\quad
  b_i = b_{\prime i} + b_{\prime\prime},\quad
  c_i = c_{\prime i} + c_{\prime\prime};
   \eqno {\rm (41.)}$$
together with the differentiated expressions
$$x_i' = x_{\prime i}' + x_{\prime\prime}',\quad
  y_i' = y_{\prime i}' + y_{\prime\prime}',\quad
  z_i' = z_{\prime i}' + z_{\prime\prime}',
   \eqno {\rm (42.)}$$
and
$$a_i' = a_{\prime i}' + a_{\prime\prime}',\quad
  b_i' = b_{\prime i}' + b_{\prime\prime}',\quad
  c_i' = c_{\prime i}' + c_{\prime\prime}'.
   \eqno {\rm (43.)}$$
Introducing the expressions (42.) for the rectangular components
of velocity, we find that the value given by (4.) for the living
force $2T$ decomposes itself into the three following parts,
$$\left. \eqalign{
2T
   &=   \sum \mathbin{.} m (x'^2 + y'^2 + z'^2) \cr
   &=   \sum \mathbin{.} m (x_\prime'^2 + y_\prime'^2 + z_\prime'^2)
      + 2 (x_{\prime\prime}' \sum \mathbin{.} m x_\prime'
         + y_{\prime\prime}' \sum \mathbin{.} m y_\prime'
         + z_{\prime\prime}' \sum \mathbin{.} m z_\prime') \cr
   &\mathrel{\phantom{=}} \mathord{}
      + (x_{\prime\prime}'^2 + y_{\prime\prime}'^2
                             + z_{\prime\prime}'^2) \sum m;\cr}
   \right\}
   \eqno {\rm (44.)}$$
if then we establish, as we may, the three equations of condition,
$$\sum \mathbin{.} m x_\prime = 0,\quad
  \sum \mathbin{.} m y_\prime = 0,\quad
  \sum \mathbin{.} m z_\prime = 0,
   \eqno {\rm (45.)}$$
which give by (40.),
$$x_{\prime\prime} = {\sum \mathbin{.} mx \over \sum m},\quad
  y_{\prime\prime} = {\sum \mathbin{.} my \over \sum m},\quad
  z_{\prime\prime} = {\sum \mathbin{.} mz \over \sum m},
   \eqno {\rm (46.)}$$
so that
$x_{\prime\prime}$~$y_{\prime\prime}$~$z_{\prime\prime}$
are now the coordinates of the point which is
called the centre of gravity of the system, we may reduce the
function~$T$ to the form
$$T = T_\prime + T_{\prime\prime},
   \eqno {\rm (47.)}$$
in which
$$T_\prime = {\textstyle {1 \over 2}} \sum \mathbin{.} m
      (x_\prime'^2 + y_\prime'^2 + z_\prime'^2),
   \eqno {\rm (48.)}$$
and
$$T_{\prime\prime} = {\textstyle {1 \over 2}}
      (x_{\prime\prime}'^2 + y_{\prime\prime}'^2
                           + z_{\prime\prime}'^2) \sum m.
   \eqno {\rm (49.)}$$

By this known decomposition, the whole living force $2T$ of the
system is resolved into the two parts $2 T_\prime$ and
$2 T_{\prime\prime}$, of which the former, $2T_\prime$, may be
called the {\it relative living force}, being that which results
solely from the relative velocities of the points of the system, in
their motions about their common centre of gravity
$x_{\prime\prime}$~$y_{\prime\prime}$~$z_{\prime\prime}$;
while the latter part, $2T_{\prime\prime}$, results only from the
absolute motion of that centre of gravity in space, and is the same
as if all the masses of the system were united in that common
centre.  At the same time, the law of living force, $T = U + H$,
(6.), resolves itself by the law of motion of the centre of gravity
into the two following separate equations,
$$T_\prime = U + H_\prime,
   \eqno {\rm (50.)}$$
and
$$T_{\prime\prime} = H_{\prime\prime};
   \eqno {\rm (51.)}$$
$H_\prime$ and $H_{\prime\prime}$ being two new constants
independent of the time~$t$, and such that their sum
$$H_\prime + H_{\prime\prime} = H.
   \eqno {\rm (52.)}$$

And we may in like manner decompose the action, or accumulated
living force $V$, which is equal to the definite integral
$\displaystyle \int_0^t 2T \,dt$, into the two following analogous
parts,
$$V = V_\prime + V_{\prime\prime},
   \eqno {\rm (E^1.)}$$
determined by the two equations,
$$V_\prime = \int_0^t 2T_\prime \,dt,
   \eqno {\rm (F^1.)}$$
and
$$V_{\prime\prime} = \int_0^t 2T_{\prime\prime} \,dt.
   \eqno {\rm (G^1.)}$$
The last equation gives by (51.),
$$V_{\prime\prime} = 2 H_{\prime\prime} t;
   \eqno {\rm (53.)}$$
a result which, by the law of motion of the centre of gravity, may
be thus expressed,
$$V_{\prime\prime} = \sqrt{
        (x_{\prime\prime} - a_{\prime\prime})^2
      + (y_{\prime\prime} - b_{\prime\prime})^2
      + (z_{\prime\prime} - c_{\prime\prime})^2}
      \mathbin{.} \sqrt{2H_{\prime\prime} \sum m}:
   \eqno {\rm (H^1.)}$$
$a_{\prime\prime}$~$b_{\prime\prime}$~$c_{\prime\prime}$ being
the initial coordinates of the centre of gravity, so that
$$a_{\prime\prime} = {\sum \mathbin{.} ma \over \sum m},\quad
  b_{\prime\prime} = {\sum \mathbin{.} mb \over \sum m},\quad
  c_{\prime\prime} = {\sum \mathbin{.} mc \over \sum m}.
   \eqno {\rm (54.)}$$
And for the variation $\delta V$ of the whole function~$V$, the
rule of the last number gives
$$\left. \eqalign{
\delta V
   &= \sum \mathbin{.} m
        (   x_\prime' \, \delta x_\prime - a_\prime' \, \delta a_\prime
          + y_\prime' \, \delta y_\prime - b_\prime' \, \delta b_\prime
          + z_\prime' \, \delta z_\prime - c_\prime' \, \delta c_\prime
        ) \cr
   &\mathrel{\phantom{=}} \mathord{} \mathord{}
      + (   x_{\prime\prime}' \, \delta x_{\prime\prime}
          - a_{\prime\prime}' \, \delta a_{\prime\prime}
          + y_{\prime\prime}' \, \delta y_{\prime\prime}
          - b_{\prime\prime}' \, \delta b_{\prime\prime}
          + z_{\prime\prime}' \, \delta z_{\prime\prime}
          - c_{\prime\prime}' \, \delta c_{\prime\prime}
        ) \sum m \cr
   &\mathrel{\phantom{=}} \mathord{} \mathord{}
      + t \, \delta H
      + \lambda_1 \sum \mathbin{.} m \, \delta x_\prime
      + \lambda_2 \sum \mathbin{.} m \, \delta y_\prime
      + \lambda_3 \sum \mathbin{.} m \, \delta z_\prime \cr
   &\mathrel{\phantom{=}} \mathord{} \mathord{}
      + \Lambda_1 \sum \mathbin{.} m \, \delta a_\prime
      + \Lambda_2 \sum \mathbin{.} m \, \delta b_\prime
      + \Lambda_3 \sum \mathbin{.} m \, \delta c_\prime;\cr}
   \right\}
   \eqno {\rm (I^1.)}$$
while the variation of the part $V_{\prime\prime}$, determined by
the equation (H${}^1$.), is easily shown to be equivalent to the
part
$$\delta V_{\prime\prime}
   = (   x_{\prime\prime}' \, \delta x_{\prime\prime}
       - a_{\prime\prime}' \, \delta a_{\prime\prime}
       + y_{\prime\prime}' \, \delta y_{\prime\prime}
       - b_{\prime\prime}' \, \delta b_{\prime\prime}
       + z_{\prime\prime}' \, \delta z_{\prime\prime}
       - c_{\prime\prime}' \, \delta c_{\prime\prime}) \sum m
       + t \, \delta H_{\prime\prime};
   \eqno {\rm (K^1.)}$$
the variation of the other part $V_\prime$ may therefore be thus
expressed,
$$\left. \eqalign{
\delta V_\prime
   &= \sum \mathbin{.} m
        (   x_\prime' \, \delta x_\prime - a_\prime' \, \delta a_\prime
          + y_\prime' \, \delta y_\prime - b_\prime' \, \delta b_\prime
          + z_\prime' \, \delta z_\prime - c_\prime' \, \delta c_\prime
        ) \cr
   &\mathrel{\phantom{=}} \mathord{} \mathord{}
      + t \, \delta H_\prime
      + \lambda_1 \sum \mathbin{.} m \, \delta x_\prime
      + \lambda_2 \sum \mathbin{.} m \, \delta y_\prime
      + \lambda_3 \sum \mathbin{.} m \, \delta z_\prime \cr
   &\mathrel{\phantom{=}} \mathord{} \mathord{}
      + \Lambda_1 \sum \mathbin{.} m \, \delta a_\prime
      + \Lambda_2 \sum \mathbin{.} m \, \delta b_\prime
      + \Lambda_3 \sum \mathbin{.} m \, \delta c_\prime:\cr}
   \right\}
   \eqno {\rm (L^1.)}$$
and it resolves itself into the following separate expressions,
in which the part $V_\prime$ is considered as a function of the
$6n + 1$ quantities
$x_{\prime i}$~$y_{\prime i}$~$z_{\prime i}$
$a_{\prime i}$~$b_{\prime i}$~$c_{\prime i}$
$H_\prime$, of which, however, only $6n - 5$ are really
independent:\hfil\break
first group,
$$\left. \multieqalign{
{\delta V_\prime \over \delta x_{\prime 1}}
   &= m_1 x_{\prime 1}' + \lambda_1 m_1; \quad \cdots &
{\delta V_\prime \over \delta x_{\prime n}}
   &= m_n x_{\prime n}' + \lambda_1 m_n; \cr
{\delta V_\prime \over \delta y_{\prime 1}}
   &= m_1 y_{\prime 1}' + \lambda_2 m_1; \quad \cdots &
{\delta V_\prime \over \delta y_{\prime n}}
   &= m_n y_{\prime n}' + \lambda_2 m_n; \cr
{\delta V_\prime \over \delta z_{\prime 1}}
   &= m_1 z_{\prime 1}' + \lambda_3 m_1; \quad \cdots &
{\delta V_\prime \over \delta z_{\prime n}}
   &= m_n z_{\prime n}' + \lambda_3 m_n; \cr}
   \right\}
   \eqno {\rm (M^1.)}$$
second group,
$$\left. \multieqalign{
{\delta V_\prime \over \delta a_{\prime 1}}
   &= - m_1 a_{\prime 1}' + \Lambda_1 m_1; \quad \cdots &
{\delta V_\prime \over \delta a_{\prime n}}
   &= - m_n a_{\prime n}' + \Lambda_1 m_n; \cr
{\delta V_\prime \over \delta b_{\prime 1}}
   &= - m_1 b_{\prime 1}' + \Lambda_2 m_1; \quad \cdots &
{\delta V_\prime \over \delta b_{\prime n}}
   &= - m_n b_{\prime n}' + \Lambda_2 m_n; \cr
{\delta V_\prime \over \delta c_{\prime 1}}
   &= - m_1 c_{\prime 1}' + \Lambda_3 m_1; \quad \cdots &
{\delta V_\prime \over \delta c_{\prime n}}
   &= - m_n c_{\prime n}' + \Lambda_3 m_n; \cr}
   \right\}
   \eqno {\rm (N^1.)}$$
and finally,
$${\delta V_\prime \over \delta H_\prime} = t.
   \eqno {\rm (O^1.)}$$

With respect to the six multipliers
$\lambda_1$~$\lambda_2$~$\lambda_3$
$\Lambda_1$~$\Lambda_2$~$\Lambda_3$
which were introduced by the 3 final equations of condition
(45.), and by the 3 analogous initial equations of condition,
$$\sum \mathbin{.} m a_\prime = 0,\quad
  \sum \mathbin{.} m b_\prime = 0,\quad
  \sum \mathbin{.} m c_\prime = 0;
   \eqno {\rm (55.)}$$
we have, by differentiating these conditions,
$$\sum \mathbin{.} m x_\prime' = 0,\quad
  \sum \mathbin{.} m y_\prime' = 0,\quad
  \sum \mathbin{.} m z_\prime' = 0,
   \eqno {\rm (56.)}$$
and
$$\sum \mathbin{.} m a_\prime' = 0,\quad
  \sum \mathbin{.} m b_\prime' = 0,\quad
  \sum \mathbin{.} m c_\prime' = 0;
   \eqno {\rm (57.)}$$
and therefore
$$\lambda_1
   = {\displaystyle \sum {\delta V_\prime \over \delta x_\prime}
      \over \sum m},\quad
  \lambda_2
   = {\displaystyle \sum {\delta V_\prime \over \delta y_\prime}
      \over \sum m},\quad
  \lambda_3
   = {\displaystyle \sum {\delta V_\prime \over \delta z_\prime}
      \over \sum m},
   \eqno {\rm (58.)}$$
and
$$\Lambda_1
   = {\displaystyle \sum {\delta V_\prime \over \delta a_\prime}
      \over \sum m},\quad
  \Lambda_2
   = {\displaystyle \sum {\delta V_\prime \over \delta b_\prime}
      \over \sum m},\quad
  \Lambda_3
   = {\displaystyle \sum {\delta V_\prime \over \delta c_\prime}
      \over \sum m}.
   \eqno {\rm (59.)}$$

\bigbreak

11.
As an example of the determination of these multipliers, we may
suppose that the part $V_\prime$, of the whole action $V$, has
been expressed, before differentiation, as a function of
$H_\prime$, and of these other $6n - 6$ independent quantities
$$\left. \multieqalign{
x_{\prime 1}   - x_{\prime n} &= \xi_1, &
x_{\prime 2}   - x_{\prime n} &= \xi_2, \quad \ldots &
x_{\prime n-1} - x_{\prime n} &= \xi_{n-1}, \cr
y_{\prime 1}   - y_{\prime n} &= \eta_1, &
y_{\prime 2}   - y_{\prime n} &= \eta_2, \quad \ldots &
y_{\prime n-1} - y_{\prime n} &= \eta_{n-1}, \cr
z_{\prime 1}   - z_{\prime n} &= \zeta_1, &
z_{\prime 2}   - z_{\prime n} &= \zeta_2, \quad \ldots &
z_{\prime n-1} - z_{\prime n} &= \zeta_{n-1}, \cr}
   \right\}
   \eqno {\rm (60.)}$$
and
$$\left. \multieqalign{
a_{\prime 1}   - a_{\prime n} &= \alpha_1, &
a_{\prime 2}   - a_{\prime n} &= \alpha_2, \quad \ldots &
a_{\prime n-1} - a_{\prime n} &= \alpha_{n-1}, \cr
b_{\prime 1}   - b_{\prime n} &= \beta_1, &
b_{\prime 2}   - b_{\prime n} &= \beta_2, \quad \ldots &
b_{\prime n-1} - b_{\prime n} &= \beta_{n-1}, \cr
c_{\prime 1}   - c_{\prime n} &= \gamma_1, &
c_{\prime 2}   - c_{\prime n} &= \gamma_2, \quad \ldots &
c_{\prime n-1} - c_{\prime n} &= \gamma_{n-1}; \cr}
   \right\}
   \eqno {\rm (61.)}$$
that is, of the {\it differences\/} only of the {\it
centrobaric\/} coordinates; or, in other words, as a function of
the coordinates (initial and final) of $n - 1$ points of the
system, referred to the $n^{\rm th}$ point, as an internal or
moveable origin: because the centrobaric coordinates
$x_{\prime i}$, $y_{\prime i}$, $z_{\prime i}$,
$a_{\prime i}$, $b_{\prime i}$, $c_{\prime i}$,
may themselves, by the equations of condition, be expressed as a
function of these, namely,
$$x_{\prime i} = \xi_i    - {\sum \mathbin{.} m \xi \over \sum m},\quad
  y_{\prime i} = \eta_i   - {\sum \mathbin{.} m \eta \over \sum m},\quad
  z_{\prime i} = \zeta_i  - {\sum \mathbin{.} m \zeta \over \sum m},
   \eqno {\rm (62.)}$$
and in like manner,
$$a_{\prime i} = \alpha_i - {\sum \mathbin{.} m \alpha \over \sum m},\quad
  b_{\prime i} = \beta_i  - {\sum \mathbin{.} m \beta \over \sum m},\quad
  c_{\prime i} = \gamma_i - {\sum \mathbin{.} m \gamma \over \sum m};
   \eqno {\rm (63.)}$$
in which we are to observe, that the six quantities
$\xi_n$~$\eta_n$~$\zeta_n$ $\alpha_n$~$\beta_n$~$\gamma_n$
must be considered as separately vanishing.  When $V_\prime$ has
been thus expressed as a function of the centrobaric coordinates,
involving their differences only, it will evidently satisfy the
six partial differential equations,
$$\left. \multieqalign{
\sum {\delta V_\prime \over \delta x_\prime} &= 0, &
\sum {\delta V_\prime \over \delta y_\prime} &= 0, &
\sum {\delta V_\prime \over \delta z_\prime} &= 0, \cr
\sum {\delta V_\prime \over \delta a_\prime} &= 0, &
\sum {\delta V_\prime \over \delta b_\prime} &= 0, &
\sum {\delta V_\prime \over \delta c_\prime} &= 0; \cr}
   \right\}
   \eqno {\rm (P^1.)}$$
after this preparation, therefore, of the function $V_\prime$,
the six multipliers determined by (58.) and (59.) will vanish, so
that we shall have
$$\lambda_1 = 0,\quad
  \lambda_2 = 0,\quad
  \lambda_3 = 0,\quad
  \Lambda_1 = 0,\quad
  \Lambda_2 = 0,\quad
  \Lambda_3 = 0,
   \eqno {\rm (64.)}$$
and the groups (M${}^1$.) and (N${}^1$.) will reduce themselves
to the two following:
$$\left. \multieqalign{
{\delta V_\prime \over \delta x_{\prime 1}}
   &= m_1 x_{\prime 1}'; &
{\delta V_\prime \over \delta x_{\prime 2}}
   &= m_2 x_{\prime 2}'; \quad \cdots &
{\delta V_\prime \over \delta x_{\prime n}}
   &= m_n x_{\prime n}'; \cr
{\delta V_\prime \over \delta y_{\prime 1}}
   &= m_1 y_{\prime 1}'; &
{\delta V_\prime \over \delta y_{\prime 2}}
   &= m_2 y_{\prime 2}'; \quad \cdots &
{\delta V_\prime \over \delta y_{\prime n}}
   &= m_n y_{\prime n}'; \cr
{\delta V_\prime \over \delta z_{\prime 1}}
   &= m_1 z_{\prime 1}'; \quad \cdots &
{\delta V_\prime \over \delta z_{\prime 2}}
   &= m_2 z_{\prime 2}'; \quad \cdots &
{\delta V_\prime \over \delta z_{\prime n}}
   &= m_n z_{\prime n}'; \cr}
   \right\}
   \eqno {\rm (Q^1.)}$$
and
$$\left. \multieqalign{
{\delta V_\prime \over \delta a_{\prime 1}}
   &= - m_1 a_{\prime 1}'; &
{\delta V_\prime \over \delta a_{\prime 2}}
   &= - m_2 a_{\prime 2}'; \quad \cdots &
{\delta V_\prime \over \delta a_{\prime n}}
   &= - m_n a_{\prime n}'; \cr
{\delta V_\prime \over \delta b_{\prime 1}}
   &= - m_1 b_{\prime 1}'; &
{\delta V_\prime \over \delta b_{\prime 2}}
   &= - m_2 b_{\prime 2}'; \quad \cdots &
{\delta V_\prime \over \delta b_{\prime n}}
   &= - m_n b_{\prime n}'; \cr
{\delta V_\prime \over \delta c_{\prime 1}}
   &= - m_1 c_{\prime 1}'; \quad \cdots &
{\delta V_\prime \over \delta c_{\prime 2}}
   &= - m_2 c_{\prime 2}'; \quad \cdots &
{\delta V_\prime \over \delta c_{\prime n}}
   &= - m_n c_{\prime n}'; \cr}
   \right\}
   \eqno {\rm (R^1.)}$$
analogous in all respects to the groups (C.) and (D.).  We find,
therefore, for the relative motion of a system about its own
centre of gravity, equations of the same form as those which we
had obtained before for the absolute motion of the same system of
points in space,  And we see that in investigating such relative
motion only, it is useful to confine ourselves to the part
$V_\prime$ of our whole characteristic function, that is, to the
{\it relative action\/} of the system, or accumulated living
force of the motion about the centre of gravity; and to consider
this part as the {\it characteristic function\/} of such relative
motion, in a sense analogous to that which has been already
explained.

This relative action, or part $V_\prime$, may, however, be
otherwise expressed, and even in an infinite variety of ways, on
account of the six equations of condition which connect the $6n$
centrobaric coordinates; and every different preparation of its
form will give a different set of values for the six multipliers
$\lambda_1$~$\lambda_2$~$\lambda_3$
$\Lambda_1$~$\Lambda_2$~$\Lambda_3$.
For example, we might eliminate, by a previous preparation, the
six centrobaric coordinates of the point $m_n$ from the
expression of $V_\prime$, so as to make this expression involve
only the centrobaric coordinates of the other $n - 1$ points of
the system, and then we should have
$${\delta V_\prime \over \delta x_{\prime n}} = 0,\quad
  {\delta V_\prime \over \delta y_{\prime n}} = 0,\quad
  {\delta V_\prime \over \delta z_{\prime n}} = 0,\quad
  {\delta V_\prime \over \delta a_{\prime n}} = 0,\quad
  {\delta V_\prime \over \delta b_{\prime n}} = 0,\quad
  {\delta V_\prime \over \delta c_{\prime n}} = 0,
   \eqno {\rm (S^1.)}$$
and therefore, by the six last equations of the groups (M${}^1$.)
and (N${}^1$.), the multipliers would take the values
$$\lambda_1 = - x_{\prime n}',\quad
  \lambda_2 = - y_{\prime n}',\quad
  \lambda_3 = - z_{\prime n}',\quad
  \Lambda_1 =   a_{\prime n}',\quad
  \Lambda_2 =   b_{\prime n}',\quad
  \Lambda_3 =   c_{\prime n}',\quad
   \eqno {\rm (65.)}$$
and would reduce, by (60.) and (61.), the preceding $6n - 6$
equations of the same groups (M${}^1$.) and (N${}^1$.), to the
forms
$$\left. \multieqalign{
{\delta V_\prime \over \delta x_{\prime 1}}
   &= m_1 \xi_1', &
{\delta V_\prime \over \delta x_{\prime 2}}
   &= m_2 \xi_2', \quad \cdots &
{\delta V_\prime \over \delta x_{\prime n-1}}
   &= m_{n-1} \xi_{n-1}', \cr
{\delta V_\prime \over \delta y_{\prime 1}}
   &= m_1 \eta_1', &
{\delta V_\prime \over \delta y_{\prime 2}}
   &= m_2 \eta_2', \quad \cdots &
{\delta V_\prime \over \delta y_{\prime n-1}}
   &= m_{n-1} \eta_{n-1}', \cr
{\delta V_\prime \over \delta z_{\prime 1}}
   &= m_1 \zeta_1', \quad \cdots &
{\delta V_\prime \over \delta z_{\prime 2}}
   &= m_2 \zeta_2', \quad \cdots &
{\delta V_\prime \over \delta z_{\prime n-1}}
   &= m_{n-1} \zeta_{n-1}', \cr}
   \right\}
   \eqno {\rm (T^1.)}$$
and
$$\left. \multieqalign{
{\delta V_\prime \over \delta a_{\prime 1}}
   &= - m_1 \alpha_1', &
{\delta V_\prime \over \delta a_{\prime 2}}
   &= - m_2 \alpha_2', \quad \cdots &
{\delta V_\prime \over \delta a_{\prime n-1}}
   &= - m_{n-1} \alpha_{n-1}', \cr
{\delta V_\prime \over \delta b_{\prime 1}}
   &= - m_1 \beta_1', &
{\delta V_\prime \over \delta b_{\prime 2}}
   &= - m_2 \beta_2', \quad \cdots &
{\delta V_\prime \over \delta b_{\prime n-1}}
   &= - m_{n-1} \beta_{n-1}', \cr
{\delta V_\prime \over \delta c_{\prime 1}}
   &= - m_1 \gamma_1', \quad \cdots &
{\delta V_\prime \over \delta c_{\prime 2}}
   &= - m_2 \gamma_2', \quad \cdots &
{\delta V_\prime \over \delta c_{\prime n-1}}
   &= - m_{n-1} \gamma_{n-1}'. \cr}
   \right\}
   \eqno {\rm (U^1.)}$$

\bigbreak

12.
We might also express the relative action $V_\prime$, not as a
function of the centrobaric, but of some other internal
coordinates, or marks of relative position.  We might, for
instance, express it and its variation as functions of the
$6n - 6$ independent internal coordinates
$\xi$~$\eta$~$\zeta$ $\alpha$~$\beta$~$\gamma$
already mentioned, and of their variations, defining these
without any reference to the centre of gravity, by the equations
$$\left. \multieqalign{
\xi_i    &= x_i - x_n, &
\eta_i   &= y_i - y_n, &
\zeta_i  &= z_i - z_n, \cr
\alpha_i &= a_i - a_n, &
\beta_i  &= b_i - b_n, &
\gamma_i &= c_i - c_n. \cr}
   \right\}
   \eqno {\rm (66.)}$$
For all such transformations of $\delta V_\prime$ it is easy to
establish a rule or law, which may be called the {\it law of
varying relative action\/} (exactly analogous to the rule
(B${}^1$.)), namely, the following:
$$\delta  V_\prime
   =  \sum \mathbin{.}
      \left( {\delta T_\prime \over \delta \eta_\prime'} \right)
            \delta \eta_\prime
    - \sum \mathbin{.}
      \left( {\delta T_{\prime 0} \over \delta e_\prime'}  \right)
            \delta e_\prime
    + t \, \delta H_\prime
    + \sum \mathbin{.} \lambda_\prime \, \delta \phi_\prime
    + \sum \mathbin{.} \Lambda_\prime \, \delta \Phi_\prime;
   \eqno {\rm (V^1.)}$$
which implies that we are to express the half $T_\prime$ of the
relative living force of the system as a function of the rates of
increase $\eta_\prime'$ of any marks of relative position; and
after taking its variation with respect to these rates, to change
their variations to the variations of the marks of position
themselves; then to subtract the initial from the final value of
the result, and to add the variations of the final and initial
functions $\phi_\prime$~$\Phi_\prime$, which enter into the
equations of condition, if any, of the form $\phi_\prime = 0$,
$\Phi_\prime = 0$, (connecting the final and initial marks of
relative position,) multiplied respectively by undetermined
factors $\lambda_\prime$~$\Lambda_\prime$; and lastly, to equate
the whole result to $\delta V_\prime - t \, \delta H_\prime$,
$H_\prime$ being the quantity independent of the time in the
equation (50.) of relative living force, and $V_\prime$ being the
relative action, of which we desired to express the variation.
It is not necessary to dwell here on the demonstration of this
new rule (V${}^1$.),  which may easily be deduced from the
principles already laid down; or by the calculus of variations
from the law of relative living force, combined with the
differential equations of second order of relative motion.

But to give an example of its application, let us resume the
problem already mentioned, namely to express $\delta V_\prime$ by
means of the $6n - 5$ independent variations
$\delta \xi_i$~$\delta \eta_i$~$\delta \zeta_i$
$\delta \alpha_i$~$\delta \beta_i$~$\delta \gamma_i$
$\delta H_\prime$.
For this purpose we shall employ a known transformation of the
relative living force $2T_\prime$, multiplied by the sum of the
masses of the system, namely the following:
$$2 T_\prime \sum m
   = \sum \mathbin{.} m_i m_k \{
        (x_i' - x_k')^2 + (y_i' - y_k')^2 + (z_i' - z_k')^2 \}:
   \eqno {\rm (67.)}$$
the sign of summation $\sum$ extending, in the second member, to
all the combinations of points two by two, which can be formed
without repetition.  This transformation gives, by (66.),
$$\left. \eqalign{
2 T_\prime \sum m
  &= m_n \sum_\prime \mathbin{.} m (\xi'^2 + \eta'^2 + \zeta'^2) \cr
  &\mathrel{\phantom{=}} \mathord{}
      + \sum_\prime \mathbin{.} m_i m_k \{
        (\xi_i' - \xi_k')^2 + (\eta_i' - \eta_k')^2
           + (\zeta_i' - \zeta_k')^2 \};\cr}
   \right\}
   \eqno {\rm (68.)}$$
the sign of summation $\sum_\prime$ extending only to the first
$n - 1$ points of the system.  Applying, therefore, our general
rule or law of varying relative action, and observing that the
$6n - 6$ internal coordinates
$\xi$~$\eta$~$\zeta$ $\alpha$~$\beta$~$\gamma$
are independent, we find the following new expression:
$$\left. \eqalign{
\delta V_\prime
   &= t \, \delta H_\prime
      + {m_n \over \sum m} \mathbin{.} \sum_\prime \mathbin{.} m
         ( \xi'   \, \delta \xi   - \alpha' \, \delta \alpha
         + \eta'  \, \delta \eta  - \beta'  \, \delta \beta
         + \zeta' \, \delta \zeta - \gamma' \, \delta \gamma ) \cr
   &\mathrel{\phantom{=}} \mathord{}
      + {1 \over \sum m} \mathbin{.} \sum_\prime \mathbin{.} m_i m_k \{
           (\xi_i' - \xi_k')(\delta \xi_i - \delta \xi_k)
         + (\eta_i' - \eta_k')(\delta \eta_i - \delta \eta_k) \cr
   &\hskip 9em
         + (\zeta_i' - \zeta_k')(\delta \zeta_i - \delta \zeta_k)
         \} \cr
   &\mathrel{\phantom{=}} \mathord{}
      - {1 \over \sum m} \mathbin{.} \sum_\prime \mathbin{.} m_i m_k \{
           (\alpha_i' - \alpha_k')(\delta \alpha_i - \delta \alpha_k)
         + (\beta_i' - \beta_k')(\delta \beta_i - \delta \beta_k) \cr
   &\hskip 9em
         + (\gamma_i' - \gamma_k')(\delta \gamma_i - \delta \gamma_k)
         \}:\cr}
   \right\}
   \eqno {\rm (W^1.)}$$
which gives, besides the equation (O${}^1$.), the following groups:
$$\left. \eqalign{
{\delta V_\prime \over \delta \xi_i}
   &= {m_i \over \sum m} \mathbin{.} \sum \mathbin{.} m (\xi_i' - \xi')
      = m_i \left( \xi_i'
         - {\sum_\prime m \xi' \over \sum m} \right),\cr
{\delta V_\prime \over \delta \eta_i}
   &= {m_i \over \sum m} \mathbin{.} \sum \mathbin{.} m (\eta_i' - \eta')
      = m_i \left( \eta_i'
         - {\sum_\prime m \eta' \over \sum m} \right),\cr
{\delta V_\prime \over \delta \zeta_i}
   &= {m_i \over \sum m} \mathbin{.} \sum \mathbin{.} m (\zeta_i' - \zeta')
      = m_i \left( \zeta_i'
         - {\sum_\prime m \zeta' \over \sum m} \right),\cr}
   \right\}
   \eqno {\rm (X^1.)}$$
and
$$\left. \eqalign{
{\delta V_\prime \over \delta \alpha_i}
   &= {- m_i \over \sum m} \mathbin{.}
         \sum \mathbin{.} m (\alpha_i' - \alpha')
      = - m_i \left( \alpha_i'
         - {\sum_\prime m \alpha' \over \sum m} \right),\cr
{\delta V_\prime \over \delta \beta_i}
   &= {- m_i \over \sum m} \mathbin{.}
         \sum \mathbin{.} m (\beta_i' - \beta')
      = - m_i \left( \beta_i'
         - {\sum_\prime m \beta' \over \sum m} \right),\cr
{\delta V_\prime \over \delta \gamma_i}
   &= {- m_i \over \sum m} \mathbin{.}
         \sum \mathbin{.} m (\gamma_i' - \gamma')
      = - m_i \left( \gamma_i'
         - {\sum_\prime m \gamma' \over \sum m} \right);\cr}
   \right\}
   \eqno {\rm (Y^1.)}$$
results which may be thus summed up:
$$\left. \eqalign{
\delta V_\prime
   &= t \, \delta H_\prime + \sum_\prime \mathbin{.} m
         ( \xi'   \, \delta \xi   - \alpha' \, \delta \alpha
         + \eta'  \, \delta \eta  - \beta'  \, \delta \beta
         + \zeta' \, \delta \zeta - \gamma' \, \delta \gamma ) \cr
   &\mathrel{\phantom{=}} \mathord{}
      - {1 \over \sum m}
         ( \sum_\prime m \xi'   \mathbin{.}
               \sum_\prime m \, \delta \xi
         + \sum_\prime m \eta'  \mathbin{.}
               \sum_\prime m \, \delta \eta
         + \sum_\prime m \zeta' \mathbin{.}
               \sum_\prime m \, \delta \zeta ) \cr
   &\mathrel{\phantom{=}} \mathord{}
      + {1 \over \sum m}
         ( \sum_\prime m \alpha' \mathbin{.}
               \sum_\prime m \, \delta \alpha
         + \sum_\prime m \beta'  \mathbin{.}
               \sum_\prime m \, \delta \beta
         + \sum_\prime m \gamma' \mathbin{.}
               \sum_\prime m \, \delta \gamma ),\cr}
   \right\}
   \eqno {\rm (Z^1.)}$$
and might have been otherwise deduced by our rule, from this other
known transformation of $T_\prime$,
$$T_\prime = {\textstyle {1 \over 2}}
      \sum_\prime \mathbin{.} m (\xi'^2 + \eta'^2 + \zeta'^2)
      - {(\sum_\prime m \xi')^2 + (\sum_\prime m \eta')^2
            + (\sum_\prime m \zeta')^2
      \over 2 \sum m}.
   \eqno {\rm (69.)}$$
And to obtain, with any set of internal or relative marks of
position, the two partial differential equations which the
characteristic function $V_\prime$ of relative motion must
satisfy, and which offer (as we shall find) the chief means of
discovering its form, namely, the equations analogous to those
marked (F.) and (G.), we have only to eliminate the rates of
increase of the marks of position of the system, which determine
the final and initial components of the relative velocities of
its points, by the law of varying relative action, from the final
and initial expressions of the law of relative living force;
namely, from the following equations:
$$T_\prime = U + H_\prime,
   \eqno {\rm (50.)}$$
and
$$T_{\prime 0} = U_0 + H_\prime.
   \eqno {\rm (70.)}$$

The law of areas, or the property respecting rotation which was
expressed by the partial differential equations (P.), will also
always admit of being expressed in relative coordinates, and will
assist in discovering the form of the characteristic function
$V_\prime$; by showing that this function involves only such
internal coordinates (in number $6n - 9$) as do not alter by any
common rotation of all points final and initial, round the centre
of gravity, or round any other internal origin; that origin being
treated as fixed, and the quantity $H_\prime$ as constant, in
determining the effects of this rotation.  The general problem of
dynamics, respecting the motions of a free system of $n$ points
attracting or repelling one another, is therefore reduced, in the
last analysis, by the method of the present essay, to the
research and differentiation of a function $V_\prime$, depending
on $6n - 9$ internal or relative coordinates, and on the quantity
$H_\prime$, and satisfying a pair of partial differential
equations of the first order and second degree; in integrating
which equations, we are to observe, that at the assumed origin of
the motion, namely at the moment when $t = 0$, the final or
variable coordinates are equal to their initial values, and the
partial differential coefficient
$\displaystyle {\delta V_\prime \over \delta H_\prime}$
vanishes; and, that at a moment infinitely little distant, the
differential alterations of the coordinates have ratios connected
with the other partial differential coefficients of the
characteristic function $V_\prime$, by the law of varying
relative action.  It may be here observed, that, although the
consideration of the point, called usually the centre of gravity,
is very simply suggested by the process of the tenth number, yet
this internal centre is even more simply indicated by our early
corollaries from the law of varying action; which show that the
components of relative final velocities, in any system of
attracting or repelling points, may be expressed by the
differences of quantities of the form
$\displaystyle {1 \over m} {\delta V \over \delta x}$,
$\displaystyle {1 \over m} {\delta V \over \delta y}$,
$\displaystyle {1 \over m} {\delta V \over \delta z}$:
and that therefore in calculating these relative velocities, it
is advantageous to introduce the final sums
$\sum m x$, $\sum m y$, $\sum m z$,
and, for an analogous reason, the initial sums
$\sum m a$, $\sum m b$, $\sum m c$,
among the marks of the extreme positions of the system, in the
expression of the characteristic function~$V$; because, in
differentiating that expression for the calculation of relative
velocities, those sums may be treated as constant.

\bigbreak

{\sectiontitle
On Systems of two Points, in general; Characteristic Function
of the motion of any Binary System.\par}

\nobreak\bigskip

13.
To illustrate the foregoing principles, which extend to any free
system of points, however numerous, attracting or repelling one
another, let us now consider, in particular, a system of two such
points.  For such a system, the known {\it force-function\/} $U$
becomes, by (2.)
$$U = m_1 m_2 f(r),
   \eqno {\rm (71.)}$$
$r$ being the mutual distance
$$r = \sqrt{ (x_1 - x_2)^2 + (y_1 - y_2)^2 + (z_1 - z_2)^2 }
   \eqno {\rm (72.)}$$
between the two points $m_1$, $m_2$, and $f(r)$ being a function
of this distance such that its derivative or differential
coefficient $f'(r)$ expresses the law of their repulsion or
attraction, according as it is positive or negative.  The known
differential equations of motion, of the second order, are now,
by (1.), comprised in the following formula:
$$m_1 (x_1'' \, \delta x_1 + y_1'' \, \delta y_1
                           + z'' \, \delta z_1) +
  m_2 (x_2'' \, \delta x_2 + y_2'' \, \delta y_2
                           + z'' \, \delta z_2)
  = m_1 m_2 \, \delta f(r);
   \eqno {\rm (73.)}$$
they are therefore, separately,
$$\left. \multieqalign{
x_1'' &= m_2 {\delta f(r) \over \delta x_1}, &
y_1'' &= m_2 {\delta f(r) \over \delta y_1}, &
z_1'' &= m_2 {\delta f(r) \over \delta z_1}, \cr
x_2'' &= m_1 {\delta f(r) \over \delta x_2}, &
y_2'' &= m_1 {\delta f(r) \over \delta y_2}, &
z_2'' &= m_1 {\delta f(r) \over \delta z_2}. \cr}
   \right\}
   \eqno {\rm (74.)}$$

The problem of integrating these equations consists in proposing
to assign, by their means, six relations between the time~$t$,
the masses $m_1$~$m_2$, the six varying coordinates
$x_1$~$y_1$~$z_1$ $x_2$~$y_2$~$z_2$,
and their initial values and initial rates of increase
$a_1$~$b_1$~$c_1$ $a_2$~$b_2$~$c_2$
$a_1'$~$b_1'$~$c_1'$ $a_2'$~$b_2'$~$c_2'$.
If we knew these six final integrals, and combined them with the
initial form of the law of living force, or of the known
intermediate integral
$${\textstyle {1 \over 2}} m_1 (x_1'^2 + y_1'^2 + z_1'^2) +
  {\textstyle {1 \over 2}} m_2 (x_2'^2 + y_2'^2 + z_2'^2)
  = m_1 m_2 f(r) + H;
   \eqno {\rm (75.)}$$
that is, with the following formula,
$${\textstyle {1 \over 2}} m_1 (a_1'^2 + b_1'^2 + c_1'^2) +
  {\textstyle {1 \over 2}} m_2 (a_2'^2 + b_2'^2 + c_2'^2)
  = m_1 m_2 f(r_0) + H,
   \eqno {\rm (76.)}$$
in which $r_0$ is the initial distance
$$r_0 = \sqrt{ (a_1 - a_2)^2 + (b_1 - b_2)^2 + (c_1 - c_2)^2 },
   \eqno {\rm (77.)}$$
and $H$ is a constant quantity, introduced by integration; we
could, by the combination of these seven relations, determine the
time~$t$, and the six initial components of velocity
$a_1'$~$b_1'$~$c_1'$ $a_2'$~$b_2'$~$c_2'$,
as functions of the twelve final and initial coordinates
$x_1$~$y_1$~$z_1$ $x_2$~$y_2$~$z_2$
$a_1$~$b_1$~$c_1$ $a_2$~$b_2$~$c_2$,
and of the quantity~$H$, (involving also the masses:) we could
therefore determine whatever else depends on the manner and time
of motion of this system of two points, as a function of the same
extreme coordinates and of the same quantity $H$.  In particular,
we could determine the action, or accumulated living force of the
system, namely,
$$V =   m_1 \int_0^t (x_1'^2 + y_1'^2 + z_1'^2) \,dt
      + m_2 \int_0^t (x_2'^2 + y_2'^2 + z_2'^2) \,dt,
   \eqno {\rm (A^2.)}$$
as a function of those thirteen quantities
$x_1$~$y_1$~$z_1$ $x_2$~$y_2$~$z_2$
$a_1$~$b_1$~$c_1$ $a_2$~$b_2$~$c_2$~$H$: and might then
calculate the variation of this function,
$$\left. \eqalign{
\delta V
   &= {\delta V \over \delta x_1} \, \delta x_1
    + {\delta V \over \delta y_1} \, \delta y_1
    + {\delta V \over \delta z_1} \, \delta z_1
    + {\delta V \over \delta x_2} \, \delta x_2
    + {\delta V \over \delta y_2} \, \delta y_2
    + {\delta V \over \delta z_2} \, \delta z_2 \cr
   &\mathrel{\phantom{=}} \mathord{}
    + {\delta V \over \delta a_1} \, \delta a_1
    + {\delta V \over \delta b_1} \, \delta b_1
    + {\delta V \over \delta c_1} \, \delta c_1
    + {\delta V \over \delta a_2} \, \delta a_2
    + {\delta V \over \delta b_2} \, \delta b_2
    + {\delta V \over \delta c_2} \, \delta c_2 \cr
   &\mathrel{\phantom{=}} \mathord{}
    + {\delta V \over \delta H} \, \delta H.\cr}
   \right\}
   \eqno {\rm (B^2.)}$$
But the essence of our method consists in {\it forming previously
the expression of this variation by our law of varying action},
namely,
$$\left. \eqalign{
\delta V
   &= m_1 ( x_1' \, \delta x_1 - a_1' \, \delta a_1
          + y_1' \, \delta y_1 - b_1' \, \delta b_1
          + z_1' \, \delta z_1 - c_1' \, \delta c_1 ) \cr
   &\mathrel{\phantom{=}} \mathord{}
    + m_2 ( x_2' \, \delta x_2 - a_2' \, \delta a_2
          + y_2' \, \delta y_2 - b_2' \, \delta b_2
          + z_2' \, \delta z_2 - c_2' \, \delta c_2 ) \cr
   &\mathrel{\phantom{=}} \mathord{}
    + t \, \delta H;\cr}
   \right\}
   \eqno {\rm (C^2.)}$$
and in {\it considering $V$ as a characteristic function of the
motion}, from the form of which may be deduced all the
intermediate and all the final integrals of the known
differential equations, by resolving the expression (C${}^2$.)
into the following separate groups, (included in (C.) and (D.),)
$$\left. \multieqalign{
{\delta V \over \delta x_1} &= m_1 x_1', &
{\delta V \over \delta y_1} &= m_1 y_1', &
{\delta V \over \delta z_1} &= m_1 z_1', \cr
{\delta V \over \delta x_2} &= m_2 x_2', &
{\delta V \over \delta y_2} &= m_2 y_2', &
{\delta V \over \delta z_2} &= m_2 z_2'; \cr}
   \right\}
   \eqno {\rm (D^2.)}$$
and
$$\left. \multieqalign{
{\delta V \over \delta a_1} &= - m_1 a_1', &
{\delta V \over \delta b_1} &= - m_1 b_1', &
{\delta V \over \delta c_1} &= - m_1 c_1', \cr
{\delta V \over \delta a_2} &= - m_2 a_2', &
{\delta V \over \delta b_2} &= - m_2 b_2', &
{\delta V \over \delta c_2} &= - m_2 c_2'; \cr}
   \right\}
   \eqno {\rm (E^2.)}$$
besides this other equation, which had occurred before,
$${\delta V \over \delta H} = t.
   \eqno {\rm (E.)}$$

By this new method, the difficulty of integrating the six known
equations of motion of the second order (74.) is reduced to the
search and differentiation of a single function~$V$; and to find
the form of this function, we are to employ the following pair of
partial differential equations of the first order:
$$\left. \eqalign{
   &{1 \over 2 m_1} \left\{
        \left( {\delta V \over \delta x_1} \right)^2
      + \left( {\delta V \over \delta y_1} \right)^2
      + \left( {\delta V \over \delta z_1} \right)^2
      \right\}
    + {1 \over 2 m_2} \left\{
        \left( {\delta V \over \delta x_2} \right)^2
      + \left( {\delta V \over \delta y_2} \right)^2
      + \left( {\delta V \over \delta z_2} \right)^2
      \right\} \cr
   &\qquad\qquad
   = m_1 m_2 f(r) + H,\cr}
   \right\}
   \eqno {\rm (F^2.)}$$
$$\left. \eqalign{
   &{1 \over 2 m_1} \left\{
        \left( {\delta V \over \delta a_1} \right)^2
      + \left( {\delta V \over \delta b_1} \right)^2
      + \left( {\delta V \over \delta c_1} \right)^2
      \right\}
    + {1 \over 2 m_2} \left\{
        \left( {\delta V \over \delta a_2} \right)^2
      + \left( {\delta V \over \delta b_2} \right)^2
      + \left( {\delta V \over \delta c_2} \right)^2
      \right\} \cr
   &\qquad\qquad
   = m_1 m_2 f(r_0) + H,\cr}
   \right\}
   \eqno {\rm (G^2.)}$$
combined with some simple considerations.  And it easily results
from the principles already laid down, that the integral of this
pair of equations, adapted to the present question, is
$$\left. \eqalign{
V
   &= \sqrt{   (x_{\prime\prime} - a_{\prime\prime})^2
             + (y_{\prime\prime} - b_{\prime\prime})^2
             + (z_{\prime\prime} - c_{\prime\prime})^2 }
      \mathbin{.} \sqrt{ 2 H_{\prime\prime} (m_1 + m_2) } \cr
   &\mathrel{\phantom{=}} \mathord{}
   + {m_1 m_2 \over m_1 + m_2} \left(
         h \vartheta + \int_{r_0}^r \rho \,dr \right);\cr}
   \right\}
   \eqno {\rm (H^2.)}$$
in which
$x_{\prime\prime}$~$y_{\prime\prime}$~$z_{\prime\prime}$
$a_{\prime\prime}$~$b_{\prime\prime}$~$c_{\prime\prime}$
denote the coordinates, final and initial, of the centre of
gravity of the system,
$$\left. \multieqalign{
x_{\prime\prime} &= {m_1 x_1 + m_2 x_2 \over m_1 + m_2}, &
y_{\prime\prime} &= {m_1 y_1 + m_2 y_2 \over m_1 + m_2}, &
z_{\prime\prime} &= {m_1 z_1 + m_2 z_2 \over m_1 + m_2}, \cr
a_{\prime\prime} &= {m_1 a_1 + m_2 a_2 \over m_1 + m_2}, &
b_{\prime\prime} &= {m_1 b_1 + m_2 b_2 \over m_1 + m_2}, &
c_{\prime\prime} &= {m_1 c_1 + m_2 c_2 \over m_1 + m_2}, \cr}
   \right\}
   \eqno {\rm (78.)}$$
and $\vartheta$ is the angle between the final and initial
distances $r$, $r_0$: we have also put for abridgement
$$\rho = \pm \sqrt{ 2( m_1 + m_2 ) \left(
         f(r) + {H_\prime \over m_1 m_2} \right)
         - {h^2 \over r^2}},
   \eqno {\rm (79.)}$$
the upper or the lower sign to be used, according as the distance
$r$ is increasing or decreasing, and have introduced three
auxiliary quantities $h$, $H_\prime$, $H_{\prime\prime}$, to be
determined by this condition,
$$0 = \vartheta + \int_{r_0}^r {\delta \rho \over \delta h} dr,
   \eqno {\rm (I^2.)}$$
combined with the two following,
$$\left. \eqalign{
{m_1 m_2 \over m_1 + m_2} \int_{r_0}^r
      {\delta \rho \over \delta H_\prime} dr
   &= \sqrt{   (x_{\prime\prime} - a_{\prime\prime})^2
             + (y_{\prime\prime} - b_{\prime\prime})^2
             + (z_{\prime\prime} - c_{\prime\prime})^2 }
      \mathbin{.} \sqrt{ {m_1 + m_2 \over 2 H_{\prime\prime}} },\cr
H_\prime + H_{\prime\prime} &= H;\cr}
   \right\}
   \eqno {\rm (K^2.)}$$
which auxiliary quantities, although in one view they are
functions of the twelve extreme coordinates, are yet to be
treated as constant in calulating the three definite integrals,
or limits of sums of numerous small elements,
$$\int_{r_0}^r \rho \, dr,\quad
  \int_{r_0}^r {\delta \rho \over \delta h} dr,\quad
  \int_{r_0}^r {\delta \rho \over \delta H_\prime} dr.$$

The form (H${}^2$.), for the {\it characteristic function of a
binary system}, may be regarded as a central or radical relation,
which includes the whole theory of the motion of such a system;
so that all the details of this motion may be deduced from it by
the application of our general method.  But because the theory of
binary systems has been brought to great perfection already, by
the labours of former writers, it may suffice to give briefly
here a few instances of such deduction.

\bigbreak

14.
The form (H${}^2$.), for the characteristic function of a binary
system, involves explicitly, when $\rho$ is changed to its value
(79.), the twelve quantities
$x_{\prime\prime}$~$y_{\prime\prime}$~$z_{\prime\prime}$
$a_{\prime\prime}$~$b_{\prime\prime}$~$c_{\prime\prime}$
$r$~$r_0$~$\vartheta$~$h$~$H_\prime$~$H_{\prime\prime}$,
(besides the masses $m_1$~$m_2$ which are always considered as
given;) its variation may therefore be thus expressed:
$$\left. \eqalign{
\delta V
   &=   {\delta V \over \delta x_{\prime\prime}} \delta x_{\prime\prime}
      + {\delta V \over \delta y_{\prime\prime}} \delta y_{\prime\prime}
      + {\delta V \over \delta z_{\prime\prime}} \delta z_{\prime\prime}
      + {\delta V \over \delta a_{\prime\prime}} \delta a_{\prime\prime}
      + {\delta V \over \delta b_{\prime\prime}} \delta b_{\prime\prime}
      + {\delta V \over \delta c_{\prime\prime}} \delta c_{\prime\prime} \cr
   &\mathrel{\phantom{=}} \mathord{}
      + {\delta V \over \delta r} \delta r
      + {\delta V \over \delta r_0} \delta r_0
      + {\delta V \over \delta \vartheta} \delta \vartheta
      + {\delta V \over \delta H_\prime} \delta H_\prime
      + {\delta V \over \delta H_{\prime\prime}} \delta H_{\prime\prime}.\cr}
   \right\}
   \eqno {\rm (L^2.)}$$

In this expression, if we put for abridgement
$$\lambda = \sqrt{ {2 H_{\prime\prime} (m_1 + m_2) \over
               (x_{\prime\prime} - a_{\prime\prime})^2
             + (y_{\prime\prime} - b_{\prime\prime})^2
             + (z_{\prime\prime} - c_{\prime\prime})^2 }},
   \eqno {\rm (80.)}$$
we shall have
$$\left. \multieqalign{
{\delta V \over \delta x_{\prime\prime}}
   &= \lambda (x_{\prime\prime} - a_{\prime\prime}), &
{\delta V \over \delta y_{\prime\prime}}
   &= \lambda (y_{\prime\prime} - b_{\prime\prime}), &
{\delta V \over \delta z_{\prime\prime}}
   &= \lambda (z_{\prime\prime} - c_{\prime\prime}),\cr
{\delta V \over \delta a_{\prime\prime}}
   &= \lambda (a_{\prime\prime} - x_{\prime\prime}), &
{\delta V \over \delta b_{\prime\prime}}
   &= \lambda (b_{\prime\prime} - y_{\prime\prime}), &
{\delta V \over \delta c_{\prime\prime}}
   &= \lambda (c_{\prime\prime} - z_{\prime\prime});\cr}
   \right\}
   \eqno {\rm (M^2.)}$$
and if we put
$$\rho_0 = \pm \sqrt{ 2(m_1 + m_2)
               \left( f(r_0) + {H_\prime \over m_1 m_2} \right)
               - {h_2 \over r_0^2} },
   \eqno {\rm (81.)}$$
the sign of the radical being determined by the same rule as that
of $\rho$, we shall have
$${\delta V \over \delta r}
      = {m_1 m_2 \rho \over m_1 + m_2},\quad
  {\delta V \over \delta r_0}
      = {- m_1 m_2 \rho_0 \over m_1 + m_2},\quad
  {\delta V \over \delta \vartheta}
      = {m_1 m_2 h \over m_1 + m_2};
   \eqno {\rm (N^2.)}$$
besides, by the equations of condition (I${}^2$.), (K${}^2$.), we
have
$${\delta V \over \delta h} = 0,
   \eqno {\rm (O^2.)}$$
and
$${\delta V \over \delta H_{\prime\prime}}
      = {\delta V \over \delta H_\prime}
      = \int_{r_0}^r {dr \over \rho},\quad
  \delta H_\prime + \delta H_{\prime\prime} = \delta H.
   \eqno {\rm (P^2.)}$$

The expression (L${}^2$.) may therefore be thus transformed:
$$\left. \eqalign{
\delta V
   &= \lambda \{
            (x_{\prime\prime} - a_{\prime\prime})
            (\delta x_{\prime\prime} - \delta a_{\prime\prime})
          + (y_{\prime\prime} - b_{\prime\prime})
            (\delta y_{\prime\prime} - \delta b_{\prime\prime})
          + (z_{\prime\prime} - c_{\prime\prime})
            (\delta z_{\prime\prime} - \delta c_{\prime\prime})
       \} \cr
   &\mathrel{\phantom{=}} \mathord{}
      + {m_1 m_2 \over m_1 + m_2}
         (\rho \, \delta r - \rho_0 \, \delta r_0 + h \, \delta \vartheta)
      + \int_{r_0}^r {\delta r \over \rho} \mathbin{.} \delta H;\cr}
   \right\}
   \eqno {\rm (Q^2.)}$$
and may be resolved by our general method into twelve separate
expressions for the final and initial components of velocities,
namely,
$$\left. \eqalign{
x_1' &= {1 \over m_1} {\delta V \over \delta x_1}
      = {\lambda \over m_1 + m_2} (x_{\prime\prime} - a_{\prime\prime})
         + {m_2 \over m_1 + m_2}
         \left( \rho {\delta r \over \delta x_1}
                + h {\delta \vartheta \over \delta x_1}
                \right),\cr
y_1' &= {1 \over m_1} {\delta V \over \delta y_1}
      = {\lambda \over m_1 + m_2} (y_{\prime\prime} - b_{\prime\prime})
         + {m_2 \over m_1 + m_2}
         \left( \rho {\delta r \over \delta y_1}
                + h {\delta \vartheta \over \delta y_1}
                \right),\cr
z_1' &= {1 \over m_1} {\delta V \over \delta z_1}
      = {\lambda \over m_1 + m_2} (z_{\prime\prime} - c_{\prime\prime})
         + {m_2 \over m_1 + m_2}
         \left( \rho {\delta r \over \delta z_1}
                + h {\delta \vartheta \over \delta z_1}
                \right),\cr
x_2' &= {1 \over m_2} {\delta V \over \delta x_2}
      = {\lambda \over m_1 + m_2} (x_{\prime\prime} - a_{\prime\prime})
         + {m_1 \over m_1 + m_2}
         \left( \rho {\delta r \over \delta x_2}
                + h {\delta \vartheta \over \delta x_2}
                \right),\cr
y_2' &= {1 \over m_2} {\delta V \over \delta y_2}
      = {\lambda \over m_1 + m_2} (y_{\prime\prime} - b_{\prime\prime})
         + {m_1 \over m_1 + m_2}
         \left( \rho {\delta r \over \delta y_2}
                + h {\delta \vartheta \over \delta y_2}
                \right),\cr
z_2' &= {1 \over m_2} {\delta V \over \delta z_2}
      = {\lambda \over m_1 + m_2} (z_{\prime\prime} - c_{\prime\prime})
         + {m_1 \over m_1 + m_2}
         \left( \rho {\delta r \over \delta z_2}
                + h {\delta \vartheta \over \delta z_2}
                \right),\cr}
   \right\}
   \eqno {\rm (R^2.)}$$
and
$$\left. \eqalign{
a_1' &= {-1 \over m_1} {\delta V \over \delta a_1}
      = {\lambda \over m_1 + m_2} (x_{\prime\prime} - a_{\prime\prime})
         + {m_2 \over m_1 + m_2}
         \left( \rho_0 {\delta r_0 \over \delta a_1}
                - h {\delta \vartheta \over \delta a_1}
                \right),\cr
b_1' &= {-1 \over m_1} {\delta V \over \delta b_1}
      = {\lambda \over m_1 + m_2} (y_{\prime\prime} - b_{\prime\prime})
         + {m_2 \over m_1 + m_2}
         \left( \rho_0 {\delta r_0 \over \delta b_1}
                - h {\delta \vartheta \over \delta b_1}
                \right),\cr
c_1' &= {-1 \over m_1} {\delta V \over \delta c_1}
      = {\lambda \over m_1 + m_2} (z_{\prime\prime} - c_{\prime\prime})
         + {m_2 \over m_1 + m_2}
         \left( \rho_0 {\delta r_0 \over \delta c_1}
                - h {\delta \vartheta \over \delta c_1}
                \right),\cr
a_2' &= {-1 \over m_2} {\delta V \over \delta a_2}
      = {\lambda \over m_1 + m_2} (x_{\prime\prime} - a_{\prime\prime})
         + {m_1 \over m_1 + m_2}
         \left( \rho_0 {\delta r_0 \over \delta a_2}
                - h {\delta \vartheta \over \delta a_2}
                \right),\cr
b_2' &= {-1 \over m_2} {\delta V \over \delta b_2}
      = {\lambda \over m_1 + m_2} (y_{\prime\prime} - b_{\prime\prime})
         + {m_1 \over m_1 + m_2}
         \left( \rho_0 {\delta r_0 \over \delta b_2}
                - h {\delta \vartheta \over \delta b_2}
                \right),\cr
c_2' &= {-1 \over m_2} {\delta V \over \delta c_2}
      = {\lambda \over m_1 + m_2} (z_{\prime\prime} - c_{\prime\prime})
         + {m_1 \over m_1 + m_2}
         \left( \rho_0 {\delta r_0 \over \delta c_2}
                - h {\delta \vartheta \over \delta c_2}
                \right);\cr}
   \right\}
   \eqno {\rm (S^2.)}$$
besides the following expression for the time of motion of the
system:
$$t = {\delta V \over \delta H} = \int_{r_0}^r {dr \over \rho},
   \eqno {\rm (T^2.)}$$
which gives by (K${}^2$.), and by (79.), (80.),
$$t = {m_1 + m_2 \over \lambda}.
   \eqno {\rm (U^2.)}$$

The six equations (R${}^2$.) give the six intermediate integrals,
and the six equations (S${}^2$.) give the six final integrals of
the six known differential equations of motion (74.) for any
binary system, if we eliminate or determine the three auxiliary
quantities $h$, $H_\prime$, $H_{\prime\prime}$, by the three
conditions (I${}^2$.) (T${}^2$.) (U${}^2$.).  Thus, if we observe
that the distances $r$, $r_0$, and the included angle
$\vartheta$, depend only on relative coordinates, which may be
thus denoted,
$$\left. \multieqalign{
   x_1 - x_2 &= \xi,    &
   y_1 - y_2 &= \eta,   &
   z_1 - z_2 &= \zeta,  \cr
   a_1 - a_2 &= \alpha, &
   b_1 - b_2 &= \beta,  &
   c_1 - c_2 &= \gamma, \cr}
   \right\}
   \eqno {\rm (82.)}$$
we obtain by easy combinations the three following intermediate
integrals for the centre of gravity of the system:
$$x_{\prime\prime}' t = x_{\prime\prime} - a_{\prime\prime},\quad
  y_{\prime\prime}' t = y_{\prime\prime} - b_{\prime\prime},\quad
  z_{\prime\prime}' t = z_{\prime\prime} - c_{\prime\prime},
   \eqno {\rm (83.)}$$
and the three following final integrals,
$$a_{\prime\prime}' t = x_{\prime\prime} - a_{\prime\prime},\quad
  b_{\prime\prime}' t = y_{\prime\prime} - b_{\prime\prime},\quad
  c_{\prime\prime}' t = z_{\prime\prime} - c_{\prime\prime},
   \eqno {\rm (84.)}$$
expressing the well-known law of the rectilinear and uniform
motion of that centre.  We obtain also the three following
intermediate integrals for the relative motion of one point of
the system about the other:
$$\left. \eqalign{
\xi'    &= \rho {\delta r \over \delta \xi}
            + h {\delta \vartheta \over \delta \xi},\cr
\eta'   &= \rho {\delta r \over \delta \eta}
            + h {\delta \vartheta \over \delta \eta},\cr
\zeta'  &= \rho {\delta r \over \delta \zeta}
            + h {\delta \vartheta \over \delta \zeta},\cr}
   \right\}
   \eqno {\rm (85.)}$$
and the three following final integrals,
$$\left. \eqalign{
\alpha' &= \rho_0 {\delta r_0 \over \delta \alpha}
            - h {\delta \vartheta \over \delta \alpha},\cr
\beta'  &= \rho_0 {\delta r_0 \over \delta \beta}
            - h {\delta \vartheta \over \delta \beta},\cr
\gamma' &= \rho_0 {\delta r_0 \over \delta \gamma}
            - h {\delta \vartheta \over \delta \gamma};\cr}
   \right\}
   \eqno {\rm (86.)}$$
in which the auxiliary quantities $h$, $H_\prime$ are to be
determined by (I${}^2$.), (T${}^2$.), and in which the dependence
of $r$, $r_0$, $\vartheta$, on $\xi$, $\eta$, $\zeta$, $\alpha$,
$\beta$, $\gamma$, is expressed by the following equations:
$$\left. \eqalign{
   & r   = \sqrt{ \xi^2 + \eta^2 + \zeta^2},\quad
     r_0 = \sqrt{ \alpha^2 + \beta^2 + \gamma^2},\cr
   & r r_0 \cos \vartheta
      = \xi \alpha + \eta \beta + \zeta \gamma.\cr}
   \right\}
   \eqno {\rm (87.)}$$
If then we put, for abridgement,
$$A = {\rho \over r} + {h \over r^2 \tan \vartheta},\quad
  B = {h \over r r_0 \sin \vartheta},\quad
  C = {- \rho_0 \over r_0} + {h \over r_0^2 \tan \vartheta},
   \eqno {\rm (88.)}$$
we shall have these three intermediate integrals,
$$\xi'    = A \xi - B \alpha,\quad
  \eta'   = A \eta - B \beta,\quad
  \zeta'  = A \zeta - B \gamma,
   \eqno {\rm (89.)}$$
and these three final integrals,
$$\alpha'  = B \xi - C \alpha,\quad
  \beta'   = B \eta - C \beta,\quad
  \gamma'  = B \zeta - C \gamma,
   \eqno {\rm (90.)}$$
of the equations of relative motion.  These integrals give,
$$\left. \eqalign{
\xi \eta' - \eta \xi'
   = \alpha \beta' - \beta \alpha'
   = B (\alpha \eta - \beta \xi),\cr
\eta \zeta' - \zeta \eta'
   = \beta \gamma' - \gamma \beta'
   = B (\beta \zeta - \gamma \eta),\cr
\zeta \xi' - \xi \zeta'
   = \gamma \alpha' - \alpha \gamma'
   = B (\gamma \xi - \alpha \zeta),\cr}
   \right\}
   \eqno {\rm (91.)}$$
and
$$\zeta (\alpha \beta' - \beta \alpha')
      + \xi (\beta \gamma' - \gamma \beta')
      + \eta (\gamma \alpha' - \alpha \gamma')
   = 0;
   \eqno {\rm (92.)}$$
they contain therefore the known law of equable description of
areas, and the law of a plane relative orbit.  If we take for
simplicity this plane for the plane $\xi$~$\eta$, the quantities
$\zeta$~$\zeta'$~$\gamma$~$\gamma'$ will vanish; and we may
put,
$$\left. \multieqalign{
\xi    &= r \cos \theta, &
\eta   &= r \sin \theta, &
\zeta  &= 0,\cr
\alpha &= r_0 \cos \theta_0, &
\beta  &= r_0 \sin \theta_0, &
\gamma &= 0,\cr}
   \right\}
   \eqno {\rm (93.)}$$
and
$$\left. \multieqalign{
\xi'    &= r' \cos \theta - \theta' r \sin \theta, &
\eta'   &= r' \sin \theta + \theta' r \cos \theta, &
\zeta'  &= 0,\cr
\alpha' &= r_0' \cos \theta_0 - \theta_0' r_0 \sin \theta_0, &
\beta'  &= r_0' \sin \theta_0 + \theta_0' r_0 \cos \theta_0, &
\gamma' &= 0,\cr}
   \right\}
   \eqno {\rm (94.)}$$
the angles $\theta$~$\theta_0$ being counted from some fixed
line in the plane, and being such that their difference
$$\theta - \theta_0 = \vartheta.
   \eqno {\rm (95.)}$$
These values give
$$\xi \eta' - \eta \xi' = r^2 \theta',\quad
  \alpha \beta' - \beta \alpha' = r_0^2 \theta_0',\quad
  \alpha \eta - \beta \xi = r r_0 \sin \vartheta,
   \eqno {\rm (96.)}$$
and therefore, by (88.) and (91.),
$$r^2 \theta' = r_0^2 \theta_0' = h;
   \eqno {\rm (97.)}$$
the quantity ${1 \over 2} h$ is therefore the constant areal
velocity in the relative motion of the system; a result which is
easily seen to be independent of the directions of the three
rectangular coordinates.  The same values (93.), (94.), give
$$\left. \multieqalign{
\xi \cos \theta + \eta \sin \theta &= r, &
\xi' \cos \theta + \eta' \sin \theta &= r', &
\alpha \cos \theta + \beta \sin \theta &= r_0 \cos \vartheta, \cr
\alpha \cos \theta_0 + \beta \sin \theta_0 &= r_0, &
\alpha' \cos \theta_0 + \beta' \sin \theta_0 &= r_0', &
\xi \cos \theta_0 + \eta \sin \theta_0 &= r \cos \vartheta, \cr}
   \right\}
   \eqno {\rm (98.)}$$
and therefore, by the intermediate and final integrals, (89.),
(90.),
$$r' = \rho,\quad r_0' = \rho_0;
   \eqno {\rm (99.)}$$
results which evidently agree with the condition (T${}^2$.), and
which give by (79.) and (81.), for all directions of coordinates,
$$r'^2 + {h^2 \over r^2} - 2(m_1 + m_2) f(r)
   = r_0'^2 + {h^2 \over r_0^2} - 2(m_1 + m_2) f(r_0)
   = 2 H_\prime \left( {1 \over m_1} + {1 \over m_2} \right);
   \eqno {\rm (100.)}$$
the other auxiliary quantity $H_\prime$ is therefore also a
constant, independent of the time, and enters as such into the
constant part in the expression for
$\displaystyle \left( r'^2 + {h^2 \over r^2} \right)$
the square of the relative velocity.  The equation of condition
(I${}^2$.), connecting these two constants $h$, $H_\prime$, with
the extreme lengths of the radius vector~$r$, and with the angle
$\vartheta$ described by this radius in revolving from its
initial to its final direction, is the equation of the plane
relative orbit; and the other equation of condition (T${}^2$.),
connecting the same two constants with the same extreme distances
and with the time, gives the law of the velocity of mutual
approach or recess.

We may remark that the part $V_\prime$ of the whole
characteristic function $V$, which represents the relative action
and determines the relative motion in the system, namely,
$$V_\prime = {m_1 m_2 \over m_1 + m_2}
      \left( h \vartheta + \int_{r_0}^r \rho \, dr \right),
   \eqno {\rm (V^2.)}$$
may be put, by (I${}^2$.), under the form
$$V_\prime = {m_1 m_2 \over m_1 + m_2} \int_{r_0}^r
         \left( \rho - h {\delta \rho \over \delta h} \right) \, dr,
   \eqno {\rm (W^2.)}$$
or finally, by (79.)
$$V_\prime
   = 2 \int_{r_0}^r {m_1 m_2 f(r) + H_\prime \over \rho} \, dr;
   \eqno {\rm (X^2.)}$$
the condition (I${}^2$.) may also itself be transformed, by
(79.), as follows:
$$\vartheta = h \int_{r_0}^r {dr \over r^2 \rho}:
   \eqno {\rm (Y^2.)}$$
results which all admit of easy verifications.  The partial
differential equations connected with the law of relative living
force, which the characteristic function $V_\prime$ of relative
motion must satisfy, may be put under the following forms:
$$\left. \eqalign{
\left( {\delta V_\prime \over \delta r} \right)^2
   + {1 \over r^2}
      \left( {\delta V_\prime \over \delta \vartheta} \right)^2
   &= {2 m_1 m_2 \over m_1 + m_2} (U + H_\prime),\cr
\left( {\delta V_\prime \over \delta r_0} \right)^2
   + {1 \over r_0^2}
      \left( {\delta V_\prime \over \delta \vartheta} \right)^2
   &= {2 m_1 m_2 \over m_1 + m_2} (U_0 + H_\prime);\cr}
   \right\}
   \eqno {\rm (Z^2.)}$$
and if the first of the equations of this pair have its variation
taken with respect to $r$ and $\vartheta$, attention being paid
to the dynamical meanings of the coefficients of the
characteristic function, it will conduct (as in former instances)
to the known differential equations of motion of the second
order.

\bigbreak

{\sectiontitle
On the undisturbed Motion of a Planet or Comet about the Sun:
Dependence of the Characteristic Function of such Motion,
on the chord and the sum of the Radii.\par}

\nobreak\bigskip

15.
To particularize still further, let
$$f(r) = {1 \over r},
   \eqno {\rm (101.)}$$
that is, let us consider a binary system, such as a planet or
comet and the sun, with the Newtonian law of attraction; and let
us put, for abridgement,
$$m_1 + m_2 = \mu,\quad
  {h^2 \over \mu} = p,\quad
  {-m_1 m_2 \over 2 H_\prime} = {\rm a}.
   \eqno {\rm (102.)}$$

The characteristic function $V_\prime$ of relative motion may now
be expressed as follows
$$V_\prime = {m_1 m_2 \over \sqrt{\mu}}
         \left( \vartheta \sqrt{p}
                + \int_{r_0}^r \pm \sqrt{
                  {2 \over r} - {1 \over {\rm a}}
                - {p \over r^2} }.dr \right);
   \eqno {\rm (A^3.)}$$
in which $p$ is to be considered as a function of the extreme
radii vectores $r$, $r_0$, and of their included angle
$\vartheta$, involving also the quantity ${\rm a}$, or the
connected quantity $H_\prime$, and determined by the condition
$$\vartheta = \int_{r_0}^r {\pm dr \over \displaystyle
               r^2 \sqrt{ {2 \over rp} - {1 \over {\rm a}p}
                  - {1 \over r^2} }}
   \eqno {\rm (B^3.)}$$
that is, by the derivative of the formula (A${}^3$.), taken with
respect to $p$; the upper sign being taken in each expression
when the distance~$r$ is increasing, and the lower sign when that
distance is diminishing, and the quantity $p$ being treated as
constant in calculating the two definite integrals.  It results
from the foregoing remarks, that this quantity $p$ is constant
also in the sense of being independent of the time, so as not to
vary in the course of the motion; and that the condition
(B${}^3$.), connecting this constant with
$r$~$r_0$~$\vartheta$~${\rm a}$,
is the equation of the plane relative orbit; which is therefore
(as it has long been known to be) an ellipse, hyperbola, or
parabola, according as the constant~${\rm a}$ is positive,
negative, or zero, the origin of $r$ being always a focus of the
curve, and $p$ being the semiparameter.  It results also, that
the time of motion may be thus expressed:
$$t = {\delta V_\prime \over \delta H_\prime}
    = {2 {\rm a}^2 \over m_1 m_2} {\delta V \over \delta {\rm a}},
   \eqno {\rm (C^3.)}$$
and therefore thus:
$$t = \int_{r_0}^r {\pm dr \over \displaystyle
               \sqrt{ {2\mu \over r} - {\mu \over {\rm a}}
                  - {\mu p \over r^2} }};
   \eqno {\rm (D^3.)}$$
which latter is a known expression.  Confining ourselves at
present to the case ${\rm a} > 0$, and introducing the known
auxiliary quantities called excentricity and excentric anomaly,
namely,
$$e = \sqrt{ 1 - {p \over {\rm a}} },
   \eqno {\rm (103.)}$$
and
$$\upsilon = \cos^{-1} \left( {{\rm a} - r \over {\rm a}e} \right),
   \eqno {\rm (104.)}$$
which give
$$\pm \sqrt{2 {\rm a} r - r^2 - p {\rm a}}
   = {\rm a}e \sin \upsilon,
   \eqno {\rm (105.)}$$
$\upsilon$ being considered as continually increasing with the time;
and therefore, as is well known,
$$\left. \eqalign{
r &= {\rm a} (1 - e \cos \upsilon),\quad
   r_0 = {\rm a} (1 - e \cos \upsilon_0),\cr
\vartheta &= 2 \tan^{-1} \left\{
      \sqrt{1 + e \over 1 - e} \tan {\upsilon \over 2} \right\}
      - 2 \tan^{-1} \left\{
      \sqrt{1 + e \over 1 - e} \tan {\upsilon_0 \over 2} \right\},\cr}
   \right\}
   \eqno {\rm (106.)}$$
and
$$t = \sqrt{ {{\rm a}^3 \over \mu} } \mathbin{.} (\upsilon - \upsilon_0
      - e \sin \upsilon + e \sin \upsilon_0);
   \eqno {\rm (107.)}$$
we find that this expression for the characteristic function of
relative motion,
$$V_\prime = {m_1 m_2 \over \sqrt{\mu}} \int_{r_0}^r
      {\displaystyle \pm \left( {2 \over r}
         - {1 \over {\rm a}} \right) \, dr
         \over \displaystyle \sqrt{ {2 \over r}
            - {1 \over {\rm a}} - {p \over r^2}} },
   \eqno {\rm (E^3.)}$$
deduced from (A${}^3$.) and (B${}^3$.), may be transformed as
follows:
$$V_\prime  = m_1 m_2 \sqrt{ {{\rm a} \over \mu} } (\upsilon - \upsilon_0
      + e \sin \upsilon - e \sin \upsilon_0):
   \eqno {\rm (F^3.)}$$
in which the excentricity $e$, and the final and initial
excentric anomalies $\upsilon$, $\upsilon_0$, are to be considered as
functions of the final and initial radii $r$, $r_0$, and of the
included angle $\vartheta$, determined by the equations (106.).
The expression (F${}^3$.) may be thus written:
$$V_\prime = 2 m_1 m_2 \sqrt{ {{\rm a} \over \mu} }
      (\upsilon_\prime + e_\prime \sin \upsilon_\prime),
   \eqno {\rm (G^3.)}$$
if we put, for abridgement,
$$\upsilon_\prime = {\upsilon - \upsilon_0 \over 2},\quad
   e_\prime = e \cos {\upsilon + \upsilon_0 \over 2};
   \eqno {\rm (108.)}$$
for the complete determination of the characteristic function of
the present relative motion, it remains therefore to determine
the two variables $\upsilon_\prime$ and $e_\prime$, as functions of
$r$~$r_0$~$\vartheta$, or of some other set of quantities which
mark the shape and size of the plane triangle bounded by the final
and initial elliptic radii vectores and by the elliptic chord.

For this purpose it is convenient to introduce this elliptic
chord itself, which we shall call $\pm \tau$, so that
$$\tau^2 = r^2 + r_0^2 - 2 r r_0 \cos \vartheta;
   \eqno {\rm (109.)}$$
because this chord may be expressed as a function of the two
variables $\upsilon_\prime$, $e_\prime$, (involving also the mean
distance $a_\prime$,) as follows.  The value (106.) for the angle
$\vartheta$, that is, by (95.), for $\theta - \theta_0$, gives
$$\theta - 2 \tan^{-1} \left\{
      \sqrt{ {1 + e \over 1 - e} } \tan {\upsilon \over 2} \right\}
  = \theta_0 - 2 \tan^{-1} \left\{
      \sqrt{ {1 + e \over 1 - e} } \tan {\upsilon_0 \over 2} \right\}
   = \varpi,
   \eqno {\rm (110.)}$$
$\varpi$ being a new constant independent of the time, namely,
one of the values of the polar angle $\theta$, which correspond
to the minimum of radius vector; and therefore, by (106.),
$$\left. \multieqalign{
r \cos (\theta - \varpi)
   &= {\rm a} (\cos \upsilon - e), &
r \sin (\theta - \varpi)
   &= {\rm a} \sqrt{ 1 - e^2 } \sin \upsilon, \cr
r_0 \cos (\theta_0 - \varpi)
   &= {\rm a} (\cos \upsilon_0 - e), &
r_0 \sin (\theta_0 - \varpi)
   &= {\rm a} \sqrt{ 1 - e^2 } \sin \upsilon_0; \cr}
   \right\}
   \eqno {\rm (111.)}$$
expressions which give the following value for the square of the
elliptic chord:
$$\left. \eqalign{
\tau^2
   &= \{ r \cos (\theta - \varpi)
         - r_0 \cos (\theta_0 - \varpi) \}^2
      + \{ r \sin (\theta - \varpi)
         - r_0 \sin (\theta_0 - \varpi) \}^2 \cr
   &= {\rm a}^2 \{ (\cos \upsilon - \cos \upsilon_0)^2
         + (1 - e^2) (\sin \upsilon - \sin \upsilon_0)^2 \} \cr
   &= 4 {\rm a}^2 \sin \upsilon_\prime^2 \left\{
      \left( \sin {\upsilon + \upsilon_0 \over 2} \right)^2
      + (1 - e^2) \left( \cos {\upsilon + \upsilon_0 \over 2} \right)^2
      \right\} \cr
   &= 4 {\rm a}^2 (1 - e_\prime^2) \sin \upsilon_\prime^2:\cr}
   \right\}
   \eqno {\rm (112.)}$$
we may also consider $\tau$ as having the same sign with
$\sin \upsilon_\prime$, if we consider it as alternately positive and
negative, in the successive elliptic periods or revolutions,
beginning with the initial position.

Besides, if we denote by $\sigma$ the sum of the two elliptic
radii vectores, final and initial, so that
$$\sigma = r + r_0,
   \eqno {\rm (113.)}$$
we shall have, with our present abridgements,
$$\sigma = 2 {\rm a} (1 - e_\prime \cos \upsilon_\prime );
   \eqno {\rm (114.)}$$
the variables $\upsilon_\prime$~$e_\prime$ are therefore functions of
$\sigma$, $\tau$, ${\rm a}$, and consequently the characteristic
function $V_\prime$ is itself a function of those three
quantities.  We may therefore put
$$V_\prime = {m_1 m_2 w \over m_1 + m_2},
   \eqno {\rm (H^3.)}$$
$w$ being a function of $\sigma$, $\tau$, ${\rm a}$, of which the
form is to be determined by eliminating $\upsilon_\prime$~$e_\prime$
between the three equations,
$$\left. \eqalign{
w &= 2 \sqrt{\mu {\rm a}}
      (\upsilon_\prime + e_\prime \sin \upsilon_\prime),\cr
\sigma &= 2 {\rm a} (1 - e_\prime \cos \upsilon_\prime),\cr
\tau &= 2 {\rm a} (1 - e_\prime^2)^{1 \over 2}
      \sin \upsilon_\prime;\cr}
   \right\}
   \eqno {\rm (I^3.)}$$
and we may consider this new function $w$ as itself a
characteristic function of elliptic motion; the law of its
variation being expressed as follows, in the notation of the
present essay:
$$\delta w = \xi' \, \delta \xi - \alpha' \, \delta \alpha
      + \eta' \, \delta \eta - \beta' \, \delta \beta
      + \zeta' \, \delta \zeta - \gamma' \, \delta \gamma
      + {t \mu \, \delta {\rm a} \over 2 {\rm a}^2}.
   \eqno {\rm (K^3.)}$$
In this expression $\xi$~$\eta$~$\zeta$ are the relative
coordinates of the point $m_1$, at the time $t$, referred to the
other attracting point $m_2$ as an origin, and to any three
rectangular axes; $\xi'$~$\eta'$~$\zeta'$ are their rates of
increase, or the three rectangular components of final relative
velocity; $\alpha$~$\beta$~$\gamma$ $\alpha'$~$\beta'$~$\gamma'$
are the initial values, or values at the time zero, of
these relative coordinates and components of relative velocity;
${\rm a}$ is a quantity independent of the time, namely, the mean
distance of the two points $m_1$, $m_2$; and $\mu$ is the sum of
their masses.  And all the properties of the undisturbed elliptic
motion of a planet or comet about the sun may be deduced in a new
way, from the simplified characteristic function $w$, by
comparing its variation (K${}^3$.) with the following other form,
$$\delta w
   =   {\delta w \over \delta \sigma}  \delta \sigma
     + {\delta w \over \delta \tau}    \delta \tau
     + {\delta w \over \delta {\rm a}} \delta {\rm a};
   \eqno {\rm (L^3.)}$$
in which we are to observe that
$$\left. \eqalign{
\sigma &= \sqrt{ \xi^2 + \eta^2 + \zeta^2 }
          + \sqrt{ \alpha^2 + \beta^2 + \gamma^2 },\cr
\tau &= \pm \sqrt{ (\xi - \alpha)^2 + (\eta - \beta)^2
                   + (\zeta - \gamma)^2 }.\cr}
   \right\}
   \eqno {\rm (M^3.)}$$
By this comparison we are brought back to the general integral
equations of the relative motion of a binary system, (89.) and
(90.); but we have now the following particular values for the
coefficients $A$, $B$, $C$:
$$A = {1 \over r} {\delta w \over \delta \sigma}
      + {1 \over \tau} {\delta w \over \delta \tau},\quad
  B = {1 \over \tau} {\delta w \over \delta \tau},\quad
  C = {1 \over r_0} {\delta w \over \delta \sigma}
      + {1 \over \tau} {\delta w \over \delta \tau};
   \eqno {\rm (N^3.)}$$
and with respect to the three partial differential coefficients,
$\displaystyle {\delta w \over \delta \sigma}$,
$\displaystyle {\delta w \over \delta \tau}$,
$\displaystyle {\delta w \over \delta {\rm a}}$,
we have the following relation between them:
$${\rm a} {\delta w \over \delta {\rm a}}
   + \sigma {\delta w \over \delta \sigma}
   + \tau {\delta w \over \delta \tau}
   = {w \over 2},
   \eqno {\rm (O^3.)}$$
the function $w$ being homogeneous of the dimension ${1 \over 2}$
with respect to the three quantities ${\rm a}$, $\sigma$,
$\tau$; we have also, by (I${}^3$.),
$${\delta w \over \delta \sigma}
   = \sqrt{ \mu \over {\rm a} } \mathbin{.} {\sin \upsilon_\prime
      \over e_\prime - \cos \upsilon_\prime},\quad
  {\delta w \over \delta \tau}
   = \sqrt{ \mu \over {\rm a} } \mathbin{.} {\sqrt{ 1 - e_\prime^2 }
      \over \cos \upsilon_\prime - e_\prime},
   \eqno {\rm (P^3.)}$$
and therefore
$${\delta w \over \delta \sigma} {\delta w \over \delta \tau}
    = {-2 \mu \tau \over \sigma^2 - \tau^2},\quad
  \left( {\delta w \over \delta \sigma} \right)^2
    + \left( {\delta w \over \delta \tau} \right)^2
    + {\mu \over {\rm a}}
   = {4 \mu \sigma \over \sigma^2 - \tau^2},
   \eqno {\rm (Q^3.)}$$
from which may be deduced the following remarkable expressions:
$$\left. \eqalign{
\left( {\delta w \over \delta \sigma}
       + {\delta w \over \delta \tau} \right)^2
   &= {4 \mu \over \sigma + \tau} - {\mu \over {\rm a}},\cr
\left( {\delta w \over \delta \tau}
       - {\delta w \over \delta \sigma} \right)^2
   &= {4 \mu \over \sigma - \tau} - {\mu \over {\rm a}}.\cr}
   \right\}
   \eqno {\rm (R^3.)}$$
These expressions will be found to be important in the
application of the present method to the theory of elliptic
motion.

\bigbreak

16.
We shall not enter, on this occasion, into any details of such
application; but we may remark, that the circumstance of the
characteristic function involving only the elliptic chord and the
sum of the extreme radii, (besides the mean distance and the sum
of the masses,) affords, by our general method, a new proof of
the well-known theorem that the elliptic time also depends on the
same chord and sum of radii; and gives a new expression for the
law of this dependence, namely,
$$t = {2 {\rm a}^2 \over \mu} {\delta w \over \delta {\rm a}}.
   \eqno {\rm (S^3.)}$$
We may remark also, that the same form of the characteristic
function of elliptic motion conducts, by our general method, to
the following curious, but not novel property, of the ellipse,
that if any two tangents be drawn to such a curve, from any
common point outside, these tangents subtend equal angles at one
focus; they subtend also equal angles at the other.
Reciprocally, if any plane curve possess this property, when
referred to a fixed point in its own plane, which may be taken
as the origin of polar coordinates $r$, $\theta$, the curve must
satisfy the following equation in mixed differences:
$$\mathop{\rm cotan} \left( {\Delta\theta \over 2} \right)
      \mathbin{.} \Delta {1 \over r}
   = (\Delta + 2) {d \over d\theta} {1 \over r},
   \eqno {\rm (115.)}$$
which may be brought to the following form,
$$\left( {d \over d\theta} + {d^3 \over d\theta^3} \right)
      {1 \over r} = 0,
   \eqno {\rm (116.)}$$
and therefore gives, by integration,
$$r = {p \over 1 + e \cos (\theta - \varpi) };
   \eqno {\rm (117.)}$$
the curve is, consequently, a conic section, and the fixed point
is one of its foci.

The properties of parabolic are included as limiting cases in
those of elliptic motion, and may be deduced from them by making
$$H_\prime = 0,\quad \hbox{or} \quad {\rm a} = \infty;
   \eqno {\rm (118.)}$$
and therefore the characteristic function~$w$ and the time~$t$,
in parabolic as well as in elliptic motion, are functions of the
chord and of the sum of the radii.  By thus making ${\rm a}$
infinite in the foregoing expressions, we find, for parabolic
motion, the partial differential equations
$$\left( {\delta w \over \delta \sigma}
       + {\delta w \over \delta \tau} \right)^2
   = {4 \mu \over \sigma + \tau},\quad
  \left( {\delta w \over \delta \sigma}
       - {\delta w \over \delta \tau} \right)^2
   = {4 \mu \over \sigma - \tau};
   \eqno {\rm (T^3.)}$$
and in fact the parabolic form of the simplified characteristic
function~$w$ may easily be shown to be
$$w = 2 \sqrt{\mu} ( \sqrt{\sigma + \tau}
         \mp \sqrt{\sigma - \tau}),
   \eqno {\rm (U^3.)}$$
$\tau$ being, as before, the chord, and $\sigma$ the sum of the
radii; while the analogous limit of the expression (S${}^3$.), for
the time, is
$$t = {1 \over 6 \sqrt{\mu}} \{ (\sigma + \tau)^{3 \over 2}
         \mp (\sigma - \tau)^{3 \over 2} \}:
   \eqno {\rm (V^3.)}$$
which latter is a known expression.

The formul{\ae} (K${}^3.$) and (L${}^3$.), to the comparison of
which we have reduced the study of elliptic motion, extend to
hyperbolic motion also; and in any binary system, with
{\sc Newton}'s law of attraction, the simplified characteristic
function~$w$ may be expressed by the definite integral
$$w= \int_{-\tau}^\tau \sqrt{ {\mu \over \sigma + \tau}
         - {\mu \over 4 {\rm a}} } \mathbin{.} d \tau,
   \eqno {\rm (W^3.)}$$
this function~$w$ being still connected with the relative
action~$V_\prime$ by the equation (H${}^3$.); while the time $t$,
which may always be deduced from this function, by the law of
varying action, is represented by this other connected integral,
$$t = {1 \over 4} \int_{-\tau}^\tau \left(
          {\mu \over \sigma + \tau} - {\mu \over 4 {\rm a}}
          \right)^{-{1 \over 2}} d\tau:
   \eqno {\rm (X^3.)}$$
provided that, within the extent of these integrations, the
radical does not vanish nor become infinite.  When this
condition is not satisfied, we may still express the simplified
characteristic function~$w$, and the time~$t$, by the following
analogous integrals:
$$w = \int_{\tau_\prime}^{\sigma_\prime}
         \pm \sqrt{ {2 \mu \over \sigma_\prime}
            - {\mu \over {\rm a}} } \, d \sigma_\prime,
   \eqno {\rm (Y^3.)}$$
and
$$t = \int_{\tau_\prime}^{\sigma_\prime}
         \pm \left( {2 \mu \over \sigma_\prime}
            - {\mu \over {\rm a}} \right)^{-{1 \over 2}}
         \, d \sigma_\prime,
   \eqno {\rm (Z^3.)}$$
in which we have put for abridgement
$$\sigma_\prime = {\sigma + \tau \over 2},\quad
  \tau_\prime   = {\sigma - \tau \over 2},
   \eqno {\rm (119.)}$$
and in which it is easy to determine the signs of the radicals.
But to treat fully of these various transformations would carry us
too far at present, for it is time to consider the properties of
systems with more points than two.

\bigbreak

{\sectiontitle
On Systems of three Points, in general; and on their
Characteristic Functions.\par}

\nobreak\bigskip

17.
For any system of three points, the known differential equations
of motion of the 2nd order are included in the following formula:
$$\left. \eqalign{
  &  m_1 (x_1'' \, \delta x_1 + y_1'' \, \delta y_1
      + z_1'' \, \delta z_1)
   + m_2 (x_2'' \, \delta x_2 + y_2'' \, \delta y_2
      + z_2'' \, \delta z_2) \cr
  &+ m_3 (x_3'' \, \delta x_3 + y_3'' \, \delta y_3
      + z_3'' \, \delta z_3)
    = \delta U,\cr}
   \right\}
   \eqno {\rm (120.)}$$
the known force-function $U$ having the form
$$U = m_1 m_2 f^{(1,2)} + m_1 m_3 f^{(1,3)} + m_2 m_3 f^{(2,3)},
   \eqno {\rm (121.)}$$
in which $f^{(1,2)}$, $f^{(1,3)}$, $f^{(2,3)}$, are functions
respectively of the three following mutual distances of the
points of the system:
$$\left. \eqalign{
r^{(1,2)} &= \sqrt{(x_1 - x_2)^2 + (y_1 - y_2)^2 - (z_1 - z_2)^2},\cr
r^{(1,3)} &= \sqrt{(x_1 - x_3)^2 + (y_1 - y_3)^2 - (z_1 - z_3)^2},\cr
r^{(2,3)} &= \sqrt{(x_2 - x_3)^2 + (y_2 - y_3)^2 - (z_2 - z_3)^2}:\cr}
   \right\}
   \eqno {\rm (122.)}$$
the known differential equations of motion are therefore,
separately, for the point $m_1$,
$$\left. \eqalign{
x_1'' &=   m_2 {\delta f^{(1,2)} \over \delta x_1}
         + m_3 {\delta f^{(1,3)} \over \delta x_1},\cr
y_1'' &=   m_2 {\delta f^{(1,2)} \over \delta y_1}
         + m_3 {\delta f^{(1,3)} \over \delta y_1},\cr
z_1'' &=   m_2 {\delta f^{(1,2)} \over \delta z_1}
         + m_3 {\delta f^{(1,3)} \over \delta z_1},\cr}
   \right\}
   \eqno {\rm (123.)}$$
with six other analogous equations for the points $m_2$ and
$m_3$: $x_1''$, \&c., denoting the component accelerations of the
three points $m_1$~$m_2$~$m_3$, or the second differential
coefficients of their coordinates, taken with respect to the
time.  To integrate these equation is to assign, by their means,
nine relations between the time~$t$, the three masses $m_1$
$m_2$~$m_3$, the nine varying coordinates
$x_1$~$y_1$~$z_1$ $x_2$~$y_2$~$z_2$ $x_3$~$y_3$~$z_3$,
and their nine initial values and nine initial rates of increase,
which may be thus denoted,
$a_1$~$b_1$~$c_1$ $a_2$~$b_2$~$c_2$ $a_3$~$b_3$~$c_3$
$a_1'$~$b_1'$~$c_1'$ $a_2'$~$b_2'$~$c_2'$ $a_3'$~$b_3'$~$c_3'$.
The known intermediate integral containing the law of living
force, namely,
$$\left. \eqalign{
   &  {\textstyle {1 \over 2}} m_1 (x_1'^2 + y_1'^2 + z_1'^2)
    + {\textstyle {1 \over 2}} m_2 (x_2'^2 + y_2'^2 + z_2'^2)
    + {\textstyle {1 \over 2}} m_3 (x_3'^2 + y_3'^2 + z_3'^2) \cr
   &\qquad\qquad
  = m_1 m_2 f^{(1,2)} + m_1 m_3 f^{(1,3)} + m_2 m_3 f^{(2,3)}
      + H,\cr}
   \right\}
   \eqno {\rm (124.)}$$
gives the following initial relation:
$$\left. \eqalign{
   &  {\textstyle {1 \over 2}} m_1 (a_1'^2 + b_1'^2 + c_1'^2)
    + {\textstyle {1 \over 2}} m_2 (a_2'^2 + b_2'^2 + c_2'^2)
    + {\textstyle {1 \over 2}} m_3 (a_3'^2 + b_3'^2 + c_3'^2) \cr
   &\qquad\qquad
  = m_1 m_2 f_0^{(1,2)} + m_1 m_3 f_0^{(1,3)} + m_2 m_3 f_0^{(2,3)}
      + H,\cr}
   \right\}
   \eqno {\rm (125.)}$$
in which $f_0^{(1,2)}$, $f_0^{(1,3)}$, $f_0^{(2,3)}$ are composed
of the initial coordinates, in the same manner as
$f^{(1,2)}$~$f^{(1,3)}$~$f^{(2,3)}$ are composed of the final
coordinates.  If then we knew the nine final integrals of the
equations of motion of this ternary system, and combined them
with the initial form (125.) of the law of living force, we
should have ten relations to determine the ten quantities
$t$~$a_1'$~$b_1'$~$c_1'$ $a_2'$~$b_2'$~$c_2'$ $a_3'$~$b_3'$~$c_3'$,
namely, the time and the nine initial components of the
velocities of the three points, as functions of the nine final
and nine initial coordinates, and of the quantity $H$, involving
also the masses; we could therefore determine whatever else
depends on the manner and time of motion of the system, from its
initial to its final position, as a function of the same extreme
coordinates, and of $H$.  In particular, we could determine the
action~$V$, or the accumulated living force of the system, namely,
$$V =   m_1 \int_0^t ( x_1'^2 + y_1'^2 + z_1'^2 ) \, dt
      + m_2 \int_0^t ( x_2'^2 + y_2'^2 + z_2'^2 ) \, dt
      + m_3 \int_0^t ( x_3'^2 + y_3'^2 + z_3'^2 ) \, dt,
   \eqno {\rm (A^4.)}$$
as a function of these nineteen quantities,
$x_1$~$y_1$~$z_1$ $x_2$~$y_2$~$z_2$ $x_3$~$y_3$~$z_3$
$a_1$~$b_1$~$c_1$ $a_2$~$b_2$~$c_2$ $a_3$~$b_3$~$c_3$~$H$;
and might then calculate the variation of this function,
$$\left. \eqalign{
\delta V
   &=   {\delta V \over \delta x_1} \, \delta x_1
      + {\delta V \over \delta y_1} \, \delta y_1
      + {\delta V \over \delta z_1} \, \delta z_1
      + {\delta V \over \delta a_1} \, \delta a_1
      + {\delta V \over \delta b_1} \, \delta b_1
      + {\delta V \over \delta c_1} \, \delta c_1 \cr
   &  + {\delta V \over \delta x_2} \, \delta x_2
      + {\delta V \over \delta y_2} \, \delta y_2
      + {\delta V \over \delta z_2} \, \delta z_2
      + {\delta V \over \delta a_2} \, \delta a_2
      + {\delta V \over \delta b_2} \, \delta b_2
      + {\delta V \over \delta c_2} \, \delta c_2 \cr
   &  + {\delta V \over \delta x_3} \, \delta x_3
      + {\delta V \over \delta y_3} \, \delta y_3
      + {\delta V \over \delta z_3} \, \delta z_3
      + {\delta V \over \delta a_3} \, \delta a_3
      + {\delta V \over \delta b_3} \, \delta b_3
      + {\delta V \over \delta c_3} \, \delta c_3 \cr
   &  + {\delta V \over \delta H} \, \delta H.\cr}
   \right\}
   \eqno {\rm (B^4.)}$$
But the law of varying action gives, {\it previously}, the
following expression for this variation:
$$\left. \eqalign{
\delta V
   &=   m_1 (   x_1' \, \delta x_1 - a_1' \, \delta a_1
              + y_1' \, \delta y_1 - b_1' \, \delta b_1
              + z_1' \, \delta z_1 - c_1' \, \delta c_1 ) \cr
   &  + m_2 (   x_2' \, \delta x_2 - a_2' \, \delta a_2
              + y_2' \, \delta y_2 - b_2' \, \delta b_2
              + z_2' \, \delta z_2 - c_2' \, \delta c_2 ) \cr
   &  + m_3 (   x_3' \, \delta x_3 - a_3' \, \delta a_3
              + y_3' \, \delta y_3 - b_3' \, \delta b_3
              + z_3' \, \delta z_3 - c_3' \, \delta c_3 ) \cr
   &  + t \, \delta H;\cr}
   \right\}
   \eqno {\rm (C^4.)}$$
and shows, therefore, that the research of all the intermediate
and all the final integral equations, of motion of the system,
may be reduced, reciprocally, to the search and differentiation
of this one characteristic function~$V$; because if we knew this
one function, we should have the nine intermediate integrals of
the known differential equations, under the forms
$$\left. \multieqalign{
{\delta V \over \delta x_1} &= m_1 x_1', &
{\delta V \over \delta y_1} &= m_1 y_1', &
{\delta V \over \delta z_1} &= m_1 z_1', \cr
{\delta V \over \delta x_2} &= m_2 x_2', &
{\delta V \over \delta y_2} &= m_2 y_2', &
{\delta V \over \delta z_2} &= m_2 z_2', \cr
{\delta V \over \delta x_3} &= m_3 x_3', &
{\delta V \over \delta y_3} &= m_3 y_3', &
{\delta V \over \delta z_3} &= m_3 z_3', \cr}
   \right\}
   \eqno {\rm (D^4.)}$$
and the nine final integrals under the forms
$$\left. \multieqalign{
{\delta V \over \delta a_1} &= - m_1 a_1', &
{\delta V \over \delta b_1} &= - m_1 b_1', &
{\delta V \over \delta c_1} &= - m_1 c_1', \cr
{\delta V \over \delta a_2} &= - m_2 a_2', &
{\delta V \over \delta b_2} &= - m_2 b_2', &
{\delta V \over \delta c_2} &= - m_2 c_2', \cr
{\delta V \over \delta a_3} &= - m_3 a_3', &
{\delta V \over \delta b_3} &= - m_3 b_3', &
{\delta V \over \delta c_3} &= - m_3 c_3', \cr}
   \right\}
   \eqno {\rm (E^4.)}$$
the auxiliary constant~$H$ being to be eliminated, and the
time~$t$ introduced, by this other equation, which has often
occurred in this essay,
$$t ={\delta V \over \delta H}.
   \eqno {\rm (E.)}$$

The same law of varying action suggests also a method of
investigating the form of this characteristic function~$V$, not
requiring the previous integration of the known equations of
motion; namely, the integration of a pair of partial differential
equations connected with the law of living force; which are,
$$\left. \eqalign{
   &{1 \over 2 m_1} \left\{
      \left( {\delta V \over \delta x_1} \right)^2
    + \left( {\delta V \over \delta y_1} \right)^2
    + \left( {\delta V \over \delta z_1} \right)^2
   \right\}
   + {1 \over 2 m_2} \left\{
      \left( {\delta V \over \delta x_2} \right)^2
    + \left( {\delta V \over \delta y_2} \right)^2
    + \left( {\delta V \over \delta z_2} \right)^2
   \right\} \cr
   &\quad\quad
    + {1 \over 2 m_3} \left\{
      \left( {\delta V \over \delta x_3} \right)^2
    + \left( {\delta V \over \delta y_3} \right)^2
    + \left( {\delta V \over \delta z_3} \right)^2
   \right\} \cr
   &\qquad\qquad
   = m_1 m_2 f^{(1,2)} + m_1 m_3 f^{(1,3)} + m_2 m_3 f^{(2,3)}
    + H,\cr}
   \right\}
   \eqno {\rm (F^4.)}$$
and
$$\left. \eqalign{
   &{1 \over 2 m_1} \left\{
      \left( {\delta V \over \delta a_1} \right)^2
    + \left( {\delta V \over \delta b_1} \right)^2
    + \left( {\delta V \over \delta c_1} \right)^2
   \right\}
   + {1 \over 2 m_2} \left\{
      \left( {\delta V \over \delta a_2} \right)^2
    + \left( {\delta V \over \delta b_2} \right)^2
    + \left( {\delta V \over \delta c_2} \right)^2
   \right\} \cr
   &\quad\quad
    + {1 \over 2 m_3} \left\{
      \left( {\delta V \over \delta a_3} \right)^2
    + \left( {\delta V \over \delta b_3} \right)^2
    + \left( {\delta V \over \delta c_3} \right)^2
   \right\} \cr
   &\qquad\qquad
   = m_1 m_2 f_0^{(1,2)} + m_1 m_3 f_0^{(1,3)} + m_2 m_3 f_0^{(2,3)}
    + H.\cr}
   \right\}
   \eqno {\rm (G^4.)}$$
And to diminish the difficulty of thus determining the
function~$V$, which depends on 18 coordinates, we may separate
it, by principles already explained, into a part
$V_{\prime\prime}$ depending only on the motion of the centre
of gravity of the system, and determined by the formula (H${}^1$.),
and another part $V_\prime$, depending only on the relative motions
of the points of the system about this internal centre, and equal
to the accumulated living force, connected with this relative motion
only.  In this manner the difficulty is reduced to determining
the relative action $V_\prime$; and if we introduce the relative
coordinates
$$\left. \multieqalign{
\xi_1    &= x_1 - x_3, &
\eta_1   &= y_1 - y_3, &
\zeta_1  &= z_1 - z_3, \cr
\xi_2    &= x_2 - x_3, &
\eta_2   &= y_2 - y_3, &
\zeta_2  &= z_2 - z_3, \cr}
   \right\}
   \eqno {\rm (126.)}$$
and
$$\left. \multieqalign{
\alpha_1 &= a_1 - a_3, &
\beta_1  &= b_1 - b_3, &
\gamma_1 &= c_1 - c_3, \cr
\alpha_2 &= a_2 - a_3, &
\beta_2  &= b_2 - b_3, &
\gamma_2 &= c_2 - c_3, \cr}
   \right\}
   \eqno {\rm (127.)}$$
we easily find, by the principles of the tenth and following
numbers, that the function $V_\prime$ may be considered as
depending only on these relative coordinates and on a
quantity~$H_\prime$ analogous to $H$ (besides the masses of the
system); and that it must satisfy two partial differential
equations, analogous to (F${}^4$.) and (G${}^4$.), namely
$$\left. \eqalign{
   &{1 \over 2 m_1} \left\{
      \left( {\delta V_\prime \over \delta \xi_1} \right)^2
    + \left( {\delta V_\prime \over \delta \eta_1} \right)^2
    + \left( {\delta V_\prime \over \delta \zeta_1} \right)^2
   \right\}
   + {1 \over 2 m_2} \left\{
      \left( {\delta V_\prime \over \delta \xi_2} \right)^2
    + \left( {\delta V_\prime \over \delta \eta_2} \right)^2
    + \left( {\delta V_\prime \over \delta \zeta_2} \right)^2
   \right\} \cr
   &\qquad\qquad\mathrel{\phantom{=}} \mathord{}
   + {1 \over 2 m_3} \left\{
      \left( {\delta V_\prime \over \delta \xi_1}
           + {\delta V_\prime \over \delta \xi_2} \right)^2
    + \left( {\delta V_\prime \over \delta \eta_1}
           + {\delta V_\prime \over \delta \eta_2} \right)^2
    + \left( {\delta V_\prime \over \delta \zeta_1}
           + {\delta V_\prime \over \delta \zeta_2} \right)^2
   \right\} \cr
   &\qquad\qquad
   = m_1 m_2 f^{(1,2)} + m_1 m_3 f^{(1,3)} + m_2 m_3 f^{(2,3)}
    + H_\prime;\cr}
   \right\}
   \eqno {\rm (H^4.)}$$
and
$$\left. \eqalign{
   &{1 \over 2 m_1} \left\{
      \left( {\delta V_\prime \over \delta \alpha_1} \right)^2
    + \left( {\delta V_\prime \over \delta \beta_1} \right)^2
    + \left( {\delta V_\prime \over \delta \gamma_1} \right)^2
   \right\}
   + {1 \over 2 m_2} \left\{
      \left( {\delta V_\prime \over \delta \alpha_2} \right)^2
    + \left( {\delta V_\prime \over \delta \beta_2} \right)^2
    + \left( {\delta V_\prime \over \delta \gamma_2} \right)^2
   \right\} \cr
   &\qquad\qquad\mathrel{\phantom{=}} \mathord{}
   + {1 \over 2 m_3} \left\{
      \left( {\delta V_\prime \over \delta \alpha_1}
           + {\delta V_\prime \over \delta \alpha_2} \right)^2
    + \left( {\delta V_\prime \over \delta \beta_1}
           + {\delta V_\prime \over \delta \beta_2} \right)^2
    + \left( {\delta V_\prime \over \delta \gamma_1}
           + {\delta V_\prime \over \delta \gamma_2} \right)^2
   \right\} \cr
   &\qquad\qquad
   = m_1 m_2 f_0^{(1,2)} + m_1 m_3 f_0^{(1,3)} + m_2 m_3 f_0^{(2,3)}
    + H_\prime:\cr}
   \right\}
   \eqno {\rm (I^4.)}$$
the law of the variation of this function being, by (Z${}^1$.),
$$\left. \eqalign{
\delta V_\prime
   &= t \, \delta H_\prime + m_1 (
         \xi_1' \, \delta \xi_1 - \alpha_1' \, \delta \alpha_1
       + \eta_1' \, \delta \eta_1 - \beta_1' \, \delta \beta_1
       + \zeta_1' \, \delta \zeta_1 - \gamma_1' \, \delta \gamma_1 )
   \cr
   &\mathrel{\phantom{= t \, \Delta H_\prime}} \mathord{}
   + m_2 (
         \xi_2' \, \delta \xi_2 - \alpha_2' \, \delta \alpha_2
       + \eta_2' \, \delta \eta_2 - \beta_2' \, \delta \beta_2
       + \zeta_2' \, \delta \zeta_2 - \gamma_2' \, \delta \gamma_2 )
   \cr
   &\mathrel{\phantom{=}} \mathord{}
   - {1 \over m_1 + m_2 + m_3} \left\{
   \eqalign{
         (m_1 \xi_1' + m_2 \xi_2')
            (m_1 \, \delta \xi_1 + m_2 \, \delta \xi_2) \cr
       - (m_1 \alpha_1' + m_2 \alpha_2')
            (m_1 \, \delta \alpha_1 + m_2 \, \delta \alpha_2) \cr
       + (m_1 \eta_1' + m_2 \eta_2')
            (m_1 \, \delta \eta_1 + m_2 \, \delta \eta_2) \cr
       - (m_1 \beta_1' + m_2 \beta_2')
            (m_1 \, \delta \beta_1 + m_2 \, \delta \beta_2) \cr
       + (m_1 \zeta_1' + m_2 \zeta_2')
            (m_1 \, \delta \zeta_1 + m_2 \, \delta \zeta_2) \cr
       - (m_1 \gamma_1' + m_2 \gamma_2')
            (m_1 \, \delta \gamma_1 + m_2 \, \delta \gamma_2) \cr}
      \right\} \cr}
   \right\}
   \eqno {\rm (K^4.)}$$
which resolves itself in the same manner as before into the six
intermediate and six final integrals of relative motion,
namely, into the following equations:
$$\left. \multieqalign{
{1 \over m_1} {\delta V_\prime \over \delta \xi_1}
   &= \xi_1'
      - {m_1 \xi_1' + m_2 \xi_2' \over m_1 + m_2 + m_3}; &
{1 \over m_2} {\delta V_\prime \over \delta \xi_2}
   &= \xi_2'
      - {m_1 \xi_1' + m_2 \xi_2' \over m_1 + m_2 + m_3}; \cr
{1 \over m_1} {\delta V_\prime \over \delta \eta_1}
   &= \eta_1'
      - {m_1 \eta_1' + m_2 \eta_2' \over m_1 + m_2 + m_3}; &
{1 \over m_2} {\delta V_\prime \over \delta \eta_2}
   &= \eta_2'
      - {m_1 \eta_1' + m_2 \eta_2' \over m_1 + m_2 + m_3}; \cr
{1 \over m_1} {\delta V_\prime \over \delta \zeta_1}
   &= \zeta_1'
      - {m_1 \zeta_1' + m_2 \zeta_2' \over m_1 + m_2 + m_3}; &
{1 \over m_2} {\delta V_\prime \over \delta \zeta_2}
   &= \zeta_2'
      - {m_1 \zeta_1' + m_2 \zeta_2' \over m_1 + m_2 + m_3}; \cr}
   \right\}
   \eqno {\rm (L^4.)}$$
and
$$\left. \multieqalign{
{-1 \over m_1} {\delta V_\prime \over \delta \alpha_1}
   &= \alpha_1'
      - {m_1 \alpha_1' + m_2 \alpha_2' \over m_1 + m_2 + m_3}; &
{-1 \over m_2} {\delta V_\prime \over \delta \alpha_2}
   &= \alpha_2'
      - {m_1 \alpha_1' + m_2 \alpha_2' \over m_1 + m_2 + m_3}; \cr
{-1 \over m_1} {\delta V_\prime \over \delta \beta_1}
   &= \beta_1'
      - {m_1 \beta_1' + m_2 \beta_2' \over m_1 + m_2 + m_3}; &
{-1 \over m_2} {\delta V_\prime \over \delta \beta_2}
   &= \beta_2'
      - {m_1 \beta_1' + m_2 \beta_2' \over m_1 + m_2 + m_3}; \cr
{-1 \over m_1} {\delta V_\prime \over \delta \gamma_1}
   &= \gamma_1'
      - {m_1 \gamma_1' + m_2 \gamma_2' \over m_1 + m_2 + m_3}; &
{-1 \over m_2} {\delta V_\prime \over \delta \gamma_2}
   &= \gamma_2'
      - {m_1 \gamma_1' + m_2 \gamma_2' \over m_1 + m_2 + m_3}; \cr}
   \right\}
   \eqno {\rm (M^4.)}$$
which must be combined with our old formula,
$${\delta V_\prime \over \delta H_\prime} = t.
   \eqno {\rm (O^1.)}$$

\bigbreak

18.
The quantity $H_\prime$ in $V_\prime$, and the analogous quantity
$H_{\prime\prime}$ in $V_{\prime\prime}$, are indeed independent
of the time, and do not vary in the course of the motion; but it
is required by the spirit of our method, that in deducing the
absolute action or original characteristic function $V$ from the
two parts $V_\prime$ and $V_{\prime\prime}$, we should consider
these two parts $H_\prime$ and $H_{\prime\prime}$ of the original
quantity $H_\prime$ as functions involving each the nine initial
and nine final coordinates of the points of the ternary system;
the forms of these two functions, of the eighteen coordinates and
of $H$, being determined by the two conditions,
$${\delta V_\prime \over \delta H_\prime}
     = {\delta V_{\prime\prime} \over \delta H_{\prime\prime}},\quad
  H_\prime + H_{\prime\prime} = H.
   \eqno {\rm (N^4.)}$$
However it results from these conditions, that in taking the
variation of the whole original function~$V$, of the first
order, with respect to the eighteen coordinates, we may treat the
two auxiliary quantities $H_\prime$ and $H_{\prime\prime}$ as
constant; and therefore that we have the following expressions
for the partial differential coefficients of the first order of
$V$, taken with respect to the coordinates parallel to $x$,
$$\left. \multieqalign{
{\delta V \over \delta x_1}
   &= {\delta V_\prime \over \delta \xi_1}
       + {m_1 \over m_1 + m_2 + m_3}
         {\delta V_{\prime\prime} \over \delta x_{\prime\prime}},&
{\delta V \over \delta a_1}
   &= {\delta V_\prime \over \delta \alpha_1}
       + {m_1 \over m_1 + m_2 + m_3}
         {\delta V_{\prime\prime} \over \delta a_{\prime\prime}},\cr
{\delta V \over \delta x_2}
   &= {\delta V_\prime \over \delta \xi_2}
       + {m_2 \over m_1 + m_2 + m_3}
         {\delta V_{\prime\prime} \over \delta x_{\prime\prime}},&
{\delta V \over \delta a_2}
   &= {\delta V_\prime \over \delta \alpha_2}
       + {m_2 \over m_1 + m_2 + m_3}
         {\delta V_{\prime\prime} \over \delta a_{\prime\prime}},\cr
{\delta V \over \delta x_3}
   &=  - {\delta V_\prime \over \delta \xi_1}
       - {\delta V_\prime \over \delta \xi_2}
       + {m_3 \over m_1 + m_2 + m_3}
         {\delta V_{\prime\prime} \over \delta x_{\prime\prime}},&
{\delta V \over \delta a_3}
   &=  - {\delta V_\prime \over \delta \alpha_1}
       - {\delta V_\prime \over \delta \alpha_2}
       + {m_3 \over m_1 + m_2 + m_3}
         {\delta V_{\prime\prime} \over \delta a_{\prime\prime}},\cr}
   \right\}
   \eqno {\rm (O^4.)}$$
together with analogous expressions for the partial differential
coefficients of the same order taken with respect to the other
coordinates.  Substituting these expressions in the equations of
the form (O.), namely, in the following,
$$\left. \eqalign{
        {\delta V \over \delta x_1}
      + {\delta V \over \delta x_2}
      + {\delta V \over \delta x_3}
      + {\delta V \over \delta a_1}
      + {\delta V \over \delta a_2}
      + {\delta V \over \delta a_3}
   &= 0,\cr
        {\delta V \over \delta y_1}
      + {\delta V \over \delta y_2}
      + {\delta V \over \delta y_3}
      + {\delta V \over \delta b_1}
      + {\delta V \over \delta b_2}
      + {\delta V \over \delta b_3}
   &= 0,\cr
        {\delta V \over \delta z_1}
      + {\delta V \over \delta z_2}
      + {\delta V \over \delta z_3}
      + {\delta V \over \delta c_1}
      + {\delta V \over \delta c_2}
      + {\delta V \over \delta c_3}
   &= 0,\cr}
   \right\}
   \eqno {\rm (P^4.)}$$
we find that these equations become identical, because
$$      {\delta V_{\prime\prime} \over \delta x_{\prime\prime}}
      + {\delta V_{\prime\prime} \over \delta a_{\prime\prime}}
   = 0,\quad
        {\delta V_{\prime\prime} \over \delta y_{\prime\prime}}
      + {\delta V_{\prime\prime} \over \delta b_{\prime\prime}}
   = 0,\quad
        {\delta V_{\prime\prime} \over \delta z_{\prime\prime}}
      + {\delta V_{\prime\prime} \over \delta c_{\prime\prime}}
   = 0,
   \eqno {\rm (Q^4.)}$$
But substituting, in like manner, the expressions (O${}^4$.) in
the equations of the form (P.), of which the first is, for a
ternary system,
$$\left. \eqalign{
   &\mathbin{\phantom{+}}
      x_1 {\delta V \over \delta y_1}
    - y_1 {\delta V \over \delta x_1}
    + x_2 {\delta V \over \delta y_2}
    - y_2 {\delta V \over \delta x_2}
    + x_3 {\delta V \over \delta y_3}
    - y_3 {\delta V \over \delta x_3} \cr
   &
    + a_1 {\delta V \over \delta b_1}
    - b_1 {\delta V \over \delta a_1}
    + a_2 {\delta V \over \delta b_2}
    - b_2 {\delta V \over \delta a_2}
    + a_3 {\delta V \over \delta b_3}
    - b_3 {\delta V \over \delta a_3};\cr}
   \right\}
   \eqno {\rm (R^4.)}$$
and observing that we have
$$x_{\prime\prime}
      {\delta V_{\prime\prime} \over \delta y_{\prime\prime}}
  - y_{\prime\prime}
      {\delta V_{\prime\prime} \over \delta x_{\prime\prime}}
  + a_{\prime\prime}
      {\delta V_{\prime\prime} \over \delta b_{\prime\prime}}
  - b_{\prime\prime}
      {\delta V_{\prime\prime} \over \delta a_{\prime\prime}}
   = 0,
   \eqno {\rm (S^4.)}$$
along with two other analogous conditions, we find that the part
$V_\prime$, or the characteristic function of relative motion of
the ternary system, must satisfy the three following conditions,
involving its partial differential coefficients of the first order
and in the first degree,
$$\left. \eqalign{
0 &=   \xi_1  {\delta V_\prime \over \delta \eta_1}
     - \eta_1 {\delta V_\prime \over \delta \xi_1}
     + \xi_2  {\delta V_\prime \over \delta \eta_2}
     - \eta_2 {\delta V_\prime \over \delta \xi_2}
     + \alpha_1  {\delta V_\prime \over \delta \beta_1}
     - \beta_1 {\delta V_\prime \over \delta \alpha_1}
     + \alpha_2  {\delta V_\prime \over \delta \beta_2}
     - \beta_2 {\delta V_\prime \over \delta \alpha_2},\cr
0 &=   \eta_1  {\delta V_\prime \over \delta \zeta_1}
     - \zeta_1 {\delta V_\prime \over \delta \eta_1}
     + \eta_2  {\delta V_\prime \over \delta \zeta_2}
     - \zeta_2 {\delta V_\prime \over \delta \eta_2}
     + \beta_1  {\delta V_\prime \over \delta \gamma_1}
     - \gamma_1 {\delta V_\prime \over \delta \beta_1}
     + \beta_2  {\delta V_\prime \over \delta \gamma_2}
     - \gamma_2 {\delta V_\prime \over \delta \beta_2},\cr
0 &=   \zeta_1  {\delta V_\prime \over \delta \xi_1}
     - \xi_1 {\delta V_\prime \over \delta \zeta_1}
     + \zeta_2  {\delta V_\prime \over \delta \xi_2}
     - \xi_2 {\delta V_\prime \over \delta \zeta_2}
     + \gamma_1  {\delta V_\prime \over \delta \alpha_1}
     - \alpha_1 {\delta V_\prime \over \delta \gamma_1}
     + \gamma_2  {\delta V_\prime \over \delta \alpha_2}
     - \alpha_2 {\delta V_\prime \over \delta \gamma_2},\cr}
   \right\}
   \eqno {\rm (T^4.)}$$
which show that this function can depend only on the shape and
size of a pentagon, not generally plane, formed by the point $m_3$
considered as fixed, and by the initial and final positions of
the other two points $m_1$ and $m_2$; for example, the pentagon,
of which the corners are, in order,
$m_3$ $(m_1)$ $(m_2)$ $m_2$ $m_1$; $(m_1)$ and $(m_2)$
denoting the initial positions of the points $m_1$ and $m_2$,
referred to $m_3$ as a fixed origin.  The shape and size of this
pentagon may be determined by the ten mutual distances of its
five points, that is, by the five sides and five diagonals, which
may be thus denoted:
$$\left. \multieqalign{
 m_3  (m_1) &= \sqrt{s_1}, &
(m_1) (m_2) &= \sqrt{s_2}, &
(m_2)  m_2  &= \sqrt{s_3}, &
 m_2   m_1  &= \sqrt{s_4}, &
 m_1   m_3  &= \sqrt{s_5}, \cr
 m_3  (m_2) &= \sqrt{d_1}, &
(m_1)  m_2  &= \sqrt{d_2}, &
(m_2)  m_1  &= \sqrt{d_3}, &
 m_2   m_3  &= \sqrt{d_4}, &
 m_1  (m_1) &= \sqrt{d_5}; \cr}
   \right\}
   \eqno {\rm (128.)}$$
the values of $s_1 \,\ldots\, d_5$ as functions of the twelve
relative coordinates being
$$\left. \multieqalign{
s_1 &= \alpha_1^2 + \beta_1^2 + \gamma_1^2, &
s_2 &= (\alpha_2 - \alpha_1)^2 + (\beta_2 - \beta_1)^2
         + (\gamma_2 - \gamma_1)^2, \cr
    & &
s_3 &= (\xi_2 - \alpha_2)^2 + (\eta_2 - \beta_2)^2
         + (\zeta_2 - \gamma_2)^2, \cr
s_5 &= \xi_1^2 + \eta_1^2 + \zeta_1^2, &
s_4 &= (\xi_1 - \xi_2)^2 + (\eta_1 - \eta_2)^2
         + (\zeta_1 - \zeta_2)^2, \cr
d_1 &= \alpha_2^2 + \beta_2^2 + \gamma_2^2, &
d_2 &= (\xi_2 - \alpha_1)^2 + (\eta_2 - \beta_1)^2
         + (\zeta_2 - \gamma_1)^2, \cr
    & &
d_3 &= (\xi_1 - \alpha_2)^2 + (\eta_1 - \beta_2)^2
         + (\zeta_1 - \gamma_2)^2, \cr
d_4 &= \xi_2^2 + \eta_2^2 + \zeta_2^2, &
d_5 &= (\xi_1 - \alpha_1)^2 + (\eta_1 - \beta_1)^2
         + (\zeta_1 - \gamma_1)^2. \cr}
   \right\}
   \eqno {\rm (129.)}$$
These ten distances $\sqrt{s_1}$, \&c., are not, however, all
independent, but are connected by one equation of condition,
namely,

\vfill\eject  % Page break necessary with current page size

$$\left.
\vcenter{\halign{$#$\hfil&&\hskip2pt $#$\hfil\cr
0 =   &\mathbin{\phantom{+}}
         s_1^2 s_3^2
      &+ s_2^2 s_4^2
      &+ s_3^2 s_5^2
      &+ s_4^2 s_1^2
      &+ s_5^2 s_2^2 \cr
      &+ s_1^2 d_3^2
      &+ s_2^2 d_4^2
      &+ s_3^2 d_5^2
      &+ s_4^2 d_1^2
      &+ s_5^2 d_2^2 \cr
      &+ d_1^2 d_2^2
      &+ d_2^2 d_3^2
      &+ d_3^2 d_4^2
      &+ d_4^2 d_5^2
      &+ d_5^2 d_1^2 \cr
      &- 2 s_1^2 s_3 s_4
      &- 2 s_2^2 s_4 s_5
      &- 2 s_3^2 s_5 s_1
      &- 2 s_4^2 s_1 s_2
      &- 2 s_5^2 s_2 s_3 \cr
      &- 2 s_1^2 s_3 d_3
      &- 2 s_2^2 s_4 d_4
      &- 2 s_3^2 s_5 d_5
      &- 2 s_4^2 s_1 d_1
      &- 2 s_5^2 s_2 d_2 \cr
      &- 2 s_1^2 s_4 d_3
      &- 2 s_1^2 s_5 d_4
      &- 2 s_1^2 s_1 d_5
      &- 2 s_1^2 s_2 d_1
      &- 2 s_1^2 s_3 d_2 \cr
      &- 2 s_1 d_2 d_3^2
      &- 2 s_2 d_3 d_4^2
      &- 2 s_3 d_4 d_5^2
      &- 2 s_4 d_5 d_1^2
      &- 2 s_5 d_1 d_2^2 \cr
      &- 2 s_1 d_3^2 d_4
      &- 2 s_2 d_4^2 d_5
      &- 2 s_3 d_5^2 d_1
      &- 2 s_4 d_1^2 d_2
      &- 2 s_5 d_2^2 d_3 \cr
      &- 2 d_1 d_2^2 d_3
      &- 2 d_2 d_3^2 d_4
      &- 2 d_3 d_4^2 d_5
      &- 2 d_4 d_5^2 d_1
      &- 2 d_5 d_1^2 d_2 \cr
      &- 4 s_1 s_3 s_4 d_3
      &- 4 s_2 s_4 s_5 d_4
      &- 4 s_3 s_5 s_1 d_5
      &- 4 s_4 s_1 s_2 d_1
      &- 4 s_5 s_2 s_3 d_2 \cr
      &- 4 s_1 d_2 d_3 d_4
      &- 4 s_2 d_3 d_4 d_5
      &- 4 s_3 d_4 d_5 d_1
      &- 4 s_4 d_5 d_1 d_2
      &- 4 s_5 d_1 d_2 d_3 \cr
      &- 2 s_1 s_2 s_3 d_4
      &- 2 s_2 s_3 s_4 d_5
      &- 2 s_3 s_4 s_5 d_1
      &- 2 s_4 s_5 s_1 d_2
      &- 2 s_5 s_1 s_2 d_3 \cr
      &- 2 s_1 s_3 d_1 d_2
      &- 2 s_2 s_4 d_2 d_3
      &- 2 s_3 s_5 d_3 d_4
      &- 2 s_4 s_1 d_4 d_5
      &- 2 s_5 s_2 d_5 d_1 \cr
      &- 2 s_1 d_1 d_3 d_5
      &- 2 s_2 d_2 d_4 d_1
      &- 2 s_3 d_3 d_5 d_2
      &- 2 s_4 d_4 d_1 d_3
      &- 2 s_5 d_5 d_2 d_4 \cr
      &+ 2 s_1 s_2 s_3 s_4
      &+ 2 s_2 s_3 s_4 s_5
      &+ 2 s_3 s_4 s_5 s_1
      &+ 2 s_4 s_5 s_1 s_2
      &+ 2 s_5 s_1 s_2 s_3 \cr
      &+ 2 s_1 s_2 s_4 d_3
      &+ 2 s_2 s_3 s_5 d_4
      &+ 2 s_3 s_4 s_1 d_5
      &+ 2 s_4 s_5 s_2 d_1
      &+ 2 s_5 s_1 s_3 d_2 \cr
      &+ 2 s_1 s_3 s_4 d_1
      &+ 2 s_2 s_4 s_5 d_2
      &+ 2 s_3 s_5 s_1 d_3
      &+ 2 s_4 s_1 s_2 d_4
      &+ 2 s_5 s_2 s_3 d_5 \cr
      &+ 2 s_1 s_2 d_3 d_4
      &+ 2 s_2 s_3 d_4 d_5
      &+ 2 s_3 s_4 d_5 d_1
      &+ 2 s_4 s_5 d_1 d_2
      &+ 2 s_5 s_1 d_2 d_3 \cr
      &+ 2 s_1 s_3 d_2 d_3
      &+ 2 s_2 s_4 d_3 d_4
      &+ 2 s_3 s_5 d_4 d_5
      &+ 2 s_4 s_1 d_5 d_1
      &+ 2 s_5 s_2 d_1 d_2 \cr
      &+ 2 s_1 s_4 d_1 d_2
      &+ 2 s_2 s_5 d_2 d_3
      &+ 2 s_3 s_1 d_3 d_4
      &+ 2 s_4 s_2 d_4 d_5
      &+ 2 s_5 s_3 d_5 d_1 \cr
      &+ 2 s_1 s_4 d_1 d_3
      &+ 2 s_2 s_5 d_2 d_4
      &+ 2 s_3 s_1 d_3 d_5
      &+ 2 s_4 s_2 d_4 d_1
      &+ 2 s_5 s_3 d_5 d_2 \cr
      &+ 2 s_1 s_4 d_2 d_3
      &+ 2 s_2 s_5 d_3 d_4
      &+ 2 s_3 s_1 d_4 d_5
      &+ 2 s_4 s_2 d_5 d_1
      &+ 2 s_5 s_3 d_1 d_2 \cr
      &+ 2 s_1 s_4 d_3 d_4
      &+ 2 s_2 s_5 d_4 d_5
      &+ 2 s_3 s_1 d_5 d_1
      &+ 2 s_4 s_2 d_1 d_2
      &+ 2 s_5 s_3 d_2 d_3 \cr
      &+ 2 s_1 d_1 d_2 d_3
      &+ 2 s_2 d_2 d_3 d_4
      &+ 2 s_3 d_3 d_4 d_5
      &+ 2 s_4 d_4 d_5 d_1
      &+ 2 s_5 d_5 d_1 d_2 \cr
      &+ 2 s_1 d_3 d_4 d_5
      &+ 2 s_2 d_4 d_5 d_1
      &+ 2 s_3 d_5 d_1 d_2
      &+ 2 s_4 d_1 d_2 d_3
      &+ 2 s_5 d_2 d_3 d_4 \cr
      &+ 2 d_1 d_2 d_3 d_4
      &+ 2 d_2 d_3 d_4 d_5
      &+ 2 d_3 d_4 d_5 d_1
      &+ 2 d_4 d_5 d_1 d_2
      &+ 2 d_5 d_1 d_2 d_3;\cr}}
   \right\}
   \eqno {\rm (130.)}$$
they may therefore be expressed as functions of nine independent
quantities; for example, of four lines and five angles,
$r^{(1)}$~$r_0^{(1)}$~$r^{(2)}$~$r_0^{(2)}$,
$\theta^{(1)}$~$\theta_0^{(1)}$~$\theta^{(2)}$~$\theta_0^{(2)}$~$\iota$,
on which they depend as follows:
$$\left. \eqalign{
s_1 &= r_0^{(1)2},\cr
s_2 &= r_0^{(1)2} + r_0^{(2)2}
         -2 r_0^{(1)} r_0^{(2)}
         ( \cos \theta_0^{(1)} \cos \theta_0^{(2)}
         + \sin \theta_0^{(1)} \sin \theta_0^{(2)} \cos \iota),\cr
s_3 &= r^{(2)2} + r_0^{(2)2} - 2 r^{(2)} r_0^{(2)}
         \cos (\theta^{(2)} - \theta_0^{(2)}),\cr
s_4 &= r^{(2)2} + r^{(1)2}
         - 2 r^{(2)} r^{(1)}
         ( \cos \theta^{(1)} \cos \theta^{(2)}
         + \sin \theta^{(1)} \sin \theta^{(2)} \cos \iota),\cr
s_5 & = r^{(1)2},\cr
d_1 & = r_0^{(2)2},\cr
d_2 &= r^{(2)2} + r_0^{(1)2}
         -2 r^{(2)} r_0^{(1)}
         ( \cos \theta^{(2)} \cos \theta_0^{(1)}
         + \sin \theta^{(2)} \sin \theta_0^{(1)} \cos \iota),\cr
d_3 &= r_0^{(2)2} + r^{(1)2}
         -2 r_0^{(2)} r^{(1)}
         ( \cos \theta_0^{(2)} \cos \theta^{(1)}
         + \sin \theta_0^{(2)} \sin \theta^{(1)} \cos \iota),\cr
d_4 &= r^{(2)2},\cr
d_5 &= r^{(1)2} + r_0^{(1)2} - 2 r^{(1)} r_0^{(1)}
         \cos (\theta^{(1)} - \theta_0^{(1)}),\cr}
   \right\}
   \eqno {\rm (131.)}$$
the two line-symbols $r^{(1)}$~$r^{(2)}$ denoting, for
abridgement, the same two final radii vectores which were before
denoted by $r^{(1,3)}$~$r^{(2,3)}$, and $r_0^{(1)}$~$r_0^{(2)}$
representing the initial values of these radii; while
$\theta^{(1)}$~$\theta^{(2)}$~$\theta_0^{(1)}$~$\theta_0^{(2)}$
are angles made by these four radii, with the line of
intersection of the two planes $r_0^{(1)} r^{(1)}$,
$r_0^{(2)} r^{(2)}$; and $\iota$ is the inclination of these two
planes to each other.  We may therefore consider the
characteristic function $V_\prime$ of relative motion, for any
ternary system, as depending only on these latter lines and
angles, along with the quantity $H_\prime$.

The reasoning which it has been thought useful to develope here,
for any system of three points, attracting or repelling one
another according to any functions of their distances, was
alluded to, under a more general form, in the twelth number of
this essay; and shows, for example, that the characteristic
function of relative motion in a system of four such points,
depends on the shape and size of a heptagon, and therefore only
on the mutual distances of its seven corners, which are in number
$\displaystyle \left( {7 \times 6 \over 2} = \right) \, 21$,
but are connected by six equations of condition, leaving only
fifteen independent.  It is easy to extend these remarks to any
multiple system.

\bigbreak

{\sectiontitle
General method of improving an approximate expression for the
Characteristic Function of motion of a System in any problem
of Dynamics.\par}

\nobreak\bigskip

19.
The partial differential equation (F.), which the characteristic
function~$V$ must satisfy, in every dynamical question, may
receive some useful general transformations, by the separation of
this function~$V$ into any two parts
$$V_1 + V_2 = V.
   \eqno {\rm (U^4.)}$$
For if we establish, for abridgement, the two following equations
of definition,
$$\left. \eqalign{
T_1 &= \sum \mathbin{.} {1 \over 2m} \left(
         \left( {\delta V_1 \over \delta x} \right)^2
       + \left( {\delta V_1 \over \delta y} \right)^2
       + \left( {\delta V_1 \over \delta z} \right)^2
      \right),\cr
T_2 &= \sum \mathbin{.} {1 \over 2m} \left(
         \left( {\delta V_2 \over \delta x} \right)^2
       + \left( {\delta V_2 \over \delta y} \right)^2
       + \left( {\delta V_2 \over \delta z} \right)^2
      \right),\cr}
   \right\}
   \eqno {\rm (V^4.)}$$
analogous to the relation
$$T    = \sum \mathbin{.} {1 \over 2m} \left(
         \left( {\delta V \over \delta x} \right)^2
       + \left( {\delta V \over \delta y} \right)^2
       + \left( {\delta V \over \delta z} \right)^2
      \right),
   \eqno {\rm (W^4.)}$$
which served to transform the law of living force into the
partial differential equation (F.); we shall have, by
(U${}^4$.),
$$T = T_1 + T_2 + \sum \mathbin{.} {1 \over m} \left(
         {\delta V_1 \over \delta x}
            {\delta V_2 \over \delta x}
       + {\delta V_1 \over \delta y}
            {\delta V_2 \over \delta y}
       + {\delta V_1 \over \delta z}
            {\delta V_2 \over \delta z}
      \right);
   \eqno {\rm (X^4.)}$$
and this expression may be further transformed by the help of the
formula (C.), or by the law of varying action.  For that law gives
the following symbolic equation,
$$\sum \mathbin{.} {1 \over m} \left(
         {\delta V \over \delta x} {\delta \over \delta x}
       + {\delta V \over \delta x} {\delta \over \delta y}
       + {\delta V \over \delta x} {\delta \over \delta z}
      \right)
   = {d \over dt},
   \eqno {\rm (Y^4.)}$$
the symbols in both members being prefixed to any one function of
the varying coordinates of a system, not expressly involving the
time; it gives therefore by (U${}^4$.), (V${}^4$.),
$$\sum \mathbin{.} {1 \over m} \left(
         {\delta V_1 \over \delta x}
            {\delta V_2 \over \delta x}
       + {\delta V_1 \over \delta y}
            {\delta V_2 \over \delta y}
       + {\delta V_1 \over \delta z}
            {\delta V_2 \over \delta z}
      \right)
   = {d V_2 \over dt} - 2 T_2.
   \eqno {\rm (Z^4.)}$$
In this manner we find the following general and rigorous
transformation of the equation (F.),
$${d V_2 \over dt} = T - T_1 + T_2;
   \eqno {\rm (A^5.)}$$
$T$ being here retained for the sake of symmetry and conciseness,
instead of the equal expression $U + H$.  And if we suppose, as
we may, that the part $V_1$, like the whole function $V$,
is chosen so as to vanish with the time, then the other part
$V_2$ will also have that property, and may be expressed by the
definite integral,
$$V_2 = \int_0^t (T - T_1 + T_2) \,dt.
   \eqno {\rm (B^5.)}$$

More generally, if we employ the principles of the seventh
number, and introduce any $3n$ marks
$\eta_1, \eta_2,\ldots\, \eta_{3n}$, of the varying positions of
the $n$ points of any system, (whether they be the rectangular
coordinates themselves, or any functions of them,) we shall have
$$T = F \left( {\delta V \over \delta \eta_1},
               {\delta V \over \delta \eta_2},
               \ldots\,
               {\delta V \over \delta \eta_{3n}}
               \right),
   \eqno {\rm (C^5.)}$$
and may establish by analogy the two following equations of
definition,
$$\left. \eqalign{
T_1 &= F \left( {\delta V_1 \over \delta \eta_1},
                {\delta V_1 \over \delta \eta_2},
                \ldots\,
                {\delta V_1 \over \delta \eta_{3n}}
                \right),\cr
T_2 &= F \left( {\delta V_2 \over \delta \eta_1},
                {\delta V_2 \over \delta \eta_2},
                \ldots\,
                {\delta V_2 \over \delta \eta_{3n}}
                \right),\cr}
   \right\}
   \eqno {\rm (D^5.)}$$
the function $F$ being always rational and integer, and
homogeneous of the second dimension; and being therefore such
that (besides other properties)
$$T = T_1 + T_2
          + {\delta T_1 \over
            \displaystyle \;\; \delta {\delta V_1 \over \delta \eta_1}\;\;}
               {\delta V_2 \over \delta \eta_1}
          + {\delta T_1 \;\;\over
            \displaystyle \;\; \delta {\delta V_1 \over \delta \eta_2}\;\;}
               {\delta V_2 \over \delta \eta_2}
          + \cdots
          + {\delta T_1 \over
            \displaystyle \;\; \delta {\delta V_1 \over \delta \eta_{3n}}\;\;}
               {\delta V_2 \over \delta \eta_{3n}},
   \eqno {\rm (E^5.)}$$
$${\delta T \over
      \displaystyle \;\; \delta {\delta V \over \delta \eta_1}\;\;}
   =  {\delta T_1 \over
      \displaystyle \;\; \delta {\delta V_1 \over \delta \eta_1}\;\;}
    + {\delta T_2 \over
      \displaystyle \;\; \delta {\delta V_2 \over \delta \eta_1}\;\;},
   \ldots \quad
  {\delta T \over
      \displaystyle \;\; \delta {\delta V \over \delta \eta_{3n}}\;\;}
   =  {\delta T_1 \over
      \displaystyle \;\; \delta {\delta V_1 \over \delta \eta_{3n}}\;\;}
    + {\delta T_2 \over
      \displaystyle \;\; \delta {\delta V_2 \over \delta \eta_{3n}}\;\;},
   \eqno {\rm (F^5.)}$$
and
$${\delta T_2 \over
      \displaystyle \;\; \delta {\delta V_2 \over \delta \eta_1}\;\;}
         {\delta V_2 \over \delta \eta_1}
    + {\delta T_2 \over
      \displaystyle \;\; \delta {\delta V_2 \over \delta \eta_2}\;\;}
         {\delta V_2 \over \delta \eta_2}
    + \cdots
    + {\delta T_2 \over
      \displaystyle \;\; \delta {\delta V_2 \over \delta \eta_{3n}}\;\;}
         {\delta V_2 \over \delta \eta_{3n}}
   = 2 T_2.
   \eqno {\rm (G^5.)}$$
By the principles of the eighth number, we have also,
$${\delta T \over
      \displaystyle \;\; \delta {\delta V \over \delta \eta_1}\;\;}
   = \eta_1',\quad
  {\delta T \over
      \displaystyle \;\; \delta {\delta V \over \delta \eta_2}\;\;}
   = \eta_2',\quad \ldots \quad
  {\delta T \over
      \displaystyle \;\; \delta {\delta V \over \delta \eta_{3n}}\;\;}
   = \eta_{3n}';
   \eqno {\rm (H^5.)}$$
and since the meanings of $\eta_1',\ldots\, \eta_{3n}'$ give
evidently the symbolical equation,
$$\eta_1' {\delta \over \delta \eta_1}
    + \eta_2' {\delta \over \delta \eta_2}
    + \cdots
    + \eta_{3n}' {\delta \over \delta \eta_{3n}}
   = {d \over dt},
   \eqno {\rm (I^5.)}$$
we see that the equation (A${}^5$.) still holds with the present
more general marks of position of a moving system, and gives
still the expression (B${}^5$.), supposing only, as before, that
the two parts of the whole characteristic function are chosen so
as to vanish with the time.

It may not at first sight appear, that this rigorous
transformation (B${}^5$.), of the partial differential equation
(F.), or of the analogous equation (T.) with coordinates not
rectangular, is likely to assist much in discovering the form of
the part $V_2$ of the characteristic function~$V$, (the other
part $V_1$ being supposed to have been previously assumed;)
because it involves under the sign of integration, in the term
$T_2$, the partial differential coefficients of the sought part
$V_2$.  But if we observe that these unknown coefficients enter
only by their squares and products, we shall perceive that it
offers a general method of improving an approximation in any
problem of dynamics.  For if the first part $V_1$ be an
approximate value of the whole sought function $V$, the second
part $V_2$ will be small, and the term $T_2$ will not only be
also small, but will be in general of a higher order of
smallness; we shall therefore in general improve an approximate
value $V_1$ of the characteristic function $V$, by adding to it
the definite integral,
$$V_2 = \int_0^t (T - T_1) \, dt;
   \eqno {\rm (K^5.)}$$
though this is not, like (B${}^5$.), a perfectly rigorous
expression for the remaining part of the function.  And in
calculating this integral (K${}^5$.), for the improvement of an
approximation $V_1$, we may employ the following analogous
approximations to the rigorous formul{\ae} (D.) and (E.),
$$\left. \multieqalign{
{\delta V_1 \over \delta a_1} &= - m_1 a_1'; &
{\delta V_1 \over \delta a_2} &= - m_2 a_2'; \quad \ldots &
{\delta V_1 \over \delta a_n} &= - m_n a_n'; \cr
{\delta V_1 \over \delta b_1} &= - m_1 b_1'; &
{\delta V_1 \over \delta b_2} &= - m_2 b_2'; \quad \ldots &
{\delta V_1 \over \delta b_n} &= - m_n b_n'; \cr
{\delta V_1 \over \delta c_1} &= - m_1 c_1'; &
{\delta V_1 \over \delta c_2} &= - m_2 c_2'; \quad \ldots &
{\delta V_1 \over \delta c_n} &= - m_n c_n'; \cr}
   \right\}
   \eqno {\rm (L^5.)}$$
and
$${\delta V_1 \over \delta H} = t;
   \eqno {\rm (M^5.)}$$
or with any other marks of final and initial position, (instead
of rectangular coordinates,) the following approximate forms of
the rigorous equations (S.),
$${\delta V_1 \over \delta e_1}
   = - {\delta T_0 \over \delta e_1'},\quad
  {\delta V_1 \over \delta e_2}
   = - {\delta T_0 \over \delta e_2'},\quad \ldots \quad
  {\delta V_1 \over \delta e_{3n}}
   = - {\delta T_0 \over \delta e_{3n}'},
   \eqno {\rm (N^5.)}$$
together with the formula (M${}^5$.); by which new formul{\ae} the
manner of motion of the system is approximately though not
rigorously expressed.

It is easy to extend these remarks to problems of relative
motion, and to show that in such problems we have the rigorous
transformation
$$V_{\prime 2}
   = \int_0^t (T_\prime - T_{\prime 1} + T_{\prime 2}) \, dt,
   \eqno {\rm (O^5.)}$$
and the approximate expression
$$V_{\prime 2} = \int_0^t (T_\prime - T_{\prime 1}) \, dt,
   \eqno {\rm (P^5.)}$$
$V_{\prime 1}$ being any approximate value of the function
$V_\prime$ of relative motion, and $V_{\prime 2}$ being the
correction of this value; and $T_{\prime 1}$, $T_{\prime 2}$,
being homogeneous functions of the second dimension, composed of
the partial differential coefficients of these two parts
$V_{\prime 1}$, $V_{\prime 2}$, in the same way as $T_\prime$ is
composed of the coefficients of the whole function $V_\prime$.
These general remarks may usefully be illustrated by a particular
but extensive application.

\bigbreak

{\sectiontitle
Application of the foregoing method to the case of a Ternary
or Multiple System, with any laws of attraction or repulsion,
and with one predominant mass.\par}

\nobreak\bigskip

20.
The value (68.), for the relative living force $2T_\prime$ of a
system, reduces itself successively to the following parts,
$2 T_\prime^{(1)}, 2 T_\prime^{(2)},\ldots\, 2T_\prime^{(n-1)}$,
when we suppose that all the $n - 1$ first masses vanish, with
the exception of each successively; namely, to the part
$$2 T_\prime^{(1)} = {m_1 m_n \over m_1 + m_n}
      (\xi_1'^2 + \eta_1'^2 + \zeta_1'^2),
   \eqno {\rm (132.)}$$
when only $m_1$, $m_n$, do not vanish; the part
$$2 T_\prime^{(2)} = {m_2 m_n \over m_2 + m_n}
      (\xi_2'^2 + \eta_2'^2 + \zeta_2'^2),
   \eqno {\rm (133.)}$$
when all but $m_2$, $m_n$, vanish; and so on, as far as the part
$$2 T_\prime^{(n-1)} = {m_{n-1} m_n \over m_{n-1} + m_n}
      (\xi_{n-1}'^2 + \eta_{n-1}'^2 + \zeta_{n-1}'^2),
   \eqno {\rm (134.)}$$
which remains, when only the two last masses are retained.  The
sum of these $n-1$ parts is not, in general, equal to the whole
relative living force $2T_\prime$ of the system, with all the $n$
masses retained; but it differs little from that whole when the
first $n - 1$ masses are small in comparison with the last mass
$m_n$; for the rigorous value of this difference is, by
(68.), and by (132.) (133.) (134.),
$$\left. \eqalign{
2T_\prime - 2T_\prime^{(1)} - 2T_\prime^{(2)}
         - \cdots - 2T_\prime^{(n-1)}
   \hskip-108pt \cr
   &=   {2m_1 \over m_n} (T_\prime^{(1)} - T_\prime)
      + {2m_2 \over m_n} (T_\prime^{(2)} - T_\prime)
      + \cdots
      + {2m_{n-1} \over m_n} (T_\prime^{(n-1)} - T_\prime) \cr
   &\mathrel{\phantom{=}} \mathord{}
      + {1 \over m_n} \sum_\prime \mathbin{.} m_i m_k \{
         (\xi_i' - \xi_k')^2 + (\eta_i' - \eta_k')^2
         + (\zeta_i' - \zeta_k')^2 \}:\cr}
   \right\}
   \eqno {\rm (135.)}$$
an expression which is small of the second order when the $n - 1$
first masses are small of the first order.  If, then, we denote by
$V_\prime^{(1)}, V_\prime^{(2)},\ldots\, V_\prime^{(n-1)}$, the
relative actions, or accumulated relative living forces, such as they
would be in the $n - 1$ binary systems, $(m_1 \, m_n)$,
$(m_2 \, m_n),\ldots$ $(m_{n-1} \, m_n)$, without the perturbations
of the other small masses of the entire multiple system of $n$ points;
so that
$$V_\prime^{(1)} = \int_0^t 2 T_\prime^{(1)} \, dt,\quad
  V_\prime^{(2)} = \int_0^t 2 T_\prime^{(2)} \, dt,\quad
  \cdots \quad
  V_\prime^{(n-1)} = \int_0^t 2 T_\prime^{(n-1)} \, dt,
   \eqno {\rm (Q^5.)}$$
the perturbations being neglected in calculating these $n - 1$
definite integrals; we shall have, as an approximate value for
the whole relative action $V_\prime$ of the system, the sum
$V_{\prime 1}$ of its values for these separate binary systems,
$$V_{\prime 1}
  = V_\prime^{(1)} + V_\prime^{(2)} + \cdots + V_\prime^{(n-1)}.
   \eqno {\rm (R^5.)}$$
This sum, by our theory of binary systems, may be otherwise
expressed as follows:
$$V_{\prime 1}
  =   {m_1 m_n w^{(1)} \over m_1 + m_n}
    + {m_2 m_n w^{(2)} \over m_2 + m_n}
    + \cdots
    + {m_{n-1} m_n w^{({n-1})} \over m_{n-1} + m_n},
   \eqno {\rm (S^5.)}$$
if we put for abridgement
$$\left. \eqalign{
w^{(1)}
   &= h^{(1)} \vartheta^{(1)}
      + \int_{r_0^{(1)}}^{r^{(1)}} r'^{(1)} \, dr^{(1)},\cr
w^{(2)}
   &= h^{(2)} \vartheta^{(2)}
      + \int_{r_0^{(2)}}^{r^{(2)}} r'^{(2)} \, dr^{(2)},\cr
\noalign{\hbox{$\cdots$}}
w^{(n-1)}
   &= h^{(n-1)} \vartheta^{(n-1)}
      + \int_{r_0^{(n-1)}}^{r^{(n-1)}} r'^{(n-1)} \, dr^{(n-1)}.\cr}
   \right\}
   \eqno {\rm (T^5.)}$$
In this expression,
$$\left. \eqalign{
r'^{(1)} &= \pm \sqrt{ 2(m_1 + m_n) f^{(1)} + 2 g^{(1)}
                      - {h^{(1)2} \over r^{(1)2}} },\cr
\noalign{\hbox{$\cdots$}}
r'^{(n-1)} &= \pm \sqrt{ 2(m_{n-1} + m_n) f^{(n-1)} + 2 g^{(n-1)}
                      - {h^{(n-1)2} \over r^{(n-1)2}} },\cr}
   \right\}
   \eqno {\rm (U^5.)}$$
$r^{(1)}$,~$\ldots$ $r^{(n-1)}$ being abridged expressions for the
distances $r^{(1,n)}$,~$\ldots$ $r^{(n-1,n)}$, and
$f^{(1)}$,~$\ldots$ $f^{(n-1)}$ being abridgements for the functions
$f^{(1,n)}$,~$\ldots$ $f^{(n-1,n)}$, of these distances, of which
the derivatives, according as they are negative or positive,
express the laws of attraction or repulsion: we have also
introduced $2n - 2$ auxiliary quantities
$h^{(1)}$~$g^{(1)}$~$\ldots$ $h^{(n-1)}$~$g^{(n-1)}$, to be
eliminated or determined by the following equations of condition:
$$\left. \eqalign{
0 &= \vartheta^{(1)} + \int_{r_0^{(1)}}^{r^{(1)}}
      {\delta r'^{(1)} \over \delta h^{(1)}} dr^{(1)},\cr
0 &= \vartheta^{(2)} + \int_{r_0^{(2)}}^{r^{(2)}}
      {\delta r'^{(2)} \over \delta h^{(2)}} dr^{(2)},\cr
\noalign{\hbox{$\cdots$}}
0 &= \vartheta^{(n-1)} + \int_{r_0^{(n-1)}}^{r^{(n-1)}}
      {\delta r'^{(n-1)} \over \delta h^{(n-1)}} dr^{(n-1)},\cr}
   \right\}
   \eqno {\rm (V^5.)}$$
and
$$\int_{r_0^{(1)}}^{r^{(1)}} {dr^{(1)} \over r'^{(1)}}
   = \int_{r_0^{(2)}}^{r^{(2)}} {dr^{(2)} \over r'^{(2)}}
   = \cdots
   = \int_{r_0^{(n-1)}}^{r^{(n-1)}} {dr^{(n-1)} \over r'^{(n-1)}},
   \eqno {\rm (W^5.)}$$
or
$${\delta w^{(1)} \over \delta g^{(1)}}
   = {\delta w^{(2)} \over \delta g^{(2)}}
   = \cdots
   = {\delta w^{(n-1)} \over \delta g^{(n-1)}},
   \eqno {\rm (X^5.)}$$
along with this last condition,
$${m_1 g^{(1)} \over m_1 + m_n}
      + {m_2 g^{(2)} \over m_2 + m_n}
      + {m_3 g^{(3)} \over m_3 + m_n}
      + \cdots
      + {m_{n-1} g^{(n-1)} \over m_{n-1} + m_n}
   = {H_\prime \over m_n};
   \eqno {\rm (Y^5.)}$$
and we have denoted by
$\vartheta^{(1)}$,~$\ldots$~$\vartheta^{(n-1)}$, the angles which
the final distances $r^{(1)}$,~$\ldots$ $r^{(n-1)}$, of the
first $n - 1$ points from the last or $n$th point of the system,
make respectively with the initial distances corresponding,
namely, $r_0^{(1)}$,~$\ldots$~$r_0^{(n-1)}$.  The variation of
the sum $V_{\prime 1}$ is, by (S${}^5$.),
$$\delta V_{\prime 1}
   =   {m_1 m_n \delta w^{(1)} \over m_1 + m_n}
     + {m_2 m_n \delta w^{(2)} \over m_2 + m_n}
     + \cdots
     + {m_{n-1} m_n \delta w^{(n-1)} \over m_{n-1} + m_n};
   \eqno {\rm (Z^5.)}$$
in which, by the equations of condition, we may treat all the
auxiliary quantities
$h^{(1)}$~$g^{(1)}$~$\ldots$ $h^{(n-1)}$~$g^{(n-1)}$
as constant, if $H_\prime$ be considered as given: so that the
part of this variation $\delta V_{\prime 1}$, which depends on
the variations of the final relative coordinates, may be put
under the form,
$$\left. \eqalign{
\delta_{\xi,\eta,\zeta} V_{\prime 1}
   &=   {m_1 m_n \over m_1 + m_n} \left(
           {\delta w^{(1)} \over \delta \xi_1} \delta \xi_1
         + {\delta w^{(1)} \over \delta \eta_1} \delta \eta_1
         + {\delta w^{(1)} \over \delta \zeta_1} \delta \zeta_1
           \right) \cr
     &+ {m_2 m_n \over m_2 + m_n} \left(
           {\delta w^{(2)} \over \delta \xi_2} \delta \xi_2
         + {\delta w^{(2)} \over \delta \eta_2} \delta \eta_2
         + {\delta w^{(2)} \over \delta \zeta_2} \delta \zeta_2
           \right) \cr
     &+ \cdots \cr
     &+ {m_{n-1} m_n \over m_{n-1} + m_n} \left(
           {\delta w^{(n-1)} \over \delta \xi_{n-1}} \delta \xi_{n-1}
         + {\delta w^{(n-1)} \over \delta \eta_{n-1}} \delta \eta_{n-1}
         + {\delta w^{(n-1)} \over \delta \zeta_{n-1}} \delta \zeta_{n-1}
           \right).\cr}
   \right\}
   \eqno {\rm (A^6.)}$$

By the equations (T${}^5$.) (U${}^5$.), or by the theory of
binary systems, we have, rigorously,
$$\left. \eqalign{
\left( {\delta w^{(1)} \over \delta \xi_1} \right)^2
   + \left( {\delta w^{(1)} \over \delta \eta_1} \right)^2
   + \left( {\delta w^{(1)} \over \delta \zeta_1} \right)^2
   &= 2 (m_1 + m_n) f^{(1)} + 2 g^{(1)};\cr
\left( {\delta w^{(2)} \over \delta \xi_2} \right)^2
   + \left( {\delta w^{(2)} \over \delta \eta_2} \right)^2
   + \left( {\delta w^{(2)} \over \delta \zeta_2} \right)^2
   &= 2 (m_2 + m_n) f^{(2)} + 2 g^{(2)};\cr
\noalign{\hbox{$\cdots$}}
\left( {\delta w^{(n-1)} \over \delta \xi_{n-1}} \right)^2
   + \left( {\delta w^{(n-1)} \over \delta \eta_{n-1}} \right)^2
   + \left( {\delta w^{(n-1)} \over \delta \zeta_{n-1}} \right)^2
   &= 2 (m_{n-1} + m_n) f^{(n-1)} + 2 g^{(n-1)};\cr}
   \right\}
   \eqno {\rm (B^6.)}$$
and the rigorous law of relative living force for the whole
multiple system, is
$$T_\prime = U + H_\prime,
   \eqno {\rm (50.)}$$
in which
$$U = m_n ( m_1 f^{(1)} + m_2 f^{(2)} + \cdots
            + m_{n-1} f^{(n-1)} )
            + \sum_\prime \mathbin{.} m_i m_k f^{(i,k)},
   \eqno {\rm (C^6.)}$$
and
$$\left. \eqalign{
T_\prime
   &= {1 \over 2}
      \left( {1 \over m_1} + {1 \over m_n} \right)
         \left\{
              \left( {\delta V_\prime \over \delta \xi_1} \right)^2
            + \left( {\delta V_\prime \over \delta \eta_1} \right)^2
            + \left( {\delta V_\prime \over \delta \zeta_1} \right)^2
              \right\} \cr
   &+ {1 \over 2}
      \left( {1 \over m_2} + {1 \over m_n} \right)
         \left\{
              \left( {\delta V_\prime \over \delta \xi_2} \right)^2
            + \left( {\delta V_\prime \over \delta \eta_2} \right)^2
            + \left( {\delta V_\prime \over \delta \zeta_2} \right)^2
              \right\} \cr
   &+ \cdots \cr
   &+ {1 \over 2}
      \left( {1 \over m_{n-1}} + {1 \over m_n} \right)
         \left\{
              \left( {\delta V_\prime \over \delta \xi_{n-1}} \right)^2
            + \left( {\delta V_\prime \over \delta \eta_{n-1}} \right)^2
            + \left( {\delta V_\prime \over \delta \zeta_{n-1}} \right)^2
              \right\} \cr
   &+ {1 \over m_n} \sum_\prime \left(
            {\delta V_\prime \over \delta \xi_i}
            {\delta V_\prime \over \delta \xi_k}
          + {\delta V_\prime \over \delta \eta_i}
            {\delta V_\prime \over \delta \eta_k}
          + {\delta V_\prime \over \delta \zeta_i}
            {\delta V_\prime \over \delta \zeta_k}
            \right).\cr}
   \right\}
   \eqno {\rm (D^6.)}$$

We have therefore, by changing in this last expression the
coefficients of the characteristic function $V_\prime$ to those
of its first part $V_{\prime 1}$, and by attending to the
foregoing equations,
$$\left. \eqalign{
T_{\prime 1}
   &= m_n \sum_\prime \mathbin{.} m_i f^{(i)} + H_\prime \cr
      &+ m_n \sum_\prime \mathbin{.}
         {m_i \over m_n + m_i} {m_k \over m_n + m_k} \left(
            {\delta w^{(i)} \over \delta \xi_i}
            {\delta w^{(k)} \over \delta \xi_k}
          + {\delta w^{(i)} \over \delta \eta_i}
            {\delta w^{(k)} \over \delta \eta_k}
          + {\delta w^{(i)} \over \delta \zeta_i}
            {\delta w^{(k)} \over \delta \zeta_k}
            \right);\cr}
   \right\}
   \eqno {\rm (E^6.)}$$
and consequently
$$\left. \eqalign{
T_\prime - T_{\prime 1}
   &= \sum_\prime m_i m_k \biggl\{
      f^{(i,k)} \cr
   &\mathrel{\phantom{=}} \mathord{}
       - {m_n \over (m_n + m_i) (m_n + m_k)} \left(
            {\delta w^{(i)} \over \delta \xi_i}
            {\delta w^{(k)} \over \delta \xi_k}
          + {\delta w^{(i)} \over \delta \eta_i}
            {\delta w^{(k)} \over \delta \eta_k}
          + {\delta w^{(i)} \over \delta \zeta_i}
            {\delta w^{(k)} \over \delta \zeta_k}
            \right)
      \biggr\}.\cr}
   \right\}
   \eqno {\rm (F^6.)}$$
The general transformation of the foregoing number gives
therefore, rigorously, for the remaining part $V_{\prime 2}$ of
the characteristic function $V_\prime$ of relative motion of the
multiple system, the equation
$$\left. \eqalign{
V_{\prime 2}
   &= \int_0^t T_{\prime 2} \, dt \cr
   &\mathrel{\phantom{=}} \mathord{}
      + \sum_\prime \mathbin{.} m_i m_k \int_0^t \left\{
      f^{(i,k)} -
          {\displaystyle
            {\delta w^{(i)} \over \delta \xi_i}
            {\delta w^{(k)} \over \delta \xi_k}
          + {\delta w^{(i)} \over \delta \eta_i}
            {\delta w^{(k)} \over \delta \eta_k}
          + {\delta w^{(i)} \over \delta \zeta_i}
            {\delta w^{(k)} \over \delta \zeta_k}
         \over \displaystyle
            {1 \over m_n} (m_n + m_i)(m_n + m_k)}
      \right\} dt;\cr}
   \right\}
   \eqno {\rm (G^6.)}$$
and, approximately, the expression
$$V_{\prime 2} = \sum_\prime \mathbin{.} m_i m_k \int_0^t
      \left\{ f^{(i,k)} - {1 \over m_n}
         ( \xi_i' \xi_k' + \eta_i' \eta_k'
           + \zeta_i' \zeta_k' )
           \right\} dt:
   \eqno {\rm (H^6.)}$$
with which last expression we may combine the following
approximate formul{\ae} belonging in rigour to binary systems
only,
$$\xi_i'    =   {\delta w^{(i)} \over \delta \xi_i},\quad
  \eta_i'   =   {\delta w^{(i)} \over \delta \eta_i},\quad
  \zeta_i'  =   {\delta w^{(i)} \over \delta \zeta_i},
   \eqno {\rm (I^6.)}$$
$$\alpha_i' = - {\delta w^{(i)} \over \delta \alpha_i},\quad
  \beta_i'  = - {\delta w^{(i)} \over \delta \beta_i},\quad
  \gamma_i' = - {\delta w^{(i)} \over \delta \gamma_i},
   \eqno {\rm (K^6.)}$$
and
$$t = {\delta w^{(i)} \over \delta g^{(i)}}.
   \eqno {\rm (L^6.)}$$

We have also, rigorously, for binary systems, the following
differential equations of motion of the second order,
$$\xi_i''
   = (m_n + m_i) {\delta f^{(i)} \over \delta \xi_i};\quad
  \eta_i''
   = (m_n + m_i) {\delta f^{(i)} \over \delta \eta_i};\quad
  \zeta_i''
   = (m_n + m_i) {\delta f^{(i)} \over \delta \zeta_i};
   \eqno {\rm (M^6.)}$$
which enable us to transform in various ways the approximate
expression (H${}^6$.).  Thus, in the case of a ternary system,
with any laws of attraction or repulsion, but with one
predominant mass $m_3$, the {\it disturbing part\/}
$V_{\prime 2}$ of the characteristic function $V_\prime$ of
relative motion, may be put under the form
$$V_{\prime 2} = m_1 m_2 W,
   \eqno {\rm (N^6.)}$$
in which the coefficient $W$ may be approximately be expressed as
follows:
$$W = \int_0^t \left\{ f^{(1,2)}
      - {1 \over m_3} (\xi_1' \xi_2' + \eta_1' \eta_2'
                       + \zeta_1' \zeta_2') \right\} dt,
   \eqno {\rm (O^6.)}$$
or thus:
$$\left. \eqalign{
W  &= \int_0^t \left( f^{(1,2)}
         + \xi_2    {\delta f^{(1)} \over \delta \xi_1}
         + \eta_2   {\delta f^{(1)} \over \delta \eta_1}
         + \zeta_2  {\delta f^{(1)} \over \delta \zeta_1}
         \right) \,dt \cr
   &\mathrel{\phantom{=}} \mathord{}
      - {1 \over m_3} \left(
           \xi_2    {\delta w^{(1)} \over \delta \xi_1}
         + \eta_2   {\delta w^{(1)} \over \delta \eta_1}
         + \zeta_2  {\delta w^{(1)} \over \delta \zeta_1}
         + \alpha_2 {\delta w^{(1)} \over \delta \alpha_1}
         + \beta_2  {\delta w^{(1)} \over \delta \beta_1}
         + \gamma_2 {\delta w^{(1)} \over \delta \gamma_1}
         \right),\cr}
   \right\}
   \eqno {\rm (P^6.)}$$
or finally,
$$\left. \eqalign{
W  &= \int_0^t \left( f^{(1,2)}
         + \xi_1    {\delta f^{(2)} \over \delta \xi_2}
         + \eta_1   {\delta f^{(2)} \over \delta \eta_2}
         + \zeta_1  {\delta f^{(2)} \over \delta \zeta_2}
         \right) \,dt \cr
   &\mathrel{\phantom{=}} \mathord{}
      - {1 \over m_3} \left(
           \xi_1    {\delta w^{(2)} \over \delta \xi_2}
         + \eta_1   {\delta w^{(2)} \over \delta \eta_2}
         + \zeta_1  {\delta w^{(2)} \over \delta \zeta_2}
         + \alpha_1 {\delta w^{(2)} \over \delta \alpha_2}
         + \beta_1  {\delta w^{(2)} \over \delta \beta_2}
         + \gamma_1 {\delta w^{(2)} \over \delta \gamma_2}
         \right).\cr}
   \right\}
   \eqno {\rm (Q^6.)}$$

In general, for a multiple system, we may put
$$V_{\prime 2} = \sum_\prime \mathbin{.} m_i m_k W^{(i,k)};
   \eqno {\rm (R^6.)}$$
and approximately,
$$\left. \eqalign{
W^{(i,k)}  &= \int_0^t \left( f^{(i,k)}
         + \xi_k    {\delta f^{(i)} \over \delta \xi_i}
         + \eta_k   {\delta f^{(i)} \over \delta \eta_i}
         + \zeta_k  {\delta f^{(i)} \over \delta \zeta_i}
         \right) \,dt \cr
   &\mathrel{\phantom{=}} \mathord{}
      - {1 \over m_n} \left(
           \xi_k    {\delta w^{(i)} \over \delta \xi_i}
         + \eta_k   {\delta w^{(i)} \over \delta \eta_i}
         + \zeta_k  {\delta w^{(i)} \over \delta \zeta_i}
         + \alpha_k {\delta w^{(i)} \over \delta \alpha_i}
         + \beta_k  {\delta w^{(i)} \over \delta \beta_i}
         + \gamma_k {\delta w^{(i)} \over \delta \gamma_i}
         \right),\cr}
   \right\}
   \eqno {\rm (S^6.)}$$
or
$$\left. \eqalign{
W^{(i,k)}  &= \int_0^t \left( f^{(i,k)}
         + \xi_i    {\delta f^{(k)} \over \delta \xi_k}
         + \eta_i   {\delta f^{(k)} \over \delta \eta_k}
         + \zeta_i  {\delta f^{(k)} \over \delta \zeta_k}
         \right) \,dt \cr
   &\mathrel{\phantom{=}} \mathord{}
      - {1 \over m_n} \left(
           \xi_i    {\delta w^{(k)} \over \delta \xi_k}
         + \eta_i   {\delta w^{(k)} \over \delta \eta_k}
         + \zeta_i  {\delta w^{(k)} \over \delta \zeta_k}
         + \alpha_i {\delta w^{(k)} \over \delta \alpha_k}
         + \beta_i  {\delta w^{(k)} \over \delta \beta_k}
         + \gamma_i {\delta w^{(k)} \over \delta \gamma_k}
         \right).\cr}
   \right\}
   \eqno {\rm (T^6.)}$$

\bigbreak

{\sectiontitle
Rigorous transition from the theory of Binary to that of
Multiple Systems, by means of the disturbing part of the whole
Characteristic Function; and approximate expressions for the
perturbations.\par}

\nobreak\bigskip

21.
The three equations (K${}^6$.) when the auxiliary constant
$g^{(i)}$ is eliminated by the formula (L${}^6$.) are rigorously
(by our theory) the three final integrals of the three known
equations of the second order (M${}^6$.), for the relative motion
of the binary system $(m_i m_n)$; and give, for such a system,
the three varying relative coordinates $\xi_i$~$\eta_i$~$\zeta_i$,
as functions of their initial values and initial
rates of increase $\alpha_i$~$\beta_i$~$\gamma_i$
$\alpha_i'$~$\beta_i'$~$\gamma_i'$, and of the time~$t$.  In
like manner the three equations (I${}^6$.), when $g^{(i)}$ is
eliminated by (L${}^6$.), are rigorously the three intermediate
integrals of the same known differential equations of motion of
the same binary system.  These integrals, however, cease to be
rigorous when we introduce the perturbations of the relative
motion of this partial or binary system $(m_i \, m_n)$, arising from
the attractions or repulsions of the other points $m_k$, of
the whole proposed multiple system; but they may be corrected and
rendered rigorous by employing the remaining part $V_{\prime 2}$
of the whole characteristic function of relative motion
$V_\prime$, along with the principal part or approximate value
$V_{\prime 1}$.

The equations (X${}^1$.), (Y${}^1$.) of the twelfth number, give
rigorously
$$\xi_i' = {1 \over m_i}
         {\delta V_\prime \over \delta \xi_i}
   + {1 \over m_n} \sum_\prime
         {\delta V_\prime \over \delta \xi_i},\quad
  \eta_i' = {1 \over m_i}
         {\delta V_\prime \over \delta \eta_i}
   + {1 \over m_n} \sum_\prime
         {\delta V_\prime \over \delta \eta_i},\quad
  \zeta_i' = {1 \over m_i}
         {\delta V_\prime \over \delta \zeta_i}
   + {1 \over m_n} \sum_\prime
         {\delta V_\prime \over \delta \zeta_i},
   \eqno {\rm (U^6.)}$$
and
$$- \alpha_i' = {1 \over m_i}
         {\delta V_\prime \over \delta \alpha_i}
   + {1 \over m_n} \sum_\prime
         {\delta V_\prime \over \delta \alpha_i},\quad
  - \beta_i' = {1 \over m_i}
         {\delta V_\prime \over \delta \beta_i}
   + {1 \over m_n} \sum_\prime
         {\delta V_\prime \over \delta \beta_i},\quad
  - \gamma_i' = {1 \over m_i}
         {\delta V_\prime \over \delta \gamma_i}
   + {1 \over m_n} \sum_\prime
         {\delta V_\prime \over \delta \gamma_i},
   \eqno {\rm (V^6.)}$$
and therefore, by (A${}^6$.),
$$\left. \eqalign{
{\delta w^{(i)} \over \delta \xi_i}
   &= \xi_i'
      - \sum_{\prime\prime} \mathbin{.} {m_k \over m_k + m_n}
         {\delta w^{(k)} \over \delta \xi_k}
      - {1 \over m_i}
         {\delta V_{\prime 2} \over \delta \xi_i}
      - {1 \over m_n} \sum_\prime
         {\delta V_{\prime 2} \over \delta \xi_i},\cr
{\delta w^{(i)} \over \delta \eta_i}
   &= \eta_i'
      - \sum_{\prime\prime} \mathbin{.} {m_k \over m_k + m_n}
         {\delta w^{(k)} \over \delta \eta_k}
      - {1 \over m_i}
         {\delta V_{\prime 2} \over \delta \eta_i}
      - {1 \over m_n} \sum_\prime
         {\delta V_{\prime 2} \over \delta \eta_i},\cr
{\delta w^{(i)} \over \delta \zeta_i}
   &= \zeta_i'
      - \sum_{\prime\prime} \mathbin{.} {m_k \over m_k + m_n}
         {\delta w^{(k)} \over \delta \zeta_k}
      - {1 \over m_i}
         {\delta V_{\prime 2} \over \delta \zeta_i}
      - {1 \over m_n} \sum_\prime
         {\delta V_{\prime 2} \over \delta \zeta_i},\cr}
   \right\}
   \eqno {\rm (W^6.)}$$
and similarly
$$\left. \eqalign{
- {\delta w^{(i)} \over \delta \alpha_i}
   &= \alpha_i'
      + \sum_{\prime\prime} \mathbin{.} {m_k \over m_k + m_n}
         {\delta w^{(k)} \over \delta \alpha_k}
      + {1 \over m_i}
         {\delta V_{\prime 2} \over \delta \alpha_i}
      + {1 \over m_n} \sum_\prime
         {\delta V_{\prime 2} \over \delta \alpha_i},\cr
- {\delta w^{(i)} \over \delta \beta_i}
   &= \beta_i'
      + \sum_{\prime\prime} \mathbin{.} {m_k \over m_k + m_n}
         {\delta w^{(k)} \over \delta \beta_k}
      + {1 \over m_i}
         {\delta V_{\prime 2} \over \delta \beta_i}
      + {1 \over m_n} \sum_\prime
         {\delta V_{\prime 2} \over \delta \beta_i},\cr
- {\delta w^{(i)} \over \delta \gamma_i}
   &= \gamma_i'
      + \sum_{\prime\prime} \mathbin{.} {m_k \over m_k + m_n}
         {\delta w^{(k)} \over \delta \gamma_k}
      + {1 \over m_i}
         {\delta V_{\prime 2} \over \delta \gamma_i}
      + {1 \over m_n} \sum_\prime
         {\delta V_{\prime 2} \over \delta \gamma_i},\cr}
   \right\}
   \eqno {\rm (X^6.)}$$
the sign of summation $\sum_{\prime\prime}$ referring only to the
disturbing masses $m_k$, to the exclusion of $m_i$ and $m_n$; and
these equations (W${}^6$.), (X${}^6$.) are the rigorous
formul{\ae}, corresponding to the approximate relations
(I${}^6$.), (K${}^6$.).  In like manner, the formula (L${}^6$.)
for the time of motion in a binary system, which is only an
approximation when the system is considered as multiple, may be
rigorously corrected for perturbation by adding to it an
analogous term deduced from the disturbing part $V_{\prime 2}$ of
the whole characteristic function; that is, by changing it to the
following:
$$t = {\delta w^{(i)} \over \delta g^{(i)}}
      + {\delta V_{\prime 2} \over \delta H_\prime},
   \eqno {\rm (Y^6.)}$$
which gives, for this other coefficient of $w^{(i)}$, the
corrected and rigorous expression
$${\delta w^{(i)} \over \delta g^{(i)}}
   = t - {\delta V_{\prime 2} \over \delta H_\prime}:
   \eqno {\rm (Z^6.)}$$
$V_{\prime 2}$ being here supposed so chosen as to be rigorously
the correction to $V_{\prime 1}$.  If therefore, by the theory of
binary systems, or by eliminating $g^{(i)}$ between the four
equations (K${}^6$.) (L${}^6$.), we have deduced expressions for
the three varying relative coordinates $\xi_i$~$\eta_i$~$\zeta_i$
as functions of the time~$t$, and of the six initial quantities
$\alpha_i$~$\beta_i$~$\gamma_i$
$\alpha_i'$~$\beta_i'$~$\gamma_i'$,
which may be thus denoted,
$$\left. \eqalign{
\xi_i   &= \phi_1( \alpha_i, \beta_i, \gamma_i,
                   \alpha_i', \beta_i', \gamma_i',t),\cr
\eta_i  &= \phi_2( \alpha_i, \beta_i, \gamma_i,
                   \alpha_i', \beta_i', \gamma_i',t),\cr
\zeta_i &= \phi_3( \alpha_i, \beta_i, \gamma_i,
                   \alpha_i', \beta_i', \gamma_i',t);\cr}
   \right\}
   \eqno {\rm (A^7.)}$$
we shall know that the following relations are rigorously and
{\it identically\/} true,
$$\left. \eqalign{
\xi_i   &= \phi_1 \left( \alpha_i, \beta_i, \gamma_i,
                   - {\delta w^{(i)} \over \delta \alpha_i},
                   - {\delta w^{(i)} \over \delta \beta_i},
                   - {\delta w^{(i)} \over \delta \gamma_i},
                     {\delta w^{(i)} \over \delta g^{(i)}}
                     \right),\cr
\eta_i  &= \phi_2 \left( \alpha_i, \beta_i, \gamma_i,
                   - {\delta w^{(i)} \over \delta \alpha_i},
                   - {\delta w^{(i)} \over \delta \beta_i},
                   - {\delta w^{(i)} \over \delta \gamma_i},
                     {\delta w^{(i)} \over \delta g^{(i)}}
                     \right),\cr
\zeta_i &= \phi_3 \left( \alpha_i, \beta_i, \gamma_i,
                   - {\delta w^{(i)} \over \delta \alpha_i},
                   - {\delta w^{(i)} \over \delta \beta_i},
                   - {\delta w^{(i)} \over \delta \gamma_i},
                     {\delta w^{(i)} \over \delta g^{(i)}}
                     \right),\cr}
   \right\}
   \eqno {\rm (B^7.)}$$
and consequently that these relations will still be rigorously
true when we substitute for the four coefficients of $w^{(i)}$
their rigorous values (X${}^6$.) and (Z${}^6$.) for the case of a
multiple system.  We may thus retain in rigour for any multiple
system the final integrals (A${}^7$.) of the motion of a binary
system, if only we add to the initial components
$\alpha_i'$~$\beta_i'$~$\gamma_i'$
of relative velocity, and to the time~$t$, the following
perturbational terms:
$$\left. \eqalign{
\Delta \alpha_i'
   &= \sum_{\prime\prime} \mathbin{.} {m_k \over m_k + m_n}
         {\delta w^{(k)} \over \delta \alpha_k}
      + {1 \over m_i} {\delta V_{\prime 2} \over \delta \alpha_i}
      + {1 \over m_n} \sum_\prime
           {\delta V_{\prime 2} \over \delta \alpha_i},\cr
\Delta \beta_i'
   &= \sum_{\prime\prime} \mathbin{.} {m_k \over m_k + m_n}
         {\delta w^{(k)} \over \delta \beta_k}
      + {1 \over m_i} {\delta V_{\prime 2} \over \delta \beta_i}
      + {1 \over m_n} \sum_\prime
           {\delta V_{\prime 2} \over \delta \beta_i},\cr
\Delta \gamma_i'
   &= \sum_{\prime\prime} \mathbin{.} {m_k \over m_k + m_n}
         {\delta w^{(k)} \over \delta \gamma_k}
      + {1 \over m_i} {\delta V_{\prime 2} \over \delta \gamma_i}
      + {1 \over m_n} \sum_\prime
           {\delta V_{\prime 2} \over \delta \gamma_i},\cr}
   \right\}
   \eqno {\rm (C^7.)}$$
and
$$\Delta t = - {\delta V_{\prime 2} \over \delta H_\prime}.
   \eqno {\rm (D^7.)}$$

In the same way, if the theory of binary systems, or the
elimination of $g^{(i)}$ between the four equations (I${}^6$.)
(L${}^6$.), has given three intermediate integrals, of the form
$$\left. \eqalign{
\xi_i'   &= \psi_1( \xi_i, \eta_i, \zeta_i,
                   \alpha_i, \beta_i, \gamma_i,t),\cr
\eta_i'  &= \psi_2( \xi_i, \eta_i, \zeta_i,
                   \alpha_i, \beta_i, \gamma_i,t),\cr
\zeta_i' &= \psi_3( \xi_i, \eta_i, \zeta_i,
                   \alpha_i, \beta_i, \gamma_i,t),\cr}
   \right\}
   \eqno {\rm (E^7.)}$$
we can conclude that the following equations are rigorous and
identical,
$$\left. \eqalign{
{\delta w^{(i)} \over \delta \xi_i}
   &= \psi_1 \left( \xi_i, \eta_i, \zeta_i,
         \alpha_i, \beta_i, \gamma_i,
         {\delta w^{(i)} \over \delta g^{(i)}}
         \right),\cr
{\delta w^{(i)} \over \delta \eta_i}
   &= \psi_2 \left( \xi_i, \eta_i, \zeta_i,
         \alpha_i, \beta_i, \gamma_i,
         {\delta w^{(i)} \over \delta g^{(i)}}
         \right),\cr
{\delta w^{(i)} \over \delta \zeta_i}
   &= \psi_3 \left( \xi_i, \eta_i, \zeta_i,
         \alpha_i, \beta_i, \gamma_i,
         {\delta w^{(i)} \over \delta g^{(i)}}
         \right),\cr}
   \right\}
   \eqno {\rm (F^7.)}$$
and must therefore be still true, when, in passing to a multiple
system, we change the coefficients of $w^{(i)}$ to their rigorous
values (W${}^6$.) (Z${}^6$.).  The three intermediate integrals
(E${}^7$.) of the motion of a binary system may therefore be
adapted rigorously to the case of a multiple system, by first
adding to the time~$t$ the perturbational term (D${}^7$.), and
afterwards adding to the resulting values of the final components
of relative velocity the terms
$$\left. \eqalign{
\Delta \xi_i'
   &= \sum_{\prime\prime} \mathbin{.} {m_k \over m_k + m_n}
         {\delta w^{(k)} \over \delta \xi_k}
      + {1 \over m_i} {\delta V_{\prime 2} \over \delta \xi_i}
      + {1 \over m_n} \sum_\prime
           {\delta V_{\prime 2} \over \delta \xi_i},\cr
\Delta \eta_i'
   &= \sum_{\prime\prime} \mathbin{.} {m_k \over m_k + m_n}
         {\delta w^{(k)} \over \delta \eta_k}
      + {1 \over m_i} {\delta V_{\prime 2} \over \delta \eta_i}
      + {1 \over m_n} \sum_\prime
           {\delta V_{\prime 2} \over \delta \eta_i},\cr
\Delta \zeta_i'
   &= \sum_{\prime\prime} \mathbin{.} {m_k \over m_k + m_n}
         {\delta w^{(k)} \over \delta \zeta_k}
      + {1 \over m_i} {\delta V_{\prime 2} \over \delta \zeta_i}
      + {1 \over m_n} \sum_\prime
           {\delta V_{\prime 2} \over \delta \zeta_i}.\cr}
   \right\}
   \eqno {\rm (G^7.)}$$

\bigbreak

22.
To derive now, from these rigorous results, some useful
approximate expressions, we shall neglect, in the perturbations,
the terms which are of the second order, with respect to the
small masses of the system, and with respect to the constant
$2H_\prime$ of relative living force, which is easily seen to be
small of the same order as the masses: and then the perturbations
of these coordinates, deduced by the method that has been
explained, become
$$\left. \eqalign{
\Delta \xi_i
   &=    {\delta \xi_i \over \delta \alpha_i'}
            \Delta \alpha_i'
       + {\delta \xi_i \over \delta \beta_i'}
            \Delta \beta_i'
       + {\delta \xi_i \over \delta \gamma_i'}
            \Delta \gamma_i'
       + {\delta \xi_i \over \delta t}
            \Delta t,\cr
\Delta \eta_i
   &=    {\delta \eta_i \over \delta \alpha_i'}
            \Delta \alpha_i'
       + {\delta \eta_i \over \delta \beta_i'}
            \Delta \beta_i'
       + {\delta \eta_i \over \delta \gamma_i'}
            \Delta \gamma_i'
       + {\delta \eta_i \over \delta t}
            \Delta t,\cr
\Delta \zeta_i
   &=    {\delta \zeta_i \over \delta \alpha_i'}
            \Delta \alpha_i'
       + {\delta \zeta_i \over \delta \beta_i'}
            \Delta \beta_i'
       + {\delta \zeta_i \over \delta \gamma_i'}
            \Delta \gamma_i'
       + {\delta \zeta_i \over \delta t}
            \Delta t,\cr}
   \right\}
   \eqno {\rm (H^7.)}$$
in which we may employ, instead of the rigorous values (C${}^7$.)
for $\Delta \alpha_i'$, $\Delta \beta_i'$, $\Delta \gamma_i'$,
the following approximate values:
$$\left. \eqalign{
\Delta \alpha_i'
   &= \sum_{\prime\prime} {m_k \over m_n}
         {\delta w^{(k)} \over \delta \alpha_k}
      + {1 \over m_i}
         {\delta V_{\prime 2} \over \delta \alpha_i},\cr
\Delta \beta_i'
   &= \sum_{\prime\prime} {m_k \over m_n}
         {\delta w^{(k)} \over \delta \beta_k}
      + {1 \over m_i}
         {\delta V_{\prime 2} \over \delta \beta_i},\cr
\Delta \gamma_i'
   &= \sum_{\prime\prime} {m_k \over m_n}
         {\delta w^{(k)} \over \delta \gamma_k}
      + {1 \over m_i}
         {\delta V_{\prime 2} \over \delta \gamma_i}.\cr}
   \right\}
   \eqno {\rm (I^7.)}$$

To calculate the four coefficients
$${\delta V_{\prime 2} \over \delta \alpha_i},\quad
  {\delta V_{\prime 2} \over \delta \beta_i},\quad
  {\delta V_{\prime 2} \over \delta \gamma_i},\quad
  {\delta V_{\prime 2} \over \delta H_\prime},$$
which enter into the values (I${}^7$.) (D${}^7$.), we may
consider $V_{\prime 2}$, by (R${}^6$.) (T${}^6$.), and by the
theory of binary systems, as a function of the initial and final
relative coordinates, and initial components of relative
velocities, involving also expressly the time~$t$ and the $n - 2$
auxiliary quantities of the form $g^{(k)}$; and then we are to
consider those initial components and auxiliary quantities and
the time, as depending themselves on the initial and final
coordinates, and on $H_\prime$.  But it is not difficult to
prove, by the foregoing principles, that when $t$ and $g^{(k)}$
are thus considered, their variations are, in the present order
of approximation,
$$\delta t = {\displaystyle \sum_\prime \mathbin{.} m
      \left( {\delta^2 w \over \delta g^2} \right)^{-1}
      \delta_\prime {\delta w \over \delta g}
      + \delta H_\prime
   \over \displaystyle
      \sum_\prime \mathbin{.} m
      \left( {\delta^2 w \over \delta g^2} \right)^{-1} }
   \eqno {\rm (K^7.)}$$
and
$$\delta g^{(k)}
   = \left( {\delta^2 w^{(k)} \over \delta g^{(k)2}} \right)^{-1}
      \left( \delta t
         - \delta_\prime {\delta w^{(k)} \over \delta g^{(k)}}
         \right),
   \eqno {\rm (L^7.)}$$
the sign of variation $\delta_\prime$ referring only to the
initial and final coordinates; and also that
$${\delta^2 w^{(i)} \over \delta g^{(i)2}}
         {\delta \xi_i \over \delta t}
   =   {\delta^2 w^{(i)} \over \delta \alpha_i \, \delta g^{(i)}}
         {\delta \xi_i \over \delta \alpha_i'}
     + {\delta^2 w^{(i)} \over \delta \beta_i \, \delta g^{(i)}}
         {\delta \xi_i \over \delta \beta_i'} 
     + {\delta^2 w^{(i)} \over \delta \gamma_i \, \delta g^{(i)}}
         {\delta \xi_i \over \delta \gamma_i'},
   \eqno {\rm (M^7.)}$$
along with two other analogous relations between the coefficients
of the two other coordinates $\eta_i$, $\zeta_i$; from which it
follows that $t$ and $g^{(k)}$, and therefore
$\alpha_k'$~$\beta_k'$~$\gamma_k'$,
may be treated as constant, in taking the variation of the
disturbing part $V_{\prime 2}$, for the purpose of calculating
the perturbations (H${}^7$.): and that the terms involving
$\Delta t$ are destroyed by other terms.  We may therefore put
simply
$$\left. \eqalign{
\Delta \xi_i
   &=   {\delta \xi_i \over \delta \alpha_i'} \Delta \alpha_i'
      + {\delta \xi_i \over \delta \beta_i'} \Delta \beta_i'
      + {\delta \xi_i \over \delta \gamma_i'} \Delta \gamma_i',\cr
\Delta \eta_i
   &=   {\delta \eta_i \over \delta \alpha_i'} \Delta \alpha_i'
      + {\delta \eta_i \over \delta \beta_i'} \Delta \beta_i'
      + {\delta \eta_i \over \delta \gamma_i'} \Delta \gamma_i',\cr
\Delta \zeta_i
   &=   {\delta \zeta_i \over \delta \alpha_i'} \Delta \alpha_i'
      + {\delta \zeta_i \over \delta \beta_i'} \Delta \beta_i'
      + {\delta \zeta_i \over \delta \gamma_i'} \Delta \gamma_i',\cr}
   \right\}
   \eqno {\rm (N^7.)}$$
employing for $\Delta \alpha_i'$ the following new expression,
$$\left. \eqalign{
\Delta \alpha_i'
   = \sum_{\prime\prime} \mathbin{.} m_k \biggl\{
    &    \int_0^t {\delta R^{(i,k)} \over \delta \alpha_i} dt
      + {\delta \alpha_i' \over \delta \alpha_i}
         \int_0^t {\delta R^{(i,k)} \over \delta \alpha_i'} dt \cr
    & + {\delta \beta_i' \over \delta \alpha_i}
         \int_0^t {\delta R^{(i,k)} \over \delta \beta_i'} dt
      + {\delta \gamma_i' \over \delta \alpha_i}
         \int_0^t {\delta R^{(i,k)} \over \delta \gamma_i'} dt
      \biggr\} \cr}
   \right\}
   \eqno {\rm (O^7.)}$$
together with analogous expressions for $\Delta \beta_i'$,
$\Delta \gamma_i'$, in which the sign of summation
$\sum_{\prime\prime}$ refers to the disturbing masses, and in
which the quantity
$$R^{(i,k)} = f^{(i,k)}
      + \xi_i {\delta f^{(k)} \over \delta \xi_k}
      + \eta_i {\delta f^{(k)} \over \delta \eta_k}
      + \zeta_i {\delta f^{(k)} \over \delta \zeta_k}
   \eqno {\rm (P^7.)}$$
is considered as depending on
$\alpha_i$~$\beta_i$~$\gamma_i$
$\alpha_i'$~$\beta_i'$~$\gamma_i'$
$\alpha_k$~$\beta_k$~$\gamma_k$
$\alpha_k'$~$\beta_k'$~$\gamma_k'$
$t$ by the theory of binary systems, while
$\alpha_i'$~$\beta_i'$~$\gamma_i'$,
are considered as depending, by the same rules, on
$\alpha_i$~$\beta_i$~$\gamma_i$
$\xi_i$~$\eta_i$~$\zeta_i$
and $t$.

It may also be easily shown, that
$${\delta \xi_i \over \delta \alpha_i'}
      {\delta \alpha_i' \over \delta \alpha_i}
   + {\delta \xi_i \over \delta \beta_i'}
      {\delta \alpha_i' \over \delta \beta_i}
   + {\delta \xi_i \over \delta \gamma_i'}
      {\delta \alpha_i' \over \delta \gamma_i}
   = - {\delta \xi_i \over \delta \alpha_i};
   \eqno {\rm (Q^7.)}$$
with other analogous equations: the perturbation of the coordinates
$\xi_i$ may therefore be thus expressed,
$$\left. \eqalign{
\Delta \xi_i = \sum_{\prime\prime} \mathbin{.} m_k \biggl\{
     &   {\delta \xi_i \over \delta \alpha_i'}
            \int_0^t {\delta R^{(i,k)} \over \delta \alpha_i} dt
       - {\delta \xi_i \over \delta \alpha_i}
            \int_0^t {\delta R^{(i,k)} \over \delta \alpha_i'} dt \cr
     & + {\delta \xi_i \over \delta \beta_i'}
            \int_0^t {\delta R^{(i,k)} \over \delta \beta_i} dt
       - {\delta \xi_i \over \delta \beta_i}
            \int_0^t {\delta R^{(i,k)} \over \delta \beta_i'} dt \cr
     & + {\delta \xi_i \over \delta \gamma_i'}
            \int_0^t {\delta R^{(i,k)} \over \delta \gamma_i} dt
       - {\delta \xi_i \over \delta \gamma_i}
            \int_0^t {\delta R^{(i,k)} \over \delta \gamma_i'} dt
   \biggr\},\cr}
   \right\}
   \eqno {\rm (R^7.)}$$
and the perturbations of the two other coordinates may be expressed
in an analogous manner.

It results from the same principles, that in taking the first
differentials of these perturbations (R${}^7$.), the integrals
may be treated as constant; and therefore that we may either
represent the change of place of the disturbed point $m_i$, in
its relative orbit about $m_n$, by altering a little the initial
components of velocity without altering the initial position, and
then employing the rules of binary systems; or calculate at once
the perturbations of place and of velocity, by employing the same
rules, and altering at once the initial position and initial
velocity.  If we adopt the former of these two methods, we are to
employ the expressions (O${}^7$.), which may be thus summed up,
$$\left. \eqalign{
\Delta \alpha_i'
   &= \sum_{\prime\prime} \mathbin{.} m_k
         {\delta \over \delta \alpha_i}
         \int_0^t R^{(i,k)} dt,\cr
\Delta \beta_i'
   &= \sum_{\prime\prime} \mathbin{.} m_k
         {\delta \over \delta \beta_i}
         \int_0^t R^{(i,k)} dt,\cr
\Delta \gamma_i'
   &= \sum_{\prime\prime} \mathbin{.} m_k
         {\delta \over \delta \gamma_i}
         \int_0^t R^{(i,k)} dt;\cr}
   \right\}
   \eqno {\rm (S^7.)}$$
and if we adopt the latter method, we are to make,
$$\left. \multieqalign{
\Delta \alpha_i' &= \sum_{\prime\prime} \mathbin{.} m_k
      \int_0^t {\delta R^{(i,k)} \over \delta \alpha_i} dt, &
\Delta \alpha_i  &= - \sum_{\prime\prime} \mathbin{.} m_k
      \int_0^t {\delta R^{(i,k)} \over \delta \alpha_i'} dt, \cr
\Delta \beta_i' &= \sum_{\prime\prime} \mathbin{.} m_k
      \int_0^t {\delta R^{(i,k)} \over \delta \beta_i} dt, &
\Delta \beta_i  &= - \sum_{\prime\prime} \mathbin{.} m_k
      \int_0^t {\delta R^{(i,k)} \over \delta \beta_i'} dt, \cr
\Delta \gamma_i' &= \sum_{\prime\prime} \mathbin{.} m_k
      \int_0^t {\delta R^{(i,k)} \over \delta \gamma_i} dt, &
\Delta \gamma_i  &= - \sum_{\prime\prime} \mathbin{.} m_k
      \int_0^t {\delta R^{(i,k)} \over \delta \gamma_i'} dt. \cr}
   \right\}
   \eqno {\rm (T^7.)}$$
The latter was the method of {\sc Lagrange}: the former is
suggested more immediately by the principles of the present essay.

\bigbreak

{\sectiontitle
General introduction of the Time into the expression of the
Characteristic Function in any dynamical problem.\par}

\nobreak\bigskip

23.
Before we conclude this sketch of our general method in dynamics,
it will be proper to notice briefly a transformation of the
characteristic function, which may be used in all applications.
This transformation consists in putting, generally,
$$V = tH + S,
   \eqno {\rm (U^7.)}$$
and considering the part $S$, namely, the definite integral
$$S = \int_0^t (T + U) \,dt,
   \eqno {\rm (V^7.)}$$
as a function of the initial and final coordinates and of the
time, of which the variation is, by our law of varying action,
$$\delta S = - H \, \delta t + \sum \mathbin{.} m
               (   x' \, \delta x - a' \, \delta a
                 + y' \, \delta y - b' \, \delta b
                 + z' \, \delta z - c' \, \delta c ).
   \eqno {\rm (W^7.)}$$
The partial differential coefficients of the first order of this
auxiliary function $S$, are hence,
$${\delta S \over \delta t} = -H;
   \eqno {\rm (X^7.)}$$
$${\delta S \over \delta x_i} = m_i x_i',\quad
  {\delta S \over \delta y_i} = m_i y_i',\quad
  {\delta S \over \delta z_i} = m_i z_i';
   \eqno {\rm (Y^7.)}$$
and
$${\delta S \over \delta a_i} = - m_i a_i',\quad
  {\delta S \over \delta b_i} = - m_i b_i',\quad
  {\delta S \over \delta c_i} = - m_i c_i'.
   \eqno {\rm (Z^7.)}$$
These last expressions (Z${}^7$.) are forms for the final
integrals of motion of any system, corresponding to the result of
elimination of $H$ between the equations (D.) and (E.); and the
expressions (Y${}^7$.) are forms for the intermediate integrals,
more convenient in many respects than the forms already employed.

\bigbreak

24.
The limits of the present essay do not permit us here to develope
the consequences of these new expressions.  We can only observe,
that the auxiliary function $S$ must satisfy the two following
equations, in partial differentials of the first order, analogous
to, and deduced from, the equations (F.) and (G.):
$${\delta S \over \delta t} + \sum \mathbin{.} {1 \over 2m} \left\{
         \left( {\delta S \over \delta x} \right)^2
       + \left( {\delta S \over \delta y} \right)^2
       + \left( {\delta S \over \delta z} \right)^2
      \right\} = U,
   \eqno {\rm (A^8.)}$$
and
$${\delta S \over \delta t} + \sum \mathbin{.} {1 \over 2m} \left\{
         \left( {\delta S \over \delta a} \right)^2
       + \left( {\delta S \over \delta b} \right)^2
       + \left( {\delta S \over \delta c} \right)^2
      \right\} = U_0;
   \eqno {\rm (B^8.)}$$
and that to correct an approximate value $S_1$ of $S$, in the
integration of these equations, or to find the remaining part
$S_2$, if
$$S = S_1 + S_2,
   \eqno {\rm (C^8.)}$$
we may employ the symbolic equation
$${d \over dt} = {\delta \over \delta t}
      + \sum \mathbin{.} {1 \over m} \left(
            {\delta S \over \delta x} {\delta \over \delta x}
          + {\delta S \over \delta y} {\delta \over \delta y}
          + {\delta S \over \delta z} {\delta \over \delta z}
         \right);
   \eqno {\rm (D^8.)}$$
which gives, rigorously,
$${d S_2 \over dt}
   = U - U_1 + \sum \mathbin{.} {1 \over 2m} \left\{
         \left( {\delta S_2 \over \delta x} \right)^2
       + \left( {\delta S_2 \over \delta y} \right)^2
       + \left( {\delta S_2 \over \delta z} \right)^2
      \right\}
   \eqno {\rm (E^8.)}$$
if we establish by analogy the definition
$$U_1 =  {\delta S_1 \over \delta t}
       + \sum \mathbin{.} {1 \over 2m} \left\{
         \left( {\delta S_1 \over \delta x} \right)^2
       + \left( {\delta S_1 \over \delta y} \right)^2
       + \left( {\delta S_1 \over \delta z} \right)^2
      \right\};
   \eqno {\rm (F^8.)}$$
and therefore approximately
$$S_2 = \int_0^t (U - U_1) \,dt,
   \eqno {\rm (G^8.)}$$
the parts $S_1$~$S_2$ being chosen so as to vanish with the
time.  These remarks may all be extended easily, so as to embrace
relative and polar coordinates, and other marks of position, and
offer a new and better way of investigating the orbits and
perturbations of a system, by a new and better form of the
function and method of this Essay.

\nobreak\bigskip

{\it March\/} 29, 1834.


\bye
